% Traduction : les traductions du « ON » ; les adverbes d'intensité (très, trop, beaucoup, assez) — Espagnol Tle A — Cours signé OnBuch+
% Programme officiel MINESEC. Compiler avec : tectonic E23-traducciones-on-adverbios-intensidad.tex
\newcommand{\DOCMATIERE}{Espagnol}
\newcommand{\DOCNIVEAU}{Tle A}
\newcommand{\DOCMODULE}{Módulo 5 : Médias, communication et culture de la paix}
\newcommand{\DOCLECON}{E23}
\newcommand{\DOCTITRE}{Traduction : les traductions du « ON » ; les adverbes d'intensité (très, trop, beaucoup, assez)}
% OnBuch+ — préambule « cours signés » ENRICHI (compilé avec tectonic / XeLaTeX)
% Système visuel « Pop » d'OnBuch+ : fond crème (jamais blanc), bordures encre
% épaisses, ombres « dures » décalées, titres Archivo Black en capitales.
% Version MATHÉMATIQUES : graphes (pgfplots/tikz), boîtes pédagogiques
% (définition, propriété, méthode, attention, le savais-tu, exercice résolu,
% exercice, TP, bilan), page de garde et signature OnBuch+ ancrée en bas.
%
% Macros attendues dans main.tex AVANT \input{preamble} :
%   \DOCMATIERE  \DOCNIVEAU  \DOCMODULE  \DOCLECON (numéro)  \DOCTITRE
\documentclass[11pt]{article}

\usepackage{fontspec}
\usepackage[spanish,french]{babel} % français = langue principale ; l'espagnol via \esp{..} / espagnol
\usepackage[a4paper,margin=2cm,top=3.4cm,bottom=2.8cm,headheight=1.4cm,footskip=1.3cm]{geometry}
\usepackage{amsmath,amssymb,amsfonts}
\usepackage{mathtools} % \xmapsto, \coloneqq... souvent utilisés par les modèles
\usepackage{cancel} % simplifications barrées (bilans de forces)
% \abs{z} (module, valeur absolue) et \norm{u} (norme), utilisés par les modèles
\providecommand{\abs}{}\let\abs\relax\DeclarePairedDelimiter{\abs}{\lvert}{\rvert}
\providecommand{\norm}{}\let\norm\relax\DeclarePairedDelimiter{\norm}{\lVert}{\rVert}
\usepackage[table]{xcolor} % [table] : fournit aussi \rowcolors (alternance de lignes)
\usepackage{pagecolor}
\usepackage{tikz}
\usetikzlibrary{arrows.meta,positioning,calc,shapes.geometric,shapes.misc,shapes.arrows,decorations.pathmorphing,decorations.pathreplacing,decorations.markings,patterns,shadows,mindmap,trees,fit,backgrounds,matrix,intersections,angles,quotes}
\usepackage{pgfplots}
\usepackage[version=4]{mhchem} % équations nucléaires \ce{^{235}_{92}U}
\usepackage[european,siunitx]{circuitikz} % schémas électriques
\pgfplotsset{compat=1.18}
\usepgfplotslibrary{fillbetween,groupplots} % aires entre courbes (calcul intégral)
\usepackage{enumitem}
\usepackage{fancyhdr}
% [most] charge toutes les bibliothèques tcolorbox (listings, minted, raster,
% documentation...) et gonfle inutilement la mémoire TeX sur un document long
% (~300 boîtes) : on ne charge que ce qui est réellement utilisé.
\usepackage[skins,breakable]{tcolorbox}
\usepackage{graphicx}
\usepackage{colortbl}
\usepackage{booktabs}
\usepackage{tabularx}
\usepackage{diagbox} % case d'en-tête diagonale des tableaux à double entrée
% Colonnes extensibles centrée / gauche / droite (C, L, R), souvent utilisées par les modèles
\newcolumntype{C}{>{\centering\arraybackslash}X}
\newcolumntype{L}{>{\raggedright\arraybackslash}X}
\newcolumntype{R}{>{\raggedleft\arraybackslash}X}
\usepackage{array}
\usepackage{multicol}
\usepackage{siunitx}
% parse-numbers=false : accepte aussi \SI{2,5 \times 10^{-3}}{...}
\sisetup{locale=FR,output-decimal-marker={,},group-minimum-digits=4,parse-numbers=false}
\DeclareSIUnit\ohm{\ensuremath{\symup{\Omega}}} % U+2126 absent de la police
\DeclareSIUnit\division{div} % réglages d'oscilloscope (ms/div, V/div)
\DeclareSIUnit\tr{tr} % tours (vitesse de rotation, ex. \SI{1500}{\tr\per\minute})
\DeclareSIUnit\molar{M} % concentration molaire (ex. \SI{3}{\molar}, \SI{1}{\micro\molar})
\DeclareSIUnit\an{an} % année (le modèle confond parfois avec \year, un compteur TeX) — ex. \SI{5}{\kilo\meter\per\an}
\DeclareSIUnit\FCFA{FCFA} % franc CFA (ex. \SI{500}{\FCFA\per\kilo\gram})
\DeclareSIUnit\degreeBrix{\textdegree Brix} % teneur en sucre d'un sirop/jus (réfractométrie)
\DeclareSIUnit\month{mois} % le modèle utilise parfois \month, un compteur TeX, comme unité de durée
\usepackage{qrcode}
% \degree : commande fréquemment utilisée par les modèles pour un angle ou une
% température en dehors de \SI{}{} (non fournie par les paquets chargés).
\providecommand{\degree}{^{\circ}}
% \qed (fin de démonstration), utilisé par les modèles sans amsthm
\providecommand{\qed}{\hfill$\square$}
% \barre : double barre des valeurs interdites dans un tableau de variations (tkz-tab)
\providecommand{\barre}{\ensuremath{\|}}
% environnement proof (amsthm non chargé) utilisé par les modèles
\ifdefined\proof\else\newenvironment{proof}[1][Démonstration]{\par\noindent\textit{#1.}\ }{\qed\par}\fi
\usepackage{lastpage}
\usepackage{hyperref}
\hypersetup{hidelinks}

% ============================================================
% Palette « Pop » OnBuch+ — mêmes hex que la classe P (plus_ui.dart)
% ============================================================
\definecolor{poporange}{HTML}{F59321}
\definecolor{popdark}{HTML}{E07A0C}
\definecolor{popcream}{HTML}{FFF6E8}
\definecolor{popink}{HTML}{1C1714}
\definecolor{poppurple}{HTML}{7A5AE0}
\definecolor{popgreen}{HTML}{2E9E6A}
\definecolor{poppink}{HTML}{E0457B}
\definecolor{popblue}{HTML}{2F7DE1}
\definecolor{popgold}{HTML}{C98A12}
\definecolor{popmuted}{HTML}{8A7F76}
\definecolor{popline}{HTML}{EADFD2}
\definecolor{popgray}{HTML}{8A7F76}
\colorlet{popgrey}{popgray}
% Alias tolérants (le modèle nomme parfois les couleurs différemment)
\colorlet{popwhite}{white}
\colorlet{popblack}{popink}
\colorlet{popred}{poppink}
\colorlet{popyellow}{popgold}
% Teintes claires (fonds de tableaux/encadrés) — le modèle nomme parfois
% les couleurs avec un suffixe L (ex. popblueL) sans les avoir définies.
\colorlet{poporangeL}{poporange!20}
\colorlet{popdarkL}{popdark!20}
\colorlet{poppurpleL}{poppurple!20}
\colorlet{popgreenL}{popgreen!20}
\colorlet{poppinkL}{poppink!20}
\colorlet{popblueL}{popblue!20}
\colorlet{popgoldL}{popgold!20}
\colorlet{popredL}{popred!20}
\colorlet{popgrayL}{popgray!20}
% Teintes claires pour fonds de tableaux / zones
\colorlet{popblueL}{popblue!10!white}
\colorlet{poporangeL}{poporange!14!white}
\colorlet{popgreenL}{popgreen!12!white}
\colorlet{poppurpleL}{poppurple!12!white}
\colorlet{poppinkL}{poppink!12!white}
\colorlet{popgoldL}{popgold!14!white}

\pagecolor{popcream}

% ============================================================
% Typographie
% ============================================================
\defaultfontfeatures{Ligatures=TeX}
\setmainfont{PlusJakartaSans-Medium.ttf}[Path=../../../fonts/,BoldFont=PlusJakartaSans-Bold.ttf]
% Maths en sans-serif (Fira Math) avec chiffres et lettres droites de Plus
% Jakarta, pour que les formules se fondent dans le texte.
\usepackage[math-style=ISO,bold-style=ISO]{unicode-math}
\setmathfont{FiraMath-Regular.otf}[Path=../../../fonts/]
\setmathfont{PlusJakartaSans-Medium.ttf}[Path=../../../fonts/,range={up/num,up/{Latin,latin},\mathup}]
\usepackage{icomma} % virgule décimale sans espace en mode math
% Caractères mathématiques Unicode absents des polices de texte (Plus Jakarta,
% Archivo Black) : rendus par leur équivalent mathématique, en texte comme en
% formule — sinon ils s'affichent en carré vide (ex. « Arithmétique dans ℤ »).
\usepackage{newunicodechar}
\newunicodechar{ℝ}{\ensuremath{\mathbb{R}}}
\newunicodechar{ℤ}{\ensuremath{\mathbb{Z}}}
\newunicodechar{ℂ}{\ensuremath{\mathbb{C}}}
\newunicodechar{ℕ}{\ensuremath{\mathbb{N}}}
\newunicodechar{ℚ}{\ensuremath{\mathbb{Q}}}
\newunicodechar{𝔻}{\ensuremath{\mathbb{D}}}
\newunicodechar{α}{\ensuremath{\alpha}}
\newunicodechar{β}{\ensuremath{\beta}}
\newunicodechar{γ}{\ensuremath{\gamma}}
\newunicodechar{δ}{\ensuremath{\delta}}
\newunicodechar{ε}{\ensuremath{\varepsilon}}
\newunicodechar{θ}{\ensuremath{\theta}}
\newunicodechar{λ}{\ensuremath{\lambda}}
\newunicodechar{μ}{\ensuremath{\mu}}
\newunicodechar{σ}{\ensuremath{\sigma}}
\newunicodechar{τ}{\ensuremath{\tau}}
\newunicodechar{φ}{\ensuremath{\varphi}}
\newunicodechar{ψ}{\ensuremath{\psi}}
\newunicodechar{ω}{\ensuremath{\omega}}
\newunicodechar{Σ}{\ensuremath{\Sigma}}
% majuscules produites par \MakeUppercase dans les titres (α-aminés -> Α)
\newunicodechar{Α}{\ensuremath{\alpha}}
\newunicodechar{Β}{\ensuremath{\beta}}
\newunicodechar{Γ}{\ensuremath{\Gamma}}
\newunicodechar{Δ}{\ensuremath{\Delta}}
\newunicodechar{Ε}{\ensuremath{\varepsilon}}
\newunicodechar{Θ}{\ensuremath{\Theta}}
\newunicodechar{Λ}{\ensuremath{\Lambda}}
\newunicodechar{Μ}{\ensuremath{\mu}}
\newunicodechar{Τ}{\ensuremath{\tau}}
\newunicodechar{Φ}{\ensuremath{\Phi}}
\newunicodechar{Ψ}{\ensuremath{\Psi}}
\newunicodechar{Ω}{\ensuremath{\Omega}}
\newunicodechar{↦}{\ensuremath{\mapsto}}
\newunicodechar{∈}{\ensuremath{\in}}
\newunicodechar{∉}{\ensuremath{\notin}}
\newunicodechar{∧}{\ensuremath{\wedge}}
\newunicodechar{⇔}{\ensuremath{\Leftrightarrow}}
\newunicodechar{⟺}{\ensuremath{\Longleftrightarrow}}
\newunicodechar{⇒}{\ensuremath{\Rightarrow}}
\newunicodechar{⟹}{\ensuremath{\Longrightarrow}}
\newunicodechar{∘}{\ensuremath{\circ}}
\newunicodechar{∪}{\ensuremath{\cup}}
\newunicodechar{∩}{\ensuremath{\cap}}
\newunicodechar{≡}{\ensuremath{\equiv}}
\newunicodechar{⊂}{\ensuremath{\subset}}
\newunicodechar{⊥}{\ensuremath{\perp}}
\newunicodechar{∃}{\ensuremath{\exists}}
\newunicodechar{∀}{\ensuremath{\forall}}
\newunicodechar{∛}{\ensuremath{\sqrt[3]{\,}}}
\newunicodechar{∜}{\ensuremath{\sqrt[4]{\,}}}
\newunicodechar{⃗}{\ensuremath{{}^{\to}}}
\newunicodechar{ⁿ}{\ifmmode^{n}\else\textsuperscript{\ensuremath{n}}\fi}
\newunicodechar{ˣ}{\ifmmode^{x}\else\textsuperscript{\ensuremath{x}}\fi}
\newunicodechar{ᵇ}{\ifmmode^{b}\else\textsuperscript{\ensuremath{b}}\fi}
\newunicodechar{ʳ}{\ifmmode^{r}\else\textsuperscript{\ensuremath{r}}\fi}
\newunicodechar{ᵘ}{\ifmmode^{u}\else\textsuperscript{\ensuremath{u}}\fi}
\newunicodechar{ʸ}{\ifmmode^{y}\else\textsuperscript{\ensuremath{y}}\fi}
\newunicodechar{ᵃ}{\ifmmode^{a}\else\textsuperscript{\ensuremath{a}}\fi}
\newunicodechar{ᵉ}{\ifmmode^{e}\else\textsuperscript{\ensuremath{e}}\fi}
\newunicodechar{ᵏ}{\ifmmode^{k}\else\textsuperscript{\ensuremath{k}}\fi}
\newunicodechar{ᵖ}{\ifmmode^{p}\else\textsuperscript{\ensuremath{p}}\fi}
\newunicodechar{ᴺ}{\ifmmode^{N}\else\textsuperscript{\ensuremath{N}}\fi}
\newunicodechar{ᶜ}{\ifmmode^{c}\else\textsuperscript{\ensuremath{c}}\fi}
\newunicodechar{ʰ}{\ifmmode^{h}\else\textsuperscript{\ensuremath{h}}\fi}
\newunicodechar{ᵠ}{\ifmmode^{q}\else\textsuperscript{\ensuremath{q}}\fi}
\newunicodechar{ₙ}{\ifmmode_{n}\else\textsubscript{\ensuremath{n}}\fi}
\newunicodechar{ₐ}{\ifmmode_{a}\else\textsubscript{\ensuremath{a}}\fi}
\newunicodechar{ᵢ}{\ifmmode_{i}\else\textsubscript{\ensuremath{i}}\fi}
\newunicodechar{ᵣ}{\ifmmode_{r}\else\textsubscript{\ensuremath{r}}\fi}
\newunicodechar{ₖ}{\ifmmode_{k}\else\textsubscript{\ensuremath{k}}\fi}
\newunicodechar{ᵦ}{\ifmmode_{\beta}\else\textsubscript{\ensuremath{\beta}}\fi}
\newunicodechar{ₓ}{\ifmmode_{x}\else\textsubscript{\ensuremath{x}}\fi}
\newfontfamily\popdisplay{ArchivoBlack-Regular.ttf}[Path=../../../fonts/]
\newfontfamily\poplogo{SpaceGrotesk-Bold.ttf}[Path=../../../fonts/]
\newfontfamily\popsemi{PlusJakartaSans-SemiBold.ttf}[Path=../../../fonts/]
\newfontfamily\popxbold{PlusJakartaSans-ExtraBold.ttf}[Path=../../../fonts/]
\newfontfamily\popxboldls{PlusJakartaSans-ExtraBold.ttf}[Path=../../../fonts/,LetterSpace=3.0]

\setlist[enumerate]{leftmargin=1.7em,itemsep=5pt,topsep=4pt}
\setlist[itemize]{leftmargin=1.7em,itemsep=4pt,topsep=4pt,label={\color{poporange}\textbullet}}
\setlength{\parindent}{0pt}
\setlength{\parskip}{5pt}
\renewcommand{\arraystretch}{1.25}

% pgfplots : style Pop par défaut
\pgfplotsset{
  every axis/.append style={axis line style={popink,line width=1pt},tick style={popink},
    grid style={popline,line width=0.5pt},label style={font=\small},tick label style={font=\footnotesize}},
  popcurve/.style={poporange,line width=1.8pt,no markers,smooth},
}
% Même style disponible dans un \draw TikZ simple (hors pgfplots)
\tikzset{popcurve/.style={poporange,line width=1.8pt}}
\tikzset{popgrid/.style={popline,very thin}}
\colorlet{popcurve}{poporange} % le nom du style est parfois utilisé comme couleur

% ============================================================
% Pill / squiggle
% ============================================================
\newcommand{\poppill}[3]{%
  \tikz[baseline=-0.55ex]\node[draw=popink,line width=1.2pt,rounded corners=2.6pt,%
    fill=#2,inner xsep=6.5pt,inner ysep=3pt]{%
    \popxboldls\fontsize{7}{7}\selectfont\color{#3}\MakeUppercase{#1}};%
}
\newcommand{\popsquiggle}[1]{%
  \tikz[baseline=-0.4ex]{
    \draw[#1,line width=2.6pt,line cap=round]
      plot [smooth] coordinates {(0,0.09) (0.34,0.01) (0.68,0.21) (1.02,0.01) (1.36,0.10)};
  }%
}
% Mot-clé surligné (marqueur orange sous le texte)
\newcommand{\cle}[1]{{\popxbold\color{popdark}#1}}

% ============================================================
% En-tête / pied
% ============================================================
\pagestyle{fancy}
\fancyhf{}
\renewcommand{\headrulewidth}{0pt}
\renewcommand{\footrulewidth}{0pt}
\lhead{\raisebox{-0.15ex}{\poplogo\fontsize{13}{13}\selectfont\color{popink}OnBuch\color{poporange}+}}
\rhead{\poppill{\DOCMATIERE\ \textbullet\ \DOCNIVEAU\ \textbullet\ Leçon \DOCLECON}{white}{popink}}
\lfoot{\popxbold\fontsize{7.6}{7.6}\selectfont\color{popmuted}\MakeUppercase{Cours signé OnBuch+}\ \textbullet\ {\popsemi\DOCTITRE}}
\rfoot{\tikz[baseline=-0.6ex]\node[rounded corners=8.5pt,fill=popink,minimum height=17pt,minimum width=17pt,inner xsep=5pt,inner sep=0]{\popxbold\fontsize{8}{8}\selectfont\color{popcream}\thepage\,/\,\pageref*{LastPage}};}

% ============================================================
% Titres
% ============================================================
\newcommand{\chaptitre}[2]{%
  \begin{tcolorbox}[enhanced,colback=poporange,colframe=popink,boxrule=2.6pt,arc=11pt,
      drop shadow=popink,left=15pt,right=15pt,top=12pt,bottom=11pt,before skip=3pt,after skip=18pt]
    \poppill{#2}{popink}{popcream}\par\vspace{9pt}
    {\raggedright\hyphenpenalty=10000\exhyphenpenalty=10000\popdisplay\fontsize{22}{24}\selectfont\color{popcream}\MakeUppercase{#1}\par}
  \end{tcolorbox}
}
% Section numérotée : pastille orange avec le numéro + titre Archivo
\newcounter{popsec}
\newcommand{\coursec}[1]{%
  \stepcounter{popsec}\setcounter{popsub}{0}%
  \par\vspace{16pt}\noindent
  \begin{minipage}{\linewidth}
  \tikz[baseline=-0.7ex]\node[draw=popink,line width=1.6pt,fill=poporange,rounded corners=4pt,minimum size=22pt,inner sep=0]{\popdisplay\fontsize{12}{12}\selectfont\color{popcream}\thepopsec};%
  \hspace{8pt}{\popdisplay\fontsize{15}{17}\selectfont\color{popink}\MakeUppercase{#1}}\par
  \vspace{4pt}\noindent{\color{poporange}\rule{36pt}{3.6pt}}
  \end{minipage}\par\nopagebreak\vspace{8pt}
}
\newcounter{popsub}
\newcommand{\courssub}[1]{%
  \stepcounter{popsub}%
  \par\vspace{9pt}\noindent{\popxbold\fontsize{11.5}{13}\selectfont\color{popink}{\color{poporange}\thepopsec.\thepopsub}\hspace{6pt}#1}\par\nopagebreak\vspace{4pt}
}

% ============================================================
% Boîtes pédagogiques — même recette : fond blanc, bordure épaisse colorée,
% ombre dure encre, badge plein. Titre optionnel [#1].
% ============================================================
\tcbset{popbase/.style={breakable,enhanced,colback=white,boxrule=2.3pt,arc=9pt,
  shadow={3pt}{-3pt}{0pt}{popink},left=14pt,right=14pt,top=11pt,bottom=11pt,before skip=11pt,after skip=15pt,
  fontupper=\popsemi\fontsize{10.6}{14.5}\selectfont\color{popink},
  fontlower=\popsemi\fontsize{10.6}{14.5}\selectfont\color{popink}}}
% Coche dessinée (le glyphe ✓ n'existe pas dans les polices du document)
\newcommand{\popcheck}[1]{\tikz[baseline=-0.5ex]\draw[#1,line width=1.6pt,line cap=round,line join=round](-0.13,0.0)--(-0.04,-0.09)--(0.13,0.1);}
\providecommand{\coche}{\popcheck{popgreen}}
\providecommand{\resultat}[1]{\par\smallskip\noindent\fcolorbox{popgreen}{popgreenL}{\parbox{\dimexpr\linewidth-2\fboxsep-2\fboxrule\relax}{\textbf{Résultat.} #1}}\par\smallskip}
\newcommand{\popboxhead}[3]{\poppill{#1}{#2}{white}\ifx\relax#3\relax\else\hspace{6pt}{\popxbold\fontsize{10.6}{12}\selectfont\color{popink}#3}\fi\par\vspace{7pt}}

\newtcolorbox{aretenir}[1][]{popbase,colframe=popgreen,before upper={\popboxhead{À retenir}{popgreen}{#1}}}
% Texte en espagnol (document de lecture, dialogue, poème, extrait) : cadre
% d'étude. Un titre optionnel (source, auteur) s'affiche dans l'en-tête.
\newtcolorbox{textebox}[1][]{popbase,colframe=popblue,before upper={\popboxhead{Texto}{popblue}{#1}\begin{otherlanguage}{spanish}},after upper={\end{otherlanguage}}}
% Marques ✓ / ✗ (forme correcte / fautive) souvent utilisées par les modèles
\providecommand{\cross}{\textcolor{poppink}{\ensuremath{\times}}}
\providecommand{\xmark}{\cross}\providecommand{\crossmark}{\cross}\providecommand{\cmark}{\ensuremath{\checkmark}}
% Ligne de réponse / justification à compléter (inventée par les modèles)
\providecommand{\justification}[1]{\par\noindent\textit{Justification~:} \dotfill\par}
% \textipa (package tipa non chargé) : la police ne couvre pas l'API -> simple passage
\providecommand{\textipa}[1]{#1}
% Soulignement (ulem non chargé) : \uline, \ul
\providecommand{\uline}[1]{\underline{#1}}\providecommand{\ul}[1]{\underline{#1}}
% Mot ou phrase en espagnol à l'intérieur d'un paragraphe français
\newcommand{\esp}[1]{\ifmmode\text{\textit{#1}}\else\begingroup\selectlanguage{spanish}\itshape #1\endgroup\fi}
\newtcolorbox{exemplebox}[1][]{popbase,colframe=poporange,before upper={\popboxhead{Exemple}{poporange}{#1}}}
\newtcolorbox{definition}[1][]{popbase,colframe=popblue,before upper={\popboxhead{Définition}{popblue}{#1}}}
\newtcolorbox{propriete}[1][]{popbase,colframe=poppurple,before upper={\popboxhead{Propriété}{poppurple}{#1}}}
\newtcolorbox{methode}[1][]{popbase,colframe=popgold,before upper={\popboxhead{Méthode}{popgold}{#1}}}
\newtcolorbox{attention}[1][]{popbase,colframe=poppink,before upper={\popboxhead{Attention}{poppink}{#1}}}
\newtcolorbox{experience}[1][]{popbase,colframe=popblue,colback=popblueL,before upper={\popboxhead{Expérience}{popblue}{#1}}}
% « Le savais-tu ? » : carte violette pleine (contexte, culture, Cameroun)
\newtcolorbox{savaistu}[1][]{popbase,colback=poppurple,colframe=popink,
  fontupper=\popsemi\fontsize{10.4}{14.2}\selectfont\color{white},
  before upper={\poppill{Le savais-tu\,?}{white}{poppurple}\ifx\relax#1\relax\else\hspace{6pt}{\popxbold\color{white}#1}\fi\par\vspace{7pt}}}
% Exercice résolu : énoncé en haut, \tcblower puis la solution sur fond crème
\newcounter{popexo}\newcounter{popexr}
\newtcolorbox{exoresolu}[1][]{popbase,colframe=poporange,colbacklower=popcream,
  segmentation style={popink,line width=1.2pt,dashed},
  before upper={\stepcounter{popexr}\popboxhead{Exercice résolu \thepopexr}{poporange}{#1}},
  before lower={\poppill{Solution}{popink}{popcream}\par\vspace{6pt}}}
% Exercice (non résolu en place) — la correction est dans la partie Corrigés
\newtcolorbox{exercice}[1][]{popbase,colframe=popink,
  before upper={\stepcounter{popexo}\popboxhead{Exercice \thepopexo}{popink}{#1}}}
% Corrigé
\newtcolorbox{corrige}[1][]{popbase,colframe=popgreen,colback=popgreenL,
  before upper={\popboxhead{Corrigé}{popgreen}{#1}}}
% Objectifs / prérequis / bilan
\newtcolorbox{objectifs}{popbase,colframe=popink,colback=poporangeL,
  before upper={\popboxhead{Objectifs de la leçon}{popink}{}}}
\newtcolorbox{prerequis}{popbase,colframe=popink,colback=popblueL,
  before upper={\popboxhead{Prérequis}{popblue}{}}}
\newtcolorbox{bilan}{popbase,colframe=popink,colback=white,boxrule=2.8pt,
  before upper={\popboxhead{Fiche bilan}{poporange}{L'essentiel à connaître pour l'examen}}}
% Figure encadrée avec légende
\newtcolorbox{popfigure}[1][]{enhanced,breakable=false,colback=white,colframe=popline,boxrule=1.4pt,arc=8pt,
  left=10pt,right=10pt,top=10pt,bottom=8pt,before skip=10pt,after skip=12pt,halign=center,
  lower separated=false,halign lower=center,
  fontlower=\popsemi\fontsize{9}{11.5}\selectfont\color{popmuted},
  #1}
\newcommand{\legende}[1]{\par\vspace{6pt}{\popsemi\fontsize{9}{11.5}\selectfont\color{popmuted}#1}}

% ============================================================
% Signature OnBuch+
% ============================================================
\newcommand{\poplogoinline}[1]{{\poplogo\fontsize{#1}{#1}\selectfont\color{popink}OnBuch\color{poporange}+}}
\newcommand{\signatureonbuch}{%
  \begin{tcolorbox}[enhanced,colback=popink,colframe=popink,boxrule=2.2pt,arc=10pt,
      drop shadow=poporange,left=14pt,right=14pt,top=10pt,bottom=10pt,before skip=0pt,after skip=0pt]
    \tikz[baseline=-0.7ex]\node[circle,fill=popgreen,draw=popcream,line width=1.3pt,minimum size=22pt,inner sep=0]{\popcheck{white}};%
    \hspace{9pt}\begin{minipage}[c]{0.62\linewidth}
      {\popxbold\fontsize{12}{13}\selectfont\color{popcream}Signé\ \textbullet\ {\poplogo OnBuch\color{poporange}+}}\par\vspace{2pt}
      {\raggedright\popsemi\fontsize{8}{10.5}\selectfont\color{popline}\MakeUppercase{Équipe pédagogique OnBuch+}\\\MakeUppercase{Conforme au programme officiel MINESEC}\par}
    \end{minipage}\hfill
    {\poplogo\fontsize{22}{22}\selectfont\color{popcream}OnBuch\color{poporange}+}
  \end{tcolorbox}%
}

% ============================================================
% Page de garde
% #1 titre · #2 matière · #3 niveau · #4 module · #5 n° leçon · #6 durée ·
% #7 description / objectifs (texte ou itemize)
% ============================================================
\newcommand{\pagegarde}[7]{%
  \thispagestyle{empty}
  \newgeometry{a4paper,margin=1.7cm,top=1.6cm,bottom=1.4cm}%
  \begin{tikzpicture}[remember picture,overlay]
    \fill[poporange!18] ([xshift=-3.2cm,yshift=-3.6cm]current page.north east) circle (4.6cm);
    \fill[poppurple!14] ([xshift=2.4cm,yshift=7.6cm]current page.south west) circle (3.8cm);
  \end{tikzpicture}%
  {\poplogo\fontsize{30}{30}\selectfont\color{popink}OnBuch\color{poporange}+}\hfill
  \raisebox{0.6ex}{\poppill{Cours signé \textbullet\ Premium}{popink}{popcream}}\par
  \vspace{1pt}\popsquiggle{poppurple}\par
  \vspace{8pt}
  \begin{tcolorbox}[enhanced,colback=poporange,colframe=popink,boxrule=2.8pt,arc=13pt,
      drop shadow=popink,left=18pt,right=18pt,top=12pt,bottom=13pt,after skip=10pt]
    \poppill{#2 \textbullet\ #3}{popink}{popcream}\hspace{5pt}\poppill{#4}{white}{popink}\par\vspace{8pt}
    {\popdisplay\fontsize{14}{14}\selectfont\color{popink}LEÇON #5}\par\vspace{4pt}
    {\raggedright\hyphenpenalty=10000\exhyphenpenalty=10000\popdisplay\fontsize{25}{27}\selectfont\color{popcream}\MakeUppercase{#1}\par}
    \vspace{8pt}
    {\popxbold\fontsize{9.5}{9.5}\selectfont\color{popink}\MakeUppercase{Durée conseillée : #6}}
  \end{tcolorbox}
  \begin{tcolorbox}[enhanced,colback=white,colframe=popink,boxrule=2.2pt,arc=10pt,
      drop shadow=popink,left=16pt,right=16pt,top=10pt,bottom=10pt,after skip=8pt]
    \poppill{Dans cette leçon}{poppurple}{white}\par\vspace{5pt}
    \popsemi\fontsize{9.3}{12.6}\selectfont\color{popink}#7
  \end{tcolorbox}
  \noindent\begin{minipage}[t]{0.487\linewidth}
    \begin{tcolorbox}[enhanced,colback=popgreen,colframe=popink,boxrule=2.2pt,arc=10pt,
        drop shadow=popink,left=11pt,right=11pt,top=10pt,bottom=10pt,height=3.1cm,valign=center]
      \tcbox[colback=white,colframe=popink,boxrule=1.3pt,arc=5pt,boxsep=0pt,nobeforeafter,
        left=3pt,right=3pt,top=3pt,bottom=3pt,valign=center]{\qrcode[height=1.9cm]{https://play.google.com/store/apps/details?id=com.luvvix.onbuch&hl=fr}}%
      \hspace{8pt}\begin{minipage}[c]{0.5\linewidth}
        {\popdisplay\fontsize{11}{12.5}\selectfont\color{white}\MakeUppercase{Télécharge OnBuch}}\par\vspace{4pt}
        {\raggedright\popsemi\fontsize{8.4}{11}\selectfont\color{white}Gratuit \textbullet\ Google Play\\Tuteur IA, annales, concours, résultats\par}
      \end{minipage}
    \end{tcolorbox}
  \end{minipage}\hfill
  \begin{minipage}[t]{0.487\linewidth}
    \begin{tcolorbox}[enhanced,colback=poppurple,colframe=popink,boxrule=2.2pt,arc=10pt,
        drop shadow=popink,left=11pt,right=11pt,top=10pt,bottom=10pt,height=3.1cm,valign=center]
      \tikz[baseline=-0.7ex]\node[circle,fill=white,draw=popink,line width=1.3pt,minimum size=1.6cm,inner sep=0]{\popdisplay\fontsize{24}{24}\selectfont\color{poppurple}+};%
      \hspace{8pt}\begin{minipage}[c]{0.6\linewidth}
        {\popdisplay\fontsize{11}{12.5}\selectfont\color{white}\MakeUppercase{Passe à OnBuch+}}\par\vspace{4pt}
        {\raggedright\popsemi\fontsize{8.4}{11}\selectfont\color{white}Tous les cours signés, Tuteur IA illimité, fiches PDF — onglet OnBuch+ de l'app\par}
      \end{minipage}
    \end{tcolorbox}
  \end{minipage}
  \vspace{8pt}
  \popcontenu
  \vfill
  \signatureonbuch
  \restoregeometry
  \newpage
}

% Rangée « Ce cours contient » (page de garde)
\newcommand{\poptile}[3]{% #1 couleur #2 gros texte #3 légende
  \begin{tcolorbox}[enhanced,colback=#1,colframe=popink,boxrule=2pt,arc=9pt,shadow={2.5pt}{-2.5pt}{0pt}{popink},
      left=6pt,right=6pt,top=9pt,bottom=9pt,halign=center,valign=center,height=1.75cm,width=\linewidth,nobeforeafter]
    {\popdisplay\fontsize{12.5}{13}\selectfont\color{white}\MakeUppercase{#2}}\par\vspace{3pt}
    {\centering\popsemi\fontsize{7.8}{9.5}\selectfont\color{white}#3\par}
  \end{tcolorbox}}
\newcommand{\popcontenu}{%
  \par\noindent{\popxboldls\fontsize{7.5}{7.5}\selectfont\color{popmuted}CE COURS SIGNÉ CONTIENT}\par\vspace{7pt}
  \noindent\begin{minipage}[t]{0.235\linewidth}\poptile{popblue}{Cours}{complet, illustré, pas à pas}\end{minipage}\hfill
  \begin{minipage}[t]{0.235\linewidth}\poptile{poporange}{Méthodes}{et exercices résolus}\end{minipage}\hfill
  \begin{minipage}[t]{0.235\linewidth}\poptile{poppink}{Exercices}{type examen + corrigés}\end{minipage}\hfill
  \begin{minipage}[t]{0.235\linewidth}\poptile{popgreen}{Fiche bilan}{l'essentiel à retenir}\end{minipage}\par}

% Sommaire
\newtcolorbox{sommaire}{enhanced,colback=white,colframe=popink,boxrule=2.2pt,arc=10pt,
  drop shadow=popink,left=18pt,right=18pt,top=13pt,bottom=13pt,before skip=6pt,after skip=20pt,
  before upper={\poppill{Sommaire}{poppurple}{white}\par\vspace{9pt}\popsemi\fontsize{10.6}{14.5}\selectfont\color{popink}}}

% Signature de fin de cours — poussée tout en bas de la dernière page
\newcommand{\findecours}{%
  \par\vfill\nopagebreak
  \begin{center}\popsquiggle{poporange}\end{center}\vspace{4pt}
  \signatureonbuch
}
\AtBeginDocument{\def\llbracket{\mathopen{[\mkern-3.5mu[}}\def\rrbracket{\mathclose{]\mkern-3.5mu]}}} % intervalles d'entiers (glyphe ⟦ absent de la police)
% Glyphes absents de Fira Math : redéfinis à partir de glyphes disponibles
\AtBeginDocument{%
  \def\setminus{\mathbin{\backslash}}\let\smallsetminus\setminus
  \def\vdots{\mathord{\vbox{\baselineskip3.2pt\lineskiplimit0pt\kern4pt\hbox{.}\hbox{.}\hbox{.}}}}%
  \def\ddots{\mathinner{\mkern1mu\raise6.5pt\hbox{.}\mkern2mu\raise3.6pt\hbox{.}\mkern2mu\raise0.7pt\hbox{.}\mkern1mu}}%
  \def\triangle{\mathord{\scalebox{0.9}{$\symup{\Delta}$}}}%
  \def\nabla{\mathord{\raisebox{\depth}{\scalebox{1}[-1]{$\symup{\Delta}$}}}}%
  \def\checkmark{\popcheck{popdark}}%
  \def\star{\mathbin{\ast}}\def\bigstar{\mathord{\ast}}%
  \def\diamond{\mathbin{\scalebox{0.6}{$\Diamond$}}}%
  \def\bigcup{\mathop{\mathchoice{\raisebox{-0.3em}{\scalebox{1.7}{$\cup$}}}{\scalebox{1.25}{$\cup$}}{\cup}{\cup}}\displaylimits}%
  \def\bigcap{\mathop{\mathchoice{\raisebox{-0.3em}{\scalebox{1.7}{$\cap$}}}{\scalebox{1.25}{$\cap$}}{\cap}{\cap}}\displaylimits}%
  \def\lesseqqgtr{\mathrel{\lessgtr}}\def\bigoplus{\mathop{\oplus}\displaylimits}%
}
\providecommand{\ssi}{si et seulement si} % abréviation fréquente
\AtBeginDocument{\providecommand{\wideparen}{\overparen}} % arc de cercle

%% FIGFIX-BEGIN v3 — correctifs globaux des figures (docs/onbuch-generation/tools/figfix)
\usepackage{adjustbox}\usepackage{etoolbox}
% toute figure TikZ/pgfplots plus large que la ligne est réduite à la largeur disponible
\BeforeBeginEnvironment{tikzpicture}{\begin{adjustbox}{max width=\linewidth}}
\AfterEndEnvironment{tikzpicture}{\end{adjustbox}}
% légendes pgfplots lisibles par-dessus les courbes
\pgfplotsset{every axis/.append style={legend style={fill=white,fill opacity=0.92,text opacity=1,draw=black!15,inner sep=3pt}}}
% étiquettes d'arêtes (syntaxe edge["..."]) avec fond blanc
\tikzset{every edge quotes/.append style={fill=white,inner sep=1.2pt}}
% bulles de cartes mentales : texte à la ligne dans la bulle, bulle un peu plus grande
\tikzset{every concept/.append style={text width=2.0cm,align=center,minimum size=2.3cm,inner sep=2pt}}
%% FIGFIX-END
\begin{document}

\pagegarde{Traduction : les traductions du « ON » ; les adverbes d'intensité (très, trop, beaucoup, assez)}{Espagnol}{Tle A}{Módulo 5 : Médias, communication et culture de la paix}{E23}{2 h}{Cette leçon de traduction étudie les différentes équivalences espagnoles du pronom français « on » (se impersonnel, 3e personne du pluriel, uno, la gente, nosotros) ainsi que l'emploi des adverbes d'intensité (muy, mucho, demasiado, bastante, tan, tanto et leurs apocopes). Elle vise à rendre l'élève capable de choisir la bonne traduction selon le contexte et d'éviter les confusions classiques muy/mucho.
\vspace{4pt}
\begin{multicols}{2}\small
\begin{itemize}
\item À la fin de cette leçon, je sais identifier la valeur du « on » français dans une phrase (indéfini, collectif, équivalent de « nous »)
\item je sais traduire « on » par la construction impersonnelle avec « se » + 3e personne du singulier
\item je sais traduire « on » par la 3e personne du pluriel quand le sens est « les gens, ils »
\item je sais traduire « on » par « uno », « la gente » ou « nosotros » selon le contexte
\item je sais employer correctement muy, mucho, demasiado et bastante devant l'adjectif, l'adverbe ou le verbe
\item je sais distinguer muy (devant adjectif/adverbe) et mucho (avec verbe ou nom)
\end{itemize}\end{multicols}}

\begin{sommaire}
\begin{multicols}{2}
\begin{enumerate}
\item Observation : le « on » français et ses visages espagnols
\item La construction impersonnelle avec « se »
\item Les autres traductions de « on » : 3e personne du pluriel, uno, la gente, nosotros
\item Les adverbes d'intensité : muy, mucho, demasiado, bastante
\item Tan, tanto et les apocopes
\item Entraînement à la traduction : thème et version
\item Activité d'intégration
\item Méthodes et exercices résolus
\item Exercices
\item Corrigés des exercices
\item Fiche bilan
\end{enumerate}
\end{multicols}
\end{sommaire}

\begin{objectifs}
\begin{itemize}
\item À la fin de cette leçon, je sais identifier la valeur du « on » français dans une phrase (indéfini, collectif, équivalent de « nous »)
\item je sais traduire « on » par la construction impersonnelle avec « se » + 3e personne du singulier
\item je sais traduire « on » par la 3e personne du pluriel quand le sens est « les gens, ils »
\item je sais traduire « on » par « uno », « la gente » ou « nosotros » selon le contexte
\item je sais employer correctement muy, mucho, demasiado et bastante devant l'adjectif, l'adverbe ou le verbe
\item je sais distinguer muy (devant adjectif/adverbe) et mucho (avec verbe ou nom)
\item je sais employer tan/tanto et leurs apocopes (tantos, tanta, tantas, buen, gran...)
\item je sais repérer et éviter les pièges de traduction fréquents des francophones
\item je sais traduire un court texte du français vers l'espagnol et inversement en mobilisant ces outils
\end{itemize}
\end{objectifs}

\begin{prerequis}
\begin{itemize}
\item La conjugaison des verbes au présent de l'indicatif (3e personnes)
\item Les pronoms personnels sujets espagnols (yo, tú, él, nosotros...)
\item La place de l'adjectif et l'accord en genre et en nombre
\item Les adverbes en -mente (Première)
\item La construction réflexive avec « se » (verbes pronominaux)
\end{itemize}
\end{prerequis}

\newpage

\coursec{Observation : le « on » français et ses visages espagnols}

Au marché Mokolo de Yaoundé, un guide francophone accueille un groupe de touristes mexicains. Souriant, il leur lance : « On mange bien ici, on est très accueilli, mais on marche beaucoup sous un soleil trop fort ! » Les touristes attendent la traduction... et le guide hésite. Comment rendre en espagnol ce petit mot « on » qui change de sens à chaque phrase ? Et comment doser l'intensité : \esp{muy}, \esp{mucho}, \esp{demasiado}, \esp{bastante} ? C'est tout l'art du traducteur que cette leçon va vous apprendre.

La problématique de cette première partie est simple : le français possède un seul pronom indéfini, « on », mais l'espagnol n'a pas de mot équivalent. Il faut donc, avant de traduire, se poser une question de sens : qui fait l'action ? Une personne indéfinie ? Les gens en général ? Ou bien « nous » ? Observons d'abord des exemples.

\courssub{Lecture d'exemples bilingues : cinq visages du même mot}

Lisons attentivement les paires de phrases suivantes. Pour chacune, demandons-nous : qui fait l'action ?

\begin{exemplebox}[Cinq traductions possibles de « on »]
\begin{enumerate}
\item « On mange bien au Cameroun. » $\rightarrow$ \esp{En Camerún se come bien.} (les gens en général, construction impersonnelle)
\item « On nous a dit que le match était annulé. » $\rightarrow$ \esp{Nos dijeron que el partido estaba cancelado.} (quelqu'un, des personnes non nommées : 3\textsuperscript{e} personne du pluriel)
\item « Quand on est jeune, on rêve beaucoup. » $\rightarrow$ \esp{Cuando uno es joven, sueña mucho.} (une personne quelconque : \esp{uno})
\item « On va au marché, tu viens ? » (= nous) $\rightarrow$ \esp{Vamos al mercado, ¿vienes?} (\esp{nosotros} sous-entendu)
\item « On dit que les prix vont baisser. » $\rightarrow$ \esp{La gente dice que los precios van a bajar.} ou \esp{Se dice que los precios van a bajar.} (les gens, rumeur générale)
\end{enumerate}
\end{exemplebox}

\begin{definition}[Le pronom « on » en français]
En français, « on » est un \cle{pronom indéfini} de la 3\textsuperscript{e} personne du singulier. Il peut désigner : une personne indéterminée (« on a sonné »), les gens en général (« on vit bien ici »), ou, dans la langue courante, « nous » (« on part demain »). L'espagnol ne possède \textbf{aucun mot unique} correspondant à « on » : il choisit une construction selon le sens.
\end{definition}

\courssub{Un texte d'observation : toutes les traductions réunies}

Lisons ce court texte original. Chaque occurrence du « on » français y est rendue par une construction différente. Soulignez-les mentalement et identifiez à quel sens de « on » elles correspondent.

\begin{textebox}[En Camerún se come muy bien]
En Camerún se come muy bien. La gente dice que la comida es barata, pero uno sabe que a veces es demasiado cara. En los mercados de Yaoundé, como Mokolo, se encuentran frutas, verduras y pescado fresco. Cuando uno es joven, come de todo sin pensar en la salud; pero los médicos dicen que hay que comer con moderación. Nosotros comemos mucho en familia los domingos: preparamos eru, ndolé o poulet DG, y todos están bastante contentos. También se dice que la cocina camerunesa es una de las mejores de África. ¡Y es verdad! Aquí se habla francés e inglés, pero en la mesa todos hablamos el mismo idioma: el de la buena comida.
\end{textebox}

\textbf{Glossaire :} \esp{barato/a} = bon marché ; \esp{verduras} = légumes ; \esp{pescado fresco} = poisson frais ; \esp{la salud} = la santé ; \esp{con moderación} = avec modération ; \esp{la cocina} = la cuisine (l'art culinaire) ; \esp{camerunesa} = camerounaise.

\textbf{Repérage :} \esp{se come} = on mange (impersonnel) ; \esp{la gente dice} = on dit (les gens) ; \esp{uno sabe} = on sait (chacun) ; \esp{se encuentran} = on trouve (impersonnel, verbe au pluriel devant un complément pluriel) ; \esp{los médicos dicen} = on dit (3\textsuperscript{e} personne du pluriel, sujet précisé) ; \esp{comemos} = on mange (= nous) ; \esp{se dice} = on dit (rumeur).

\courssub{Le constat : un mot français, plusieurs constructions espagnoles}

\begin{propriete}[Principe fondamental de la traduction de « on »]
Il n'existe pas de mot espagnol équivalent à « on ». Pour traduire, il faut d'abord déterminer le \cle{sens} de « on » dans la phrase française, puis choisir la construction correspondante :
\begin{itemize}
\item « on » = les gens en général, action impersonnelle $\rightarrow$ \esp{se} + verbe à la 3\textsuperscript{e} personne du singulier : \esp{Aquí se habla francés.} (« Ici on parle français. »)
\item « on » = quelqu'un, des personnes non identifiées $\rightarrow$ 3\textsuperscript{e} personne du pluriel sans sujet exprimé : \esp{Nos dijeron que...} (« On nous a dit que... »)
\item « on » = chacun, une personne quelconque (vérité générale) $\rightarrow$ \esp{uno} + verbe à la 3\textsuperscript{e} personne du singulier : \esp{Cuando uno es joven...} (« Quand on est jeune... »)
\item « on » = les gens (sujet collectif nommé) $\rightarrow$ \esp{la gente} + verbe au singulier : \esp{La gente dice que...} (« On dit que... »)
\item « on » = nous $\rightarrow$ verbe à la 1\textsuperscript{re} personne du pluriel : \esp{Vamos al mercado.} (« On va au marché. »)
\end{itemize}
\end{propriete}

\begin{center}
\begin{tabularx}{\linewidth}{|X|X|X|}
\hline
\rowcolor{popblueL} Français & Español & Sens de « on » \\
\hline
On mange bien au Cameroun. & En Camerún se come bien. & les gens en général (impersonnel) \\
\hline
On nous a dit que... & Nos dijeron que... & quelqu'un (3\textsuperscript{e} pers. pluriel) \\
\hline
Quand on est jeune... & Cuando uno es joven... & chacun, une personne quelconque \\
\hline
On va au marché. & Vamos al mercado. & nous \\
\hline
On dit que... & La gente dice que... / Se dice que... & les gens, la rumeur \\
\hline
Ici on parle français et anglais. & Aquí se habla francés e inglés. & impersonnel \\
\hline
\end{tabularx}
\end{center}

\begin{attention}[Piège majeur : il n'y a pas de pronom « on » en espagnol]
Ne cherchez jamais un mot unique pour traduire « on ». Les erreurs typiques des francophones sont : inventer un pronom, ou utiliser systématiquement \esp{se} même quand « on » signifie « nous ». Ainsi, « On est allés à Douala hier » (entre amis) ne se traduit PAS par \esp{Se fue a Douala}, mais par \esp{Fuimos a Douala ayer.} Autre piège : \esp{la gente} est un nom \textbf{singulier} en espagnol : \esp{La gente dice} (et jamais \esp{la gente dicen}), alors que « les gens disent » prend un pluriel en français.
\end{attention}

\courssub{La question-guide : qui fait l'action ?}

Avant toute traduction, posez-vous la question : qui fait l'action ? La carte mentale suivante résume les cinq visages espagnols du « on » français.

\begin{popfigure}
\begin{tikzpicture}[mindmap, grow cyclic, every node/.style={concept}, root concept/.append style={concept color=popblue, font=\Large\bfseries}, level 1/.append style={level distance=3.6cm, sibling angle=72}, level 2/.append style={level distance=2.4cm, sibling angle=45, concept color=popblueL}]
\node[root concept]{ON}
  child[concept color=poporange] { node {se + verbe (3\textsuperscript{e} sg)} child { node[align=center, font=\small] {Se come\\ bien.} } }
  child[concept color=popgreen] { node {3\textsuperscript{e} pers. pluriel} child { node[align=center, font=\small] {Nos\\ dijeron...} } }
  child[concept color=poppurple] { node {uno} child { node[align=center, font=\small] {Cuando uno\\ es joven...} } }
  child[concept color=poppink] { node {la gente} child { node[align=center, font=\small] {La gente\\ dice...} } }
  child[concept color=popgold] { node {nosotros} child { node[align=center, font=\small] {Vamos al\\ mercado.} } };
\end{tikzpicture}
\legende{Carte mentale : les cinq traductions espagnoles du « on » français, avec un mini-exemple pour chaque branche.}
\end{popfigure}

\begin{methode}[Choisir la bonne traduction de « on » en deux étapes]
\begin{enumerate}
\item \textbf{Remplacez mentalement} « on » par « nous », « les gens » ou « quelqu'un / chacun ». Si le remplacement par « nous » donne une phrase naturelle dans le contexte, c'est la 1\textsuperscript{re} personne du pluriel. Si « les gens » convient, c'est l'impersonnel (\esp{se}) ou \esp{la gente}. Si « quelqu'un » ou « chacun » convient, c'est la 3\textsuperscript{e} personne du pluriel ou \esp{uno}.
\item \textbf{Appliquez la correspondance} : nous $\rightarrow$ verbe en \esp{-mos} ; les gens (action impersonnelle) $\rightarrow$ \esp{se} + 3\textsuperscript{e} sg ; les gens (opinion répandue) $\rightarrow$ \esp{la gente} + 3\textsuperscript{e} sg ; quelqu'un (sujet inconnu) $\rightarrow$ 3\textsuperscript{e} pluriel ; chacun (vérité générale) $\rightarrow$ \esp{uno} + 3\textsuperscript{e} sg.
\end{enumerate}
Exemple guidé : « On nous a volé nos sacs. » Étape 1 : « on » = quelqu'un (des voleurs inconnus). Étape 2 : 3\textsuperscript{e} personne du pluriel $\rightarrow$ \esp{Nos robaron las bolsas.}
\end{methode}

\begin{exoresolu}[Traduisez : « Au Cameroun, on parle plus de deux cents langues, et on est très fier de cette richesse. »]
Énoncé : traduisez la phrase complète en espagnol, en justifiant le choix de chaque construction pour « on ».
\tcblower
Solution détaillée :
\begin{itemize}
\item Premier « on » : « on parle » = les gens en général, action impersonnelle $\rightarrow$ \esp{se} + 3\textsuperscript{e} personne du singulier : \esp{se habla}. Attention, « plus de deux cents langues » est le complément : le verbe reste au singulier dans la construction impersonnelle classique : \esp{se habla más de doscientas lenguas} (la forme plurielle \esp{se hablan} est aussi admise par l'usage moderne, mais le singulier est la forme scolaire sûre).
\item Deuxième « on » : « on est très fier » = nous, les Camerounais (le locuteur s'inclut) $\rightarrow$ 1\textsuperscript{re} personne du pluriel de \esp{estar} : \esp{estamos muy orgullosos}. On choisit \esp{estar} (fierté ressentie, état) et \esp{muy} devant l'adjectif.
\item Traduction complète : \esp{En Camerún se habla más de doscientas lenguas, y estamos muy orgullosos de esta riqueza.}
\end{itemize}
\end{exoresolu}

\begin{exemplebox}[Autres exemples traduits]
\esp{Aquí se habla francés e inglés.} « Ici on parle français et anglais. » (impersonnel ; notez \esp{e} devant \esp{inglés}, et non \esp{y})\par
\esp{En España se cena muy tarde.} « En Espagne, on dîne très tard. »\par
\esp{On m'a téléphoné hier soir.} $\rightarrow$ \esp{Me llamaron por teléfono anoche.} (quelqu'un : 3\textsuperscript{e} pluriel)\par
\esp{Uno nunca sabe lo que puede pasar.} « On ne sait jamais ce qui peut arriver. » (chacun)\par
\esp{La gente piensa que el fútbol es solo un juego.} « On pense que le football est seulement un jeu. »\par
\esp{¿Qué hacemos esta noche?} « Qu'est-ce qu'on fait ce soir ? » (nous)
\end{exemplebox}

\begin{exercice}[À vous de jouer]
Traduisez en espagnol, en justifiant chaque choix : 1) « On mange bien dans ce restaurant. » 2) « On a annoncé les résultats du baccalauréat. » 3) « Quand on voyage, on apprend beaucoup. » 4) « On va au stade samedi, tu viens avec nous ? » 5) « On dit que le cacao va augmenter. »
\end{exercice}

\begin{corrige}[Exercice « À vous de jouer »]
1) \esp{En este restaurante se come bien.} (impersonnel : les gens en général). 2) \esp{Anunciaron los resultados del bachillerato.} (quelqu'un, autorités non nommées : 3\textsuperscript{e} pluriel). 3) \esp{Cuando uno viaja, aprende mucho.} (chacun : \esp{uno}). 4) \esp{Vamos al estadio el sábado, ¿vienes con nosotros?} (nous : 1\textsuperscript{re} pluriel). 5) \esp{Se dice que el cacao va a subir.} ou \esp{La gente dice que el cacao va a subir.} (rumeur générale).
\end{corrige}

\begin{aretenir}[L'essentiel de cette première partie]
Le français n'a qu'un mot, « on » ; l'espagnol en a cinq constructions. Le choix dépend exclusivement du \cle{sens} : impersonnel (\esp{se} + 3\textsuperscript{e} sg), personnes inconnues (3\textsuperscript{e} pluriel), chacun (\esp{uno}), les gens (\esp{la gente} + singulier), nous (1\textsuperscript{re} pluriel). La question-guide à se poser est toujours : qui fait l'action ?
\end{aretenir}

\coursec{Les traductions du « on » : le « se » impersonnel et autres formes}

\courssub{Transition : du « on » français aux visages espagnols}

Dans la section précédente, nous avons observé que le pronom français \cle{on} est un caméléon grammatical : selon le contexte, il désigne une généralité (\emph{« On dit que… »}), un groupe précis dont le locuteur fait partie (\emph{« On y va ? »}), ou une personne indéterminée (\emph{« On a frappé à la porte »}). En espagnol, il n'existe \textbf{pas} de pronom unique équivalent. La construction la plus fréquente pour rendre l'impersonnalité, la généralité ou la voix passive réflexive est la structure \textbf{\cle{se} + verbe à la 3\up{e} personne}. C'est le pilier central de cette leçon : la maîtriser, c'est débloquer des centaines de phrases du quotidien, des panneaux de signalisation aux articles de presse.

\courssub{Observation : le « se » impersonnel en action}

Lisons d'abord un court texte authentique, typique de ce que l'on peut lire dans un journal local ou sur une affiche municipale à Yaoundé, Douala, Madrid ou Malabo. Repérez toutes les occurrences de \esp{se} + verbe.

\begin{textebox}[Extrait d'un bulletin municipal : « Mejora de la vida en el barrio »]
\emph{En el barrio de El Prado, \textbf{se vive} con tranquilidad. Desde el año pasado, \textbf{se han abierto} dos nuevos parques infantiles y \textbf{se ha renovado} el mercado central. Ahora \textbf{se venden} productos ecológicos todos los sábados y \textbf{se organizan} talleres de reciclaje para los niños. \textbf{Se prohíbe} tirar basura fuera de los contenedores y \textbf{se ruega} respetar las zonas verdes. Pronto \textbf{se construirá} un centro cultural y \textbf{se ampliará} la biblioteca. \textbf{Se busca} personal para la nueva ludoteca: \textbf{se ofrece} contrato indefinido. ¡Ven, aquí \textbf{se vive} bien!}
\end{textebox}

\textbf{Glossaire :} \esp{barrio} (quartier), \esp{tranquilidad} (tranquillité), \esp{parque infantil} (aire de jeux), \esp{renovar} (rénover), \esp{productos ecológicos} (produits bio), \esp{taller} (atelier), \esp{reciclaje} (recyclage), \esp{contenedor} (conteneur/poubelle), \esp{zonas verdes} (espaces verts), \esp{ampliar} (agrandir), \esp{ludoteca} (ludothèque/centre de jeux), \esp{personal} (personnel), \esp{contrato indefinido} (CDI / contrat à durée indéterminée).

\courssub{Analyse de la structure : \cle{se} + verbe (3\up{e} personne)}

Dans tous les exemples ci-dessus, le sujet grammatical est \textbf{absent} ou \textbf{indéterminé}. L'action est présentée comme un fait général, une règle, ou un événement dont l'auteur importe peu. C'est la définition même de l'\textbf{impersonnel}.

\begin{definition}[Construction impersonnelle avec \textit{se}]
La structure \textbf{\textit{se} + verbe conjugué à la 3\up{e} personne} (singulier ou pluriel) permet d'exprimer une action sans sujet déterminé. Elle correspond souvent au français \textbf{« on + verbe »}, à la \textbf{voix passive} (\emph{« être + participe »}) ou à une tournure impersonnelle (\emph{« il faut / il est + adjectif »}).
\end{definition}

\begin{popfigure}
\begin{tikzpicture}[
    node distance=1.2cm and 1.5cm,
    box/.style={draw=popblue, thick, rounded corners, fill=popblueL, minimum height=1.2cm, align=center, font=\small},
    arrow/.style={-Stealth, thick, popdark},
    label/.style={font=\tiny\itshape, text=popmuted, above}
]
\node[box, align=center] (se){\textbf{SE}\\(particule invariable)};
\node[box, right=of se, align=center] (verbe){\textbf{VERBE}\\(3\up{e} pers. sg \textbf{ou} pl)};
\node[box, right=of verbe, align=center] (compl){\textbf{COMPLÉMENT / NOM}\\(sujet patient)};
\node[box, below=of verbe, fill=poporangeL, draw=poporange, align=center] (accord){\textbf{RÈGLE D'ACCORD}\\Si le nom est \textbf{pluriel} $\rightarrow$ verbe au \textbf{pluriel}\\Si le nom est \textbf{singulier} ou absent $\rightarrow$ verbe au \textbf{singulier}};

\draw[arrow] (se) -- node[label] {invariable} (verbe);
\draw[arrow] (verbe) -- node[label] {accord possible} (compl);
\draw[arrow, poporange] (verbe.south) -- (accord.north);

\node[below=0.3cm of se, font=\footnotesize, text=popdark] {\esp{Se} / \esp{Se}};
\node[below=0.3cm of verbe, font=\footnotesize, text=popdark] {\esp{vive} / \esp{venden}};
\node[below=0.3cm of compl, font=\footnotesize, text=popdark] {\esp{---} / \esp{frutas}};
\end{tikzpicture}
\legende{Schéma de la structure impersonnelle \textit{se} + verbe + accord conditionnel}
\end{popfigure}

\courssub{La règle d'or : l'accord du verbe}

C'est le point \textbf{le plus testé au baccalauréat} et la source n°1 d'erreurs pour les francophones. En français, \emph{on} est toujours suivi du verbe au singulier (\emph{on vend, on vendent} $\times$). En espagnol, le verbe \textbf{s'accorde avec le nom qui suit} (le « sujet patient » ou complément d'objet direct devenu sujet grammatical de la passive réflexive).

\begin{propriete}[Accord du verbe dans la construction \textit{se} impersonnel / passif réflexif]
\begin{itemize}
    \item \textbf{Verbe intransitif} (pas de nom après) $\rightarrow$ \textbf{Toujours singulier}. \\
    \esp{Se vive bien en Camerún.} (On vit bien au Cameroun.)
    \item \textbf{Verbe transitif + Nom singulier / indénombrable} $\rightarrow$ \textbf{Singulier}. \\
    \esp{Se vende pan fresco.} (On vend du pain frais / Du pain frais est vendu.)
    \item \textbf{Verbe transitif + Nom pluriel} $\rightarrow$ \textbf{Pluriel}. \\
    \esp{Se venden frutas tropicales.} (On vend des fruits tropicaux / Des fruits tropicaux sont vendus.)
\end{itemize}
\end{propriete}

\begin{center}
\begin{tabularx}{\linewidth}{|X|X|X|}
\hline
\rowcolor{popblueL}
\textbf{Français (On + verbe sg)} & \textbf{Espagnol (Se + verbe)} & \textbf{Analyse} \\
\hline
On cherche un professeur. & \esp{Se busca un profesor.} & Nom sg $\rightarrow$ verbe sg \\
\hline
On cherche des professeurs. & \esp{Se buscan profesores.} & Nom pl $\rightarrow$ verbe \textbf{pl} \\
\hline
On mange bien ici. & \esp{Se come bien aquí.} & Intransitif $\rightarrow$ verbe sg \\
\hline
On a construit l'école. & \esp{Se construyó la escuela.} & Passé, nom sg $\rightarrow$ verbe sg \\
\hline
On a construit des écoles. & \esp{Se construyeron escuelas.} & Passé, nom pl $\rightarrow$ verbe \textbf{pl} \\
\hline
\end{tabularx}
\end{center}

\courssub{Emplois principaux : généralités, consignes, annonces, passif}

La construction \textit{se} impersonnel couvre trois grands domaines fonctionnels. Les distinguer aide à choisir le bon temps et le bon registre en traduction.

\begin{aretenir}[Trois emplois majeurs du \textit{se} impersonnel]
\begin{enumerate}
    \item \textbf{Généralités / Vérités générales} (Présent, Imparfait, Futur) : \emph{« On dit que… », « On vit bien… »} \\
    \esp{En España \textbf{se come} tarde.} (En Espagne, on mange tard.) \\
    \esp{Antiguamente \textbf{se vivía} del campo.} (Autrefois, on vivait de la terre.)
    \item \textbf{Consignes, interdictions, panneaux, annonces} (Présent, souvent infinitif \textit{se} + infinitif pour l'affichage) : \emph{« Il est interdit de… », « On demande de… »} \\
    \esp{\textbf{Se prohíbe} fumar.} (Il est interdit de fumer / Défense de fumer.) \\
    \esp{\textbf{Se ruega} el silencio.} (On prie de faire silence / Silence, s'il vous plaît.) \\
    \textit{Affiche :} \esp{Se alquila / Se vende / Se buscan.}
    \item \textbf{Passif réflexif (équivalent du passif français « être + participe »)} : l'accent est mis sur l'action subie par l'objet. Tous les temps sont possibles. \\
    \esp{La escuela \textbf{se construyó} en 2020.} (L'école a été construite en 2020.) \\
    \esp{El puente \textbf{se construirá} el año que viene.} (Le pont sera construit l'an prochain.) \\
    \esp{Los libros \textbf{se han vendido} ya.} (Les livres ont déjà été vendus.)
\end{enumerate}
\end{aretenir}

\courssub{Le « se » impersonnel à tous les temps}

La particule \textit{se} est \textbf{invariable}. Elle ne change jamais de forme. C'est le verbe qui porte la marque du temps et du nombre. Voici le paradigme complet avec le verbe \esp{construir} (construire) + nom pluriel \esp{puentes} (ponts) pour bien voir l'accord au pluriel maintenu à tous les temps.

\begin{center}
\begin{tabular}{|l|l|l|}
\hline
\rowcolor{popblueL}
\textbf{Temps / Mode} & \textbf{Forme (Nom pluriel)} & \textbf{Traduction française} \\
\hline
Présent de l'indicatif & \esp{Se construyen puentes.} & On construit des ponts / Des ponts sont construits. \\
\hline
Imparfait & \esp{Se construían puentes.} & On construisait des ponts / Des ponts étaient construits. \\
\hline
Passé simple & \esp{Se construyeron puentes.} & On a construit des ponts / Des ponts furent construits. \\
\hline
Futur simple & \esp{Se construirán puentes.} & On construira des ponts / Des ponts seront construits. \\
\hline
Conditionnel & \esp{Se construirían puentes.} & On construirait des ponts / Des ponts seraient construits. \\
\hline
Présent perfecto & \esp{Se han construido puentes.} & On a construit des ponts / Des ponts ont été construits. \\
\hline
Plus-que-parfait & \esp{Se habían construido puentes.} & On avait construit des ponts / Des ponts avaient été construits. \\
\hline
Subjonctif présent & \esp{Es bueno que \textbf{se construyan} puentes.} & Il est bien qu'on construise des ponts. \\
\hline
Infinitif (après prép./affiches) & \esp{Antes de \textbf{se construir} los puentes...} & Avant de construire les ponts... \\
\hline
Gérondif & \esp{Al \textbf{se estar construyendo} los puentes...} & En construisant les ponts / Pendant qu'on construit les ponts. \\
\hline
\end{tabular}
\end{center}

\courssub{Distinction cruciale : « Se » impersonnel vs « Se » réflexif}

C'est un piège classique. La forme \esp{se lava} existe dans les deux cas, mais le sens et la structure syntaxique diffèrent radicalement.

\begin{attention}[Piège n°1 : Confondre impersonnel et réflexif]
\begin{itemize}
    \item \textbf{Réfléchi :} Le sujet \textbf{fait l'action sur lui-même}. Il y a un sujet (yo, tú, él, nosotros...). \\
    \esp{Él \textbf{se lava} las manos.} (Il se lave les mains.) $\leftarrow$ Sujet = \esp{Él}.
    \item \textbf{Impersonnel / Passif :} \textbf{Pas de sujet déterminé}. L'action est générale ou subie par l'objet. \\
    \esp{\textbf{Se lava} la ropa en el río.} (On lave le linge à la rivière / Le linge se lave à la rivière.) $\leftarrow$ Pas de sujet, \esp{la ropa} = COD devenu sujet grammatical (accord sg).
\end{itemize}
\textbf{Test de détection :} Essayez de mettre \emph{« par quelqu'un »} ou \emph{« on »} devant le verbe français.
\begin{itemize}
    \item \esp{Él se lava} $\rightarrow$ *Il se lave par quelqu'un* ? \textbf{Non}. C'est réflexif.
    \item \esp{Se lava la ropa} $\rightarrow$ *On lave le linge / Le linge est lavé (par quelqu'un)* ? \textbf{Oui}. C'est impersonnel/passif.
\end{itemize}
\end{attention}

\courssub{Exercice résolu : Thème et Version}

\begin{exoresolu}[Traduction contextuelle : « On » $\rightarrow$ \textit{se} impersonnel]
\textbf{Énoncé :} Traduisez en espagnol les phrases suivantes en utilisant la construction \textit{se} impersonnel / passif réflexif. Justifiez l'accord du verbe (sg/pl).

1. On parle espagnol en Guinée équatoriale. (intransitif)
2. On a vendu la maison. (passé simple, nom sg)
3. On vend des mangues au marché Central. (présent, nom pl)
4. Il est interdit de stationner ici. (panneau / interdiction)
5. On construira un nouveau stade à Douala. (futur, nom sg)
6. Autrefois, on cultivait du cacao dans cette région. (imparfait, indénombrable)

\tcblower
\textbf{Solution détaillée :}

1. \esp{En Guinea Ecuatorial \textbf{se habla} español.}
   \textit{Justification :} Verbe \esp{hablar} intransitif (pas de COD après) $\rightarrow$ \textbf{singulier}.

2. \esp{\textbf{Se vendió} la casa.}
   \textit{Justification :} Passif réflexif passé simple. \esp{La casa} = sujet patient \textbf{singulier} $\rightarrow$ verbe \textbf{singulier} (\esp{vendió}, pas *vendieron*).

3. \esp{En el mercado Central \textbf{se venden} mangos.}
   \textit{Justification :} Verbe transitif \esp{vender} + \esp{mangos} (\textbf{pluriel}) $\rightarrow$ verbe au \textbf{pluriel} (\esp{venden}). \textbf{Erreur fréquente :} *Se vende mangos*.

4. \esp{\textbf{Se prohíbe} aparcar aquí.} (ou \esp{Se prohíbe el estacionamiento.})
   \textit{Justification :} Panneau / Interdiction. \esp{aparcar} (infinitif = idée sg) ou \esp{el estacionamiento} (nom sg) $\rightarrow$ \textbf{singulier}.

5. \esp{En Douala \textbf{se construirá} un nuevo estadio.}
   \textit{Justification :} Futur simple. \esp{Un nuevo estadio} = \textbf{singulier} $\rightarrow$ verbe \textbf{singulier} (\esp{construirá}).

6. \esp{Antiguamente \textbf{se cultivaba} cacao en esta región.}
   \textit{Justification :} Imparfait. \esp{Cacao} = indénombrable (masculin singulier) $\rightarrow$ verbe \textbf{singulier}.
\end{exoresolu}

\courssub{Le saviez-vous ? : Lire la ville hispanophone}

\begin{savaistu}[« Se busca », « Se vende », « Se alquila » : l'espagnol de la rue]
En vous promenant dans les rues de \textbf{Madrid}, \textbf{Buenos Aires}, \textbf{Mexico} ou \textbf{Malabo} (Guinée équatoriale), vous verrez des centaines de pancartes sur les balcons, les vitrines et les poteaux électriques. Elles utilisent \textbf{exclusivement} le \textit{se} impersonnel :
\begin{itemize}
    \item \esp{\textbf{SE VENDE}} (À vendre) — souvent suivi de \esp{piso} (appartement), \esp{coche} (voiture), \esp{terreno} (terrain).
    \item \esp{\textbf{SE ALQUILA}} (À louer) — \esp{apartamento}, \esp{local} (local commercial).
    \item \esp{\textbf{SE BUSCA}} (On cherche / Recherché) — \esp{personal} (serveurs, cuisiniers), \esp{compañero de piso} (colocataire).
    \item \esp{\textbf{SE NECESITA}} (On a besoin de) — \esp{ayuda}, \esp{voluntarios}.
\end{itemize}
\textbf{Astuce de terrain :} Si le nom est au pluriel, l'affiche correcte écrira \esp{SE VENDEN} (ex: \esp{Se venden entradas} — On vend des billets). Repérer l'accord sur les panneaux est le meilleur exercice de grammaire \emph{in situ} qui soit !
\end{savaistu}

\courssub{Exercices d'application immédiate}

\begin{exercice}[Mise en œuvre : Choisir la bonne forme]
Mettez le verbe entre parenthèses à la forme correcte (indicatif présent, sauf indication) en respectant l'accord.

1. En este restaurante \_\_\_\_\_\_\_\_\_\_ (comer) muy bien.
2. \_\_\_\_\_\_\_\_\_\_ (vender) entradas para el concierto en la taquilla.
3. \_\_\_\_\_\_\_\_\_\_ (prohibir) la entrada a los menores de 18 años.
4. El año pasado \_\_\_\_\_\_\_\_\_\_ (construir) muchas casas en la periferia. (Passé simple)
5. Para la fiesta \_\_\_\_\_\_\_\_\_\_ (necesitar) muchas sillas y mesas.
6. \_\_\_\_\_\_\_\_\_\_ (buscar) profesores de francés para el liceo. (Présent)
7. Aquí \_\_\_\_\_\_\_\_\_\_ (vivir) con mucha alegría.
8. Los exámenes \_\_\_\_\_\_\_\_\_\_ (corregir) ya por el profesor. (Présent perfecto)
\end{exercice}

\begin{corrige}[Exercice : Mise en œuvre]
1. \esp{se come} (intransitif $\rightarrow$ sg).
2. \esp{Se venden} (\esp{entradas} pl $\rightarrow$ pl).
3. \esp{Se prohíbe} (\esp{la entrada} sg $\rightarrow$ sg).
4. \esp{se construyeron} (\esp{muchas casas} pl $\rightarrow$ pl, passé simple).
5. \esp{se necesitan} (\esp{muchas sillas y mesas} pl $\rightarrow$ pl).
6. \esp{Se buscan} (\esp{profesores} pl $\rightarrow$ pl).
7. \esp{se vive} (intransitif $\rightarrow$ sg).
8. \esp{se han corregido} (\esp{Los exámenes} pl $\rightarrow$ pl, présent perfecto).
\end{corrige}

\begin{exercice}[Thème : Du français vers l'espagnol (niveau Bac)]
Traduisez en espagnol en utilisant obligatoirement la construction \textit{se} impersonnel/passif.

1. On a trouvé une solution au problème. (Passé simple)
2. Il faut respecter les règles. (Indication : \esp{Se deben respetar las normas.})
3. Dans ce pays, on produit beaucoup de café. (Présent)
4. On cherche des témoins de l'accident. (Présent)
5. Le match a été annulé à cause de la pluie. (Passé simple, \esp{suspender} ou \esp{anular})
6. Autrefois, on n'allait pas à l'école. (Imparfait, \esp{ir} $\rightarrow$ \esp{irse} ou \esp{asistir})
\end{exercice}

\begin{corrige}[Exercice : Thème]
1. \esp{Se encontró una solución al problema.}
2. \esp{Se deben respetar las normas.} (\esp{las normas} pl $\rightarrow$ \esp{se deben}). \textit{Variante :} \esp{Se exige el respeto de las normas.}
3. \esp{En este país se produce mucho café.} (\esp{café} sg indénombrable).
4. \esp{Se buscan testigos del accidente.} (\esp{testigos} pl).
5. \esp{El partido se suspendió / se anuló a causa de la lluvia.} (\esp{El partido} sg).
6. \esp{Antiguamente no se iba a la escuela / no se asistía a la escuela.} (Impersonnel intransitif $\rightarrow$ sg).
\end{corrige}

\coursec{Les autres traductions de « on » : \textit{tercera persona del plural}, \textit{uno}, \textit{la gente}, \textit{nosotros}}

\courssub{La 3\up{e} personne du pluriel : « ils / on » indéterminé}

Lorsque le « on » français désigne des personnes non identifiées mais \textbf{concrètes} (pas une généralité abstraite), l'espagnol utilise souvent la 3\up{e} personne du pluriel (\textbf{ellos}, omis le plus souvent) à l'actif. C'est l'équivalent de « \emph{they} » en anglais ou de « \emph{ils} » en français oral (\emph{« Ils disent que... »}).

\begin{propriete}[3\up{e} personne du pluriel (sujet omis) pour « on » indéterminé actif]
Emploi : Action réalisée par des gens non précisés, rumeurs, ou « on » = « les gens / ils ».
\begin{itemize}
    \item \esp{\textbf{Dicen} que va a llover.} (On dit qu'il va pleuvoir / \emph{Ils disent} qu'il va pleuvoir.)
    \item \esp{\textbf{Han cortado} el agua en el barrio.} (On a coupé l'eau dans le quartier / \emph{Ils ont coupé} l'eau.)
    \item \esp{\textbf{Llaman} a la puerta.} (On frappe à la porte / \emph{Quelqu'un frappe}.)
\end{itemize}
\textbf{Différence avec \textit{se} :} \textit{Se} + verbe = passif/impersonnel (focus sur l'action/objet). 3\up{e} pl actif = focus sur l'agent (même inconnu).
\end{propriete}

\begin{attention}[Piège n°2 : Ne pas confondre « On a coupé l'eau » (événement) et « L'eau a été coupée » (état/résultat)]
\begin{itemize}
    \item \esp{\textbf{Han cortado} el agua.} (Action ponctuelle, agents concrets : les techniciens.) $\rightarrow$ 3\up{e} pl actif.
    \item \esp{\textbf{Se ha cortado} el agua.} (Résultat, état, généralité : l'eau est coupée.) $\rightarrow$ \textit{se} passif/impersonnel.
\end{itemize}
Au Bac, si le sujet français est « on » sans précision et que le verbe est transitif avec un objet précis $\rightarrow$ hésitez entre \textit{se} + accord (passif) et 3\up{e} pl (actif). Le contexte décide.
\end{attention}

\courssub{\textit{Uno} : l'indéfini formel}

\cle{Uno} (apocope \cle{un} devant nom masc. sg) est un pronom indéfini formel, équivalent à « \emph{on} » au sens de « \emph{n'importe qui} », « \emph{une personne quelconque} ». Il s'accorde en genre (\textit{una}) et déclenche le verbe à la 3\up{e} personne du singulier. Il est plus soutenu que \textit{se}.

\begin{definition}[\textit{Uno} / \textit{Una} (pronom indéfini)]
\textbf{Forme :} \esp{Uno} (masc.), \esp{Una} (fém.). \textbf{Verbe :} 3\up{e} pers. sg. \textbf{Possessifs :} \esp{su / sus} (pas *suo*).
\end{definition}

\begin{exemplebox}[Exemples avec \textit{uno}]
\begin{itemize}
    \item \esp{\textbf{Uno} no debe mentir nunca.} (On / Quelqu'un ne doit jamais mentir.) — Règle morale, généralité.
    \item \esp{Si \textbf{uno} estudia, \textbf{uno} aprende.} (Si on étudie, on apprend.) — Conditionnel général.
    \item \esp{\textbf{Una} se siente sola a veces.} (On / Une femme se sent seule parfois.) — Accord au féminin.
\end{itemize}
\end{exemplebox}

\begin{attention}[Piège n°3 : \textit{Uno} $\neq$ \textit{Un} (chiffre) / \textit{Un} (article)]
\begin{itemize}
    \item \esp{\textbf{Uno} piensa...} (Pronom sujet = On.)
    \item \esp{Tengo \textbf{un} libro.} (Article indéfini = Un.)
    \item \esp{¿Cuántos hay? \textbf{Uno}.} (Chiffre = Un.)
\end{itemize}
\textbf{Accord des adjectifs avec \textit{uno} :} \esp{Uno solo} (masc.), \esp{Una sola} (fém.). \emph{Jamais} \esp{*uno sola}.
\end{attention}

\courssub{\textit{La gente} : le peuple, les gens (collectif singulier)}

\cle{La gente} (« les gens », « le peuple ») est un nom collectif \textbf{singulier} en espagnol. Le verbe est au singulier. C'est l'équivalent le plus naturel du « on » familier ou journalistique désignant la population, l'opinion publique, « tout le monde ».

\begin{propriete}[\textit{La gente} + verbe au singulier]
\begin{itemize}
    \item \esp{\textbf{La gente} \textbf{dice} que el examen fue difícil.} (Les gens disent / On dit que l'examen était difficile.)
    \item \esp{En Camerún, \textbf{la gente} \textbf{come} mucho plátano.} (Au Cameroun, les gens / on mange beaucoup de plantain.)
    \item \esp{\textbf{La gente joven} \textbf{usa} mucho el móvil.} (Les jeunes / La jeunesse utilise beaucoup le portable.)
\end{itemize}
\textbf{Attention :} Les adjectifs/participes qui se rapportent à \textit{la gente} s'accordent au \textbf{féminin singulier} (car \textit{gente} est fém. sg), mais si le sens est pluriel (les individus), l'accord peut se faire au masculin pluriel (syllepse) : \esp{La gente están contentos} (accepté) ou \esp{La gente está contenta} (strictement grammatical).
\end{propriete}

\courssub{\textit{Nosotros} : le « on » inclusif (nous)}

Dans le langage familier ou oral, « on » = « nous » (le locuteur inclus). L'espagnol utilise \cle{nosotros} (ou \cle{nosotras}) + verbe à la 1\up{e} personne du pluriel. C'est la traduction la plus directe pour les propositions, décisions partagées.

\begin{exemplebox}[\textit{Nosotros} = « On » inclusif]
\begin{itemize}
    \item \emph{« On y va ? »} $\rightarrow$ \esp{\textbf{¿Vamos?}} / \esp{\textbf{¿Nos vamos?}} (Verbe 1\up{e} pl, pronom omis ou \esp{nos}).
    \item \emph{« On a gagné ! »} $\rightarrow$ \esp{\textbf{Hemos ganado!}} / \esp{\textbf{Nosotros hemos ganado!}}
    \item \emph{« On se voit demain ? »} $\rightarrow$ \esp{\textbf{¿Nos vemos mañana?}} (\textit{nos} = réciproque « nous voir mutuellement »).
\end{itemize}
\end{exemplebox}

\courssub{Tableau récapitulatif : choisir la bonne traduction de « on »}

\begin{center}
\begin{tabularx}{\linewidth}{|X|X|X|X|}
\hline
\rowcolor{popblueL}
\textbf{Contexte / Sens du « on »} & \textbf{Structure espagnole} & \textbf{Exemple} & \textbf{Traduction} \\
\hline
Généralité, vérité, passif, panneau & \textbf{Se + verbe (accord)} & \esp{Se habla español.} & On parle espagnol. \\
\hline
Action par agents indéterminés (rumeur, fait divers) & \textbf{3\up{e} pers. pl (sujet omis)} & \esp{Han subido los precios.} & On a augmenté les prix. \\
\hline
« N'importe qui », règle morale, formel & \textbf{Uno / Una + verbe sg} & \esp{Uno debe respetar.} & On doit respecter. \\
\hline
« Les gens », opinion publique, population & \textbf{La gente + verbe sg} & \esp{La gente piensa que...} / \esp{La gente están locos.} & Les gens pensent que... / On est fous. \\
\hline
« Nous » (locuteur inclus), proposition & \textbf{Nosotros / Verbe 1\up{e} pl} & \esp{¿Vamos al cine?} & On va au ciné ? \\
\hline
\end{tabularx}
\end{center}

\courssub{Exercices de synthèse : « On » tous azimuts}

\begin{exercice}[Choix de la structure : \textit{se}, 3\up{e} pl, \textit{uno}, \textit{la gente}, \textit{nosotros}]
Traduisez en espagnol en justifiant votre choix structurel.

1. On a volé mon portefeuille dans le bus. (Action délictueuse, agents inconnus)
2. On ne doit pas mentir. (Règle morale générale, formel)
3. Au marché, on vend des ananas. (Activité commerciale habituelle)
4. Les gens disent qu'il fera chaud demain. (Opinion publique)
5. On part en vacances la semaine prochaine ? (Proposition entre amis)
6. Dans ce pays, on vit vieux. (Généralité statistique)
7. On a livré le colis hier. (Action accomplie par le livreur)
8. En Afrique, on mange avec les doigts. (Coutume culturelle générale)
\end{exercice}

\begin{corrige}[Exercice : Synthèse « On »]
1. \esp{Me \textbf{robaron} la cartera en el autobús.} (3\up{e} pl actif : agents concrets, vol). \textit{Alternative passif :} \esp{Me \textbf{robaron}...} est plus naturel que \esp{Se robó...} pour un vol précis.
2. \esp{\textbf{Uno} no debe mentir.} (\textit{Uno} = indéfini formel, règle morale). \textit{Alternative :} \esp{No \textbf{se debe} mentir.}
3. \esp{En el mercado \textbf{se venden} piñas.} (\textit{Se} passif/impersonnel + accord pl \textit{piñas}).
4. \esp{\textbf{La gente dice} que hará calor mañana.} (\textit{La gente} = opinion publique).
5. \esp{\textbf{¿Nos vamos} de vacaciones la semana que viene?} (\textit{Nosotros} inclusif / proposition).
6. \esp{En este país \textbf{se vive} mucho.} (\textit{Se} intransitif = généralité).
7. \esp{\textbf{Han entregado} el paquete ayer.} (3\up{e} pl actif : action précise du livreur). \textit{Alternative :} \esp{El paquete \textbf{se entregó} ayer.} (Passif focus sur colis).
8. \esp{En África \textbf{se come} con los dedos.} (\textit{Se} intransitif = coutume générale). \textit{Variante :} \esp{\textbf{La gente come} con los dedos.}
\end{corrige}

\coursec{Les autres traductions de « on » : 3e personne du pluriel, uno, la gente, nosotros}

Dans la section précédente, nous avons vu que le « on » français se traduit très souvent par la construction impersonnelle \esp{se + verbe à la 3e personne du singulier} : \esp{En Francia se bebe mucho vino} « En France, on boit beaucoup de vin ». Mais le « se » impersonnel n'est pas la seule solution. Le français possède un seul petit mot, « on », là où l'espagnol en possède \textbf{cinq}. Le bon traducteur est donc celui qui sait \cle{choisir} la bonne équivalence selon le sens exact de « on » dans la phrase. C'est l'objet de cette section.

\courssub{Observation : un même « on », quatre visages espagnols}

Lisons d'abord ce court texte de presse original, puis observons comment chaque « on » français serait rendu en espagnol.

\begin{textebox}[Crónica : Un robo en el barrio]
Ayer por la noche me robaron la moto delante de mi casa, en pleno centro de Yaoundé. Nadie vio nada, o nadie quiso ver nada. La gente dice que el barrio ya no es seguro, que antes uno podía dejar la puerta abierta sin miedo. Una vecina me contó que a ella también le robaron el teléfono la semana pasada, en el mercado. « Uno nunca sabe dónde está el peligro », me dijo con tristeza. Los vecinos piensan organizar rondas nocturnas para vigilar las calles. Mi hermano y yo hemos decidido partir mañana temprano a la comisaría para declarar el robo. La gente piensa que la policía no hace bastante, pero nosotros creemos que hay que colaborar con ella. Al final, uno comprende que la seguridad es asunto de todos : si cada uno cuida de su vecino, el barrio entero vive mejor.
\end{textebox}

\textbf{Glossaire} : \esp{robar} : voler, dérober ; \esp{en pleno centro} : en plein centre ; \esp{seguro/a} : sûr, sécurisé ; \esp{el peligro} : le danger ; \esp{las rondas nocturnas} : les rondes de nuit ; \esp{vigilar} : surveiller ; \esp{la comisaría} : le commissariat ; \esp{declarar el robo} : déclarer le vol ; \esp{asunto de todos} : l'affaire de tous ; \esp{cuidar de} : prendre soin de ; \esp{madrugar} : se lever tôt (utilisé dans les exercices).

\textbf{Questions de compréhension} :
\begin{enumerate}
\item ¿Qué le robaron al narrador y dónde ?
\item ¿Qué piensa la gente del barrio ?
\item ¿Qué han decidido hacer el narrador y su hermano ?
\end{enumerate}

Observons maintenant les équivalences cachées dans ce texte :

\begin{center}
\begin{tabularx}{\linewidth}{|X|X|}
\hline
\rowcolor{popblueL} \textbf{Français (avec « on »)} & \textbf{Espagnol du texte} \\
\hline
« On m'a volé ma moto » & \esp{Me robaron la moto} \\
\hline
« On dit que le quartier n'est plus sûr » & \esp{La gente dice que el barrio ya no es seguro} \\
\hline
« On ne sait jamais où est le danger » & \esp{Uno nunca sabe dónde está el peligro} \\
\hline
« On part demain matin » (mon frère et moi) & \esp{Partimos mañana temprano} \\
\hline
\end{tabularx}
\end{center}

Quatre occurrences de « on », quatre constructions espagnoles différentes. Analysons-les une à une.

\courssub{La 3e personne du pluriel : des agents réels mais non identifiés}

\begin{definition}[La 3e personne du pluriel à valeur indéfinie]
En espagnol, lorsque l'action est accomplie par des \cle{personnes réelles mais non identifiées} (des voleurs, des inconnus, « les gens » en général, les autorités...), on utilise le verbe à la \textbf{3e personne du pluriel sans sujet exprimé}. Le sujet sous-entendu est « ils », « les gens », « certains ».
\end{definition}

\begin{propriete}[Forme et emploi]
\textbf{Forme} : verbe à la 3e personne du pluriel, \emph{sans} pronom sujet, \emph{sans} « se ».

\textbf{Emploi} : l'action est attribuée à des agents précis dans la réalité (on sait que \emph{quelqu'un} a agi), mais on ne les nomme pas, parce qu'on les ignore ou qu'ils sont sans importance.

\esp{Me robaron la moto.} « On m'a volé ma moto. »

\esp{Me telefonearon ayer por la noche.} « On m'a téléphoné hier soir. »

\esp{Le dijeron que el examen era el lunes.} « On lui a dit que l'examen était lundi. »

\esp{En el mercado venden cacao de muy buena calidad.} « Au marché, on vend du cacao de très bonne qualité. »

\esp{Han construido un nuevo estadio en Douala.} « On a construit un nouveau stade à Douala. »
\end{propriete}

\begin{attention}[3e pluriel ou « se » impersonnel ? La nuance décisive]
Les deux constructions traduisent « on », mais la nuance est différente :
\begin{itemize}
\item \textbf{3e pluriel} : on pense à des \emph{personnes} (inconnues) qui ont agi : \esp{Me robaron la moto} (des voleurs bien réels !).
\item \textbf{Se impersonnel} : on décrit un \emph{fait général}, une coutume, sans penser à personne en particulier : \esp{En Francia se bebe mucho vino} « En France, on boit beaucoup de vin ».
\end{itemize}
\textbf{Attention} : \esp{Se me robaron la moto} n'est \emph{pas} la traduction de « On m'a volé ma moto » (voix active). C'est une structure \emph{passive réflexe} + \emph{datif d'intérêt} : « Ma moto s'est fait voler (sur moi) ». Pour traduire l'actif « On a volé ma moto », on utilise \textbf{uniquement} la 3e personne du pluriel active : \esp{Me robaron la moto}.
\end{attention}

\courssub{« Uno / una » : la vérité générale et l'expérience universalisée}

\begin{definition}[Le pronom indéfini « uno »]
\esp{Uno} (féminin \esp{una}) est un pronom indéfini qui signifie « on », « quelqu'un », « chacun ». Il s'emploie quand le locuteur énonce une \cle{vérité générale} ou raconte une \cle{expérience personnelle} qu'il présente comme valable pour tout le monde. Le verbe se met à la \textbf{3e personne du singulier}.
\end{definition}

\begin{exemplebox}[« Uno » en contexte]
\esp{Uno nunca sabe lo que puede pasar.} « On ne sait jamais ce qui peut arriver. »

\esp{Uno aprende mucho viajando.} « On apprend beaucoup en voyageant. »

\esp{Cuando uno es joven, no piensa en el futuro.} « Quand on est jeune, on ne pense pas à l'avenir. »

\esp{Aquí uno se siente como en casa.} « Ici, on se sent comme à la maison. »
\end{exemplebox}

\begin{propriete}[« Una » : quand le locuteur est une femme]
Si la personne qui parle est une \textbf{femme} et qu'elle \cle{s'inclut elle-même} dans cette généralité, elle emploie la forme féminine \esp{una}, et l'adjectif ou le participe s'accorde au féminin :

\esp{Una está cansada de esperar.} « On est fatiguée d'attendre » (dit par une femme qui parle d'elle-même).

\esp{Una no puede salir sola de noche.} « On ne peut pas sortir seule la nuit » (une femme qui inclut son propre cas).

Un homme dirait dans le même cas : \esp{Uno está cansado de esperar.} / \esp{Uno no puede salir solo de noche.}
\end{propriete}

\begin{attention}[Ne pas confondre « uno » pronom et « un » article]
\esp{Uno} pronom indéfini garde toujours son \textbf{-o} final : \esp{Uno nunca sabe}. Devant un nom masculin, l'article indéfini s'apocope : \esp{un libro}. Ne traduis donc jamais « on » par \esp{*un} tout seul.
\end{attention}

\courssub{« La gente » : « on » au sens de « les gens »}

\begin{definition}[« La gente »]
\esp{La gente} (« les gens », « le monde ») traduit « on » lorsque celui-ci désigne \cle{l'opinion générale}, ce que pensent ou font les gens en général. \textbf{Attention} : c'est un nom \textbf{féminin singulier} : le verbe et les adjectifs se mettent au singulier.
\end{definition}

\begin{exemplebox}[« La gente » en contexte]
\esp{La gente piensa que la policía no hace bastante.} « On pense que la police n'en fait pas assez. »

\esp{La gente dice que el barrio ya no es seguro.} « On dit que le quartier n'est plus sûr. »

\esp{La gente es muy amable en este pueblo.} « On est très aimable dans ce village » / « Les gens sont très aimables... »

\esp{En diciembre la gente viaja mucho.} « En décembre, on voyage beaucoup. »
\end{exemplebox}

\begin{attention}[Piège majeur : « la gente » est TOUJOURS singulier]
Le français « les gens » est pluriel ; l'espagnol \esp{la gente} est \textbf{singulier}, sans exception :

\esp{La gente es amable.} et jamais \esp{*La gente son amables.}

\esp{La gente está contenta.} et jamais \esp{*La gente están contentos.}

C'est l'une des erreurs les plus fréquentes des francophones à l'examen. Retiens : \textbf{LA GENTE + verbe au SINGULIER}.
\end{attention}

\courssub{« Nosotros » : quand « on » veut dire « nous »}

\begin{definition}[« On » = « nous »]
Dans la langue courante, le français « on » remplace très souvent « nous » : « On part demain » = « Nous partons demain ». Dans ce cas, l'espagnol n'a aucun mot spécial : il conjugue simplement le verbe à la \textbf{1re personne du pluriel}, avec ou sans le pronom \esp{nosotros/nosotras}.
\end{definition}

\begin{exemplebox}[« On » = « nous »]
\esp{Partimos mañana.} « On part demain. »

\esp{Nosotros partimos mañana temprano.} « Nous, on part demain matin » (avec insistance ou opposition).

\esp{¿Vamos al cine esta noche?} « On va au cinéma ce soir ? »

\esp{El año pasado fuimos a España.} « L'année dernière, on est allés en Espagne. »

\esp{Nosotras hemos terminado el trabajo.} « On a fini le travail » (entre femmes).
\end{exemplebox}

\begin{propriete}[Le pronom sujet est facultatif]
En espagnol, la désinence verbale suffit : \esp{partimos} contient déjà l'idée de « nous ». On n'ajoute \esp{nosotros} que pour \textbf{insister}, \textbf{opposer} ou \textbf{lever une ambiguïté} : \esp{Nosotros creemos que hay que colaborar, pero ellos no.} « Nous, on pense qu'il faut collaborer, mais eux non. »
\end{propriete}

\courssub{Tableau récapitulatif et organigramme de décision}

\begin{aretenir}[Les cinq visages espagnols de « on »]
\begin{center}
\begin{tabularx}{\linewidth}{|l|X|X|}
\hline
\rowcolor{popblueL} \textbf{Équivalence} & \textbf{Critère de choix} & \textbf{Exemple} \\
\hline
\esp{se} + 3e sg & généralité, coutume, règle impersonnelle & \esp{Aquí se habla francés.} « Ici on parle français. » \\
\hline
3e personne du pluriel & agents réels mais non identifiés & \esp{Me robaron la moto.} « On m'a volé ma moto. » \\
\hline
\esp{uno / una} & vérité générale, expérience universalisée & \esp{Uno nunca sabe.} « On ne sait jamais. » \\
\hline
\esp{la gente} + sg & « on » = « les gens », opinion générale & \esp{La gente piensa que...} « On pense que... » \\
\hline
\esp{nosotros} (1re pl.) & « on » = « nous » clairement & \esp{Partimos mañana.} « On part demain. » \\
\hline
\end{tabularx}
\end{center}
\end{aretenir}

\begin{popfigure}
\begin{tikzpicture}[
  node distance=0.9cm and 1.6cm,
  question/.style={rectangle, rounded corners, draw=popblue, fill=popblueL, align=center, text width=4.2cm, inner sep=5pt},
  reponse/.style={rectangle, rounded corners, draw=popgreen, fill=popgreenL, align=center, text width=4.6cm, inner sep=5pt},
  fleche/.style={-{Stealth[length=3mm]}, thick, draw=popdark},
  oui/.style={font=\bfseries\small, text=popgreen},
  non/.style={font=\bfseries\small, text=poporange}
]
\node[question, align=center] (q1){\textbf{« On » = « nous » ?}\\(le locuteur est inclus)};
\node[reponse, right=of q1, align=center] (r1){\textbf{1re pers. du pluriel}\\\esp{Partimos mañana.}};
\node[question, below=1.1cm of q1, align=center] (q2){\textbf{« On » = « les gens » ?}\\(opinion, comportement général)};
\node[reponse, right=of q2, align=center] (r2){\textbf{\esp{la gente} + sg ou 3e pl.}\\\esp{La gente dice que...}};
\node[question, below=1.1cm of q2, align=center] (q3){\textbf{Agents réels inconnus ?}\\(vol, appel, annonce...)};
\node[reponse, right=of q3, align=center] (r3){\textbf{3e pers. du pluriel}\\\esp{Me telefonearon ayer.}};
\node[reponse, below=1.1cm of q3, align=center] (r4){\textbf{Généralité impersonnelle}\\\esp{se} + 3e sg : \esp{Se vive bien aquí.}\\ou \esp{uno} : \esp{Uno nunca sabe.}};
\draw[fleche] (q1) -- node[oui, above]{oui} (r1);
\draw[fleche] (q1.south) -- node[non, right]{non} (q2.north);
\draw[fleche] (q2) -- node[oui, above]{oui} (r2);
\draw[fleche] (q2.south) -- node[non, right]{non} (q3.north);
\draw[fleche] (q3) -- node[oui, above]{oui} (r3);
\draw[fleche] (q3.south) -- node[non, right]{non} (r4.west);
\end{tikzpicture}
\legende{Organigramme de décision : comment choisir la traduction de « on »}
\end{popfigure}

\begin{methode}[Traduire « on » en thème : la démarche en 4 questions]
\begin{enumerate}
\item \textbf{« On » = « nous » ?} (on peut le remplacer par « nous » sans changer le sens) → 1re personne du pluriel : « On part demain » → \esp{Partimos mañana}.
\item \textbf{« On » = « les gens » ?} (on parle de l'opinion ou des habitudes générales) → \esp{la gente} + singulier, ou 3e pluriel : « On pense que... » → \esp{La gente piensa que...}.
\item \textbf{L'action est-elle faite par des inconnus bien réels ?} (vol, appel, information) → 3e personne du pluriel : « On m'a téléphoné hier » → \esp{Me telefonearon ayer}.
\item \textbf{Sinon}, c'est une généralité impersonnelle → \esp{se} + 3e singulier (coutumes, règles) ou \esp{uno} (vérité générale, ton personnel) : « En France on boit beaucoup de vin » → \esp{En Francia se bebe mucho vino} ; « On ne sait jamais » → \esp{Uno nunca sabe}.
\end{enumerate}
\end{methode}

\begin{methode}[Pour la version (espagnol → français)]
À l'inverse, quand tu traduis un texte espagnol vers le français :
\begin{enumerate}
\item Repère \esp{se + verbe} sans sujet → traduis par « on » : \esp{Se come bien aquí} → « On mange bien ici ».
\item Repère un verbe à la \textbf{3e pluriel sans sujet} → traduis par « on » : \esp{Me robaron la moto} → « On m'a volé ma moto ».
\item Repère \esp{uno / una} → traduis par « on » : \esp{Uno nunca sabe} → « On ne sait jamais ».
\item \esp{La gente} peut se traduire par « on » ou « les gens » selon la phrase.
\end{enumerate}
\end{methode}

\begin{exoresolu}[Thème guidé : choisir la bonne équivalence]
Traduire en espagnol : (1) « On m'a téléphoné hier. » (2) « En France, on boit beaucoup de vin. » (3) « On part demain, ma sœur et moi. » (4) « On ne sait jamais. » (5) « On dit que ce lycée est le meilleur de la ville. »
\tcblower
\textbf{Solution détaillée} :
\begin{enumerate}
\item Des agents réels mais inconnus (quelqu'un a téléphoné) → 3e pluriel : \esp{Me telefonearon ayer.}
\item Une coutume générale, sans agents précis → \esp{se} impersonnel : \esp{En Francia se bebe mucho vino.}
\item « On » = « nous » (le locuteur et sa sœur) → 1re pluriel : \esp{Partimos mañana, mi hermana y yo.}
\item Vérité générale, ton personnel → \esp{uno} : \esp{Uno nunca sabe.}
\item Opinion générale, « les gens disent » → \esp{la gente} + singulier : \esp{La gente dice que este liceo es el mejor de la ciudad.} (ou 3e pluriel : \esp{Dicen que este liceo es el mejor...})
\end{enumerate}
\end{exoresolu}

\begin{exercice}[À ton tour]
Traduis en espagnol en justifiant ton choix :
\begin{enumerate}
\item « On a construit un pont sur le fleuve Wouri. »
\item « Quand on est fatigué, on travaille mal. » (dit par une lycéenne)
\item « On se lève tôt au village. »
\item « On va au match samedi, toi et moi ? »
\item « On croit que les bendskins sont dangereux. »
\end{enumerate}
\end{exercice}

\begin{corrige}[Exercice « À ton tour »]
\begin{enumerate}
\item \esp{Han construido un puente sobre el río Wouri.} — agents réels non nommés → 3e pluriel.
\item \esp{Cuando una está cansada, trabaja mal.} — une femme s'inclut dans la généralité → \esp{una} + adjectif féminin (\esp{cansada}) ; \esp{trabaja} (3e sg).
\item \esp{En el pueblo se madruga.} ou \esp{En el pueblo uno se levanta temprano.} — habitude générale → \esp{se} impersonnel (verbe \esp{madrugar} = se lever tôt) ou \esp{uno}.
\item \esp{¿Vamos al partido el sábado, tú y yo?} — « on » = « nous » → 1re pluriel.
\item \esp{La gente cree que los bendskins son peligrosos.} — opinion générale → \esp{la gente} + singulier (\esp{cree}, pas \esp{*creen}) ; \esp{son} (caractéristique avec \esp{ser}).
\end{enumerate}
\end{corrige}

\begin{savaistu}[Et en Guinée équatoriale ?]
En Guinée équatoriale, seul pays hispanophone d'Afrique subsaharienne, voisin du Cameroun, l'espagnol est langue officielle. Dans la presse de Malabo comme dans celle de Madrid ou de Bogotá, on lit les mêmes tournures : \esp{Se informa que...} « On informe que... », \esp{La gente opina que...} « On est d'avis que... ». Maîtriser les cinq traductions de « on », c'est donc disposer d'un outil précieux pour lire toute la presse hispanophone, de l'autre côté de la frontière comme de l'autre côté de l'Atlantique.
\end{savaistu}

\coursec{Les adverbes d'intensité : muy, mucho, demasiado, bastante, tan, tanto}

Dans la section précédente, nous avons vu que le pronom français « on » possède plusieurs visages espagnols : \esp{se} impersonnel, troisième personne du pluriel, \esp{uno}, \esp{la gente}, \esp{nosotros}. Nous restons dans l'atelier du traducteur, mais nous changeons d'outil : après les pronoms, voici les \cle{adverbes d'intensité}, ces petits mots qui dosent la force d'une idée — « très », « beaucoup », « trop », « assez », « tant », « si ». Ils semblent simples ; ils sont pourtant la source de l'une des fautes les plus fréquentes et les plus sanctionnées au baccalauréat : la confusion entre \esp{muy} et \esp{mucho}, ou l'oubli de l'accord de \esp{demasiado} et \esp{tanto} devant un nom. Observons d'abord un texte authentique, ancré dans notre réalité camerounaise.

\courssub{Observation : un article de presse}

\begin{textebox}[La radio comunitaria de Nkolbisson — artículo original]
En el barrio de Nkolbisson, en Yaoundé, la radio comunitaria funciona muy bien. Los oyentes llaman mucho para participar en los programas, y los jóvenes escuchan bastantes emisiones sobre la cultura de la paz. El director, don Mbarga, trabaja mucho: llega muy temprano y se va muy tarde. «Tenemos muchos proyectos, pero poco dinero», explica. «El equipo es bastante bueno, aunque a veces hay demasiado ruido en el estudio y demasiadas interrupciones.»

Los programas educativos son muy populares. Hablan de temas muy importantes: la salud, la escuela, el fútbol. «Antes había muchos conflictos entre los jóvenes del barrio; ahora hay bastante menos tensión», dice una profesora. Sin embargo, algunos vecinos piensan que la música suena demasiado alto por la noche. «Es mucho más agradable escuchar las noticias que oír discusiones», añade un taxista.

La radio emite muchas horas al día y recibe bastantes cartas de los oyentes. Para don Mbarga, el balance es muy positivo: «La gente está mucho más unida. Eso vale mucho más que el dinero.»
\end{textebox}

\textbf{Glossaire :} \esp{el oyente} : l'auditeur ; \esp{la emisión} : l'émission ; \esp{temprano} : tôt ; \esp{aunque} : bien que ; \esp{el ruido} : le bruit ; \esp{los vecinos} : les voisins ; \esp{alto} (ici) : fort ; \esp{el balance} : le bilan ; \esp{unido/a} : uni(e).

\textbf{Questions de compréhension :}
\begin{enumerate}
\item ¿Dónde está la radio y quién es su director?
\item ¿Qué problemas menciona don Mbarga?
\item ¿Qué piensan algunos vecinos de la música por la noche?
\end{enumerate}

\textbf{Observation grammaticale.} Soulignons dans le texte tous les mots d'intensité et classons-les :

\begin{center}
\begin{tabularx}{\linewidth}{|X|X|X|}\hline
\rowcolor{popblueL} Exemple du texte & Mot d'intensité & Mot accompagné (nature) \\ \hline
\esp{funciona muy bien} & \esp{muy} & adverbe (\esp{bien}) \\ \hline
\esp{muy temprano / muy tarde} & \esp{muy} & adverbes \\ \hline
\esp{muy populares / muy importantes} & \esp{muy} & adjectifs \\ \hline
\esp{llaman mucho / trabaja mucho} & \esp{mucho} & verbes \\ \hline
\esp{muchos proyectos / muchas horas} & \esp{muchos/as} & noms \\ \hline
\esp{mucho más agradable} & \esp{mucho} & comparatif \\ \hline
\esp{demasiado ruido / demasiadas interrupciones} & \esp{demasiado/as} & noms \\ \hline
\esp{suena demasiado alto} & \esp{demasiado} & adjectif \\ \hline
\esp{bastante bueno / bastante menos tensión} & \esp{bastante} & adjectif / nom \\ \hline
\esp{bastantes emisiones / bastantes cartas} & \esp{bastantes} & noms \\ \hline
\end{tabularx}
\end{center}

Deux grandes lois apparaissent déjà : \esp{muy} ne change \cle{jamais} ; \esp{mucho}, \esp{demasiado} et \esp{bastante} changent de forme \cle{devant un nom}, mais restent invariables ailleurs. Précisons maintenant chaque règle.

\courssub{MUY : « très », toujours invariable}

\begin{propriete}[Muy : forme et emploi]
\esp{Muy} est un adverbe \cle{invariable}. Il se place \textbf{devant un adjectif ou devant un autre adverbe} et signifie « très ».
\begin{itemize}
\item Devant un adjectif : \esp{El cacao es muy bueno.} « Le cacao est très bon. »
\item Devant un adverbe : \esp{Habla muy rápidamente.} « Il parle très rapidement. »
\item Devant un adjectif au féminin ou au pluriel, \esp{muy} ne change pas : \esp{Las chicas están muy cansadas.} « Les filles sont très fatiguées. »
\end{itemize}
\end{propriete}

\begin{exemplebox}[Muy en contexte]
\begin{itemize}
\item \esp{Está muy cansada.} « Elle est très fatiguée. »
\item \esp{El examen fue muy difícil.} « L'examen a été très difficile. »
\item \esp{Los bendskins circulan muy deprisa.} « Les motos-taxis roulent très vite. »
\item \esp{Es una película muy interesante.} « C'est un film très intéressant. »
\end{itemize}
\end{exemplebox}

\courssub{MUCHO : « beaucoup », variable ou invariable selon le mot qui suit}

\begin{propriete}[Mucho : les trois emplois]
\textbf{1. Avec un verbe} (ou seul, comme adverbe) : \esp{mucho} est \cle{invariable}.
\esp{Trabaja mucho.} « Il travaille beaucoup. » — \esp{Me gusta mucho.} « Cela me plaît beaucoup. »

\textbf{2. Devant un nom} : \esp{mucho} fonctionne comme un adjectif quantificatif et \cle{s'accorde en genre et en nombre} avec ce nom :
\esp{mucho trabajo} « beaucoup de travail » ; \esp{mucha paciencia} « beaucoup de patience » ; \esp{muchos libros} « beaucoup de livres » ; \esp{muchas gracias} « merci beaucoup ».

\textbf{3. Devant un comparatif} : \esp{mucho} est invariable et renforce la comparaison :
\esp{Es mucho más alto que su hermano.} « Il est beaucoup plus grand que son frère. »
\end{propriete}

\begin{propriete}[Règle d'or : MUY ou MUCHO ?]
\begin{itemize}
\item \textbf{MUY} + adjectif ou adverbe : « très bon » = \esp{muy bueno} ; « très bien » = \esp{muy bien}.
\item \textbf{MUCHO} + verbe : « il travaille beaucoup » = \esp{trabaja mucho}.
\item \textbf{MUCHO/A/OS/AS} + nom : « beaucoup d'amis » = \esp{muchos amigos}.
\end{itemize}
Question-test : le mot qui suit est-il un adjectif ou un adverbe ? Alors \esp{muy}. Est-ce un verbe ou un nom ? Alors \esp{mucho}.
\end{propriete}

\begin{attention}[Erreur n°1 des francophones : *mucho bueno]
En français, « très » et « beaucoup » sont deux mots différents, mais beaucoup d'élèves traduisent « très bon » par *\esp{mucho bueno} : c'est \textbf{faux}. On dit \esp{muy bueno}. Inversement, « il parle beaucoup » ne se traduit jamais *\esp{habla muy} : on dit \esp{habla mucho}. Retenez : \esp{muy} ne peut jamais accompagner un verbe ni un nom ; \esp{mucho} ne peut jamais accompagner un adjectif (sauf comparatif : \esp{mucho más alto}).
\end{attention}

\begin{attention}[« Muy mucho » existe, « *mucho muy » non]
On entend parfois, dans la langue familière expressive, \esp{Me gusta muy mucho} « cela me plaît très beaucoup / énormément » : \esp{muy} intensifie alors l'adverbe \esp{mucho}. C'est correct mais familier ; à éviter dans une copie de Bac. En revanche, *\esp{mucho muy} est \textbf{impossible} dans tous les registres.
\end{attention}

\courssub{DEMASIADO : « trop », même double comportement que mucho}

\begin{propriete}[Demasiado : forme et emploi]
\esp{Demasiado} signifie « trop » (idée d'excès, souvent négative). Il suit exactement la logique de \esp{mucho} :
\begin{itemize}
\item \textbf{Avec un verbe} : invariable. \esp{Habla demasiado.} « Il parle trop. » \esp{Comes demasiado.} « Tu manges trop. »
\item \textbf{Devant un nom} : variable (\esp{demasiado/a/os/as}). \esp{Hay demasiado ruido.} « Il y a trop de bruit. » \esp{Demasiada gente.} « Trop de monde. » \esp{Demasiados coches.} « Trop de voitures. » \esp{Demasiadas interrupciones.} « Trop d'interruptions. »
\item \textbf{Devant un adjectif ou un adverbe} : invariable. \esp{Es demasiado caro.} « C'est trop cher. » \esp{Suena demasiado alto.} « Ça sonne trop fort. »
\end{itemize}
\end{propriete}

\begin{exemplebox}[Demasiado en contexte]
\begin{itemize}
\item \esp{El mercado está demasiado lejos.} « Le marché est trop loin. »
\item \esp{Hay demasiado ruido en la clase.} « Il y a trop de bruit dans la classe. »
\item \esp{Bebe demasiado café.} « Il boit trop de café. » (\esp{café} nom masculin singulier \rightarrow \esp{demasiado}).
\end{itemize}
\end{exemplebox}

\courssub{BASTANTE : « assez », invariable ou variable en nombre}

\begin{propriete}[Bastante : forme et emploi]
\esp{Bastante} signifie « assez » (quantité suffisante, degré notable).
\begin{itemize}
\item \textbf{Devant un adjectif ou un adverbe} : \cle{invariable}. \esp{Es bastante inteligente.} « Il est assez intelligent. » \esp{Es bastante alto.} « Il est assez grand. » \esp{Habla bastante bien.} « Il parle assez bien. »
\item \textbf{Devant un nom} : variable en \cle{nombre} seulement : \esp{bastante} (singulier) / \esp{bastantes} (pluriel), jamais de forme féminine *\esp{bastanta}. \esp{bastante dinero} « assez d'argent » ; \esp{bastante agua} « assez d'eau » ; \esp{bastantes libros} « assez de livres » ; \esp{bastantes cartas} « assez de lettres ».
\item \textbf{Avec un verbe} : invariable. \esp{Ya come bastante.} « Il mange déjà assez. »
\end{itemize}
\end{propriete}

\begin{attention}[Bastante : pas de féminin !]
Contrairement à \esp{mucho} et \esp{demasiado}, \esp{bastante} n'a \textbf{que deux formes} : \esp{bastante} / \esp{bastantes}. On écrira donc \esp{bastantes emisiones} (et non *\esp{bastantas}). De même, « assez bon » = \esp{bastante bueno} : ne dites jamais *\esp{muy bastante bueno}.
\end{attention}

\courssub{TAN, TANTO et les apocopes}

\begin{propriete}[Tan : apocope de tanto devant adjectif ou adverbe]
\esp{Tan} est la forme apocopée (raccourcie) de \esp{tanto} utilisée \textbf{uniquement devant un adjectif ou un adverbe}. Elle signifie « si », « tant » ou « aussi » (dans les comparaisons d'égalité). \esp{Tan} est \cle{toujours invariable}.
\begin{itemize}
\item \esp{Tan interesante} « si intéressant / aussi intéressant ».
\item \esp{Tan rápido} « si vite / aussi vite ».
\item \esp{Tan bien} « si bien / aussi bien ».
\end{itemize}
\end{propriete}

\begin{propriete}[Tanto : « tant (de) », « autant (de) », « tellement »]
\esp{Tanto} a deux comportements selon ce qui suit :
\begin{itemize}
\item \textbf{Devant un nom} : il fonctionne comme un adjectif quantificatif et \cle{s'accorde en genre et en nombre} : \esp{tanto tiempo}, \esp{tanta paciencia}, \esp{tanto libros}, \esp{tanta gente}, \esp{tanto dinero} (masc. sing.), \esp{tanto ruido}, \esp{tanto calor}, \esp{tanto frío}, \esp{tanto hambre}, \esp{tanto sueño}... \rightarrow \esp{tanto, tanta, tantos, tantas}.
\item \textbf{Avec un verbe} : invariable, signifie « tellement » / « autant ». \esp{Trabaja tanto.} « Il travaille tellement / autant. »
\item \textbf{Devant un adjectif ou adverbe} : on utilise \esp{tan} (voir ci-dessus), jamais *\esp{tanto}.
\end{itemize}
\end{propriete}

\begin{propriete}[Comparaisons d'égalité : tan... como / tanto... como]
\begin{itemize}
\item \textbf{Adjectif / Adverbe} : \esp{tan} + adj/adv + \esp{como}. \esp{Es tan alto como su padre.} « Il est aussi grand que son père. » \esp{Corre tan rápido como tú.} « Il court aussi vite que toi. »
\item \textbf{Nom} : \esp{tanto/a/os/as} + nom + \esp{como}. \esp{Tiene tanta paciencia como su madre.} « Il a autant de patience que sa mère. » \esp{Hay tantos alumnos como ayer.} « Il y a autant d'élèves qu'hier. »
\item \textbf{Verbe} : \esp{tanto} + verbe + \esp{como}. \esp{Estudia tanto como yo.} « Il étudie autant que moi. »
\end{itemize}
\end{propriete}

\begin{attention}[Tan vs Tanto : le test du mot suivant]
\begin{itemize}
\item Mot suivant = \textbf{adjectif ou adverbe} $\rightarrow$ \esp{TAN} (invariable). \esp{Tan bueno}, \esp{tan bien}.
\item Mot suivant = \textbf{nom} $\rightarrow$ \esp{TANTO/A/OS/AS} (accord). \esp{Tanto trabajo}, \esp{tanta gente}, \esp{tanto libros}, \esp{tanto chicas}.
\item Mot suivant = \textbf{verbe} $\rightarrow$ \esp{TANTO} (invariable). \esp{Come tanto}.
\end{itemize}
Jamais *\esp{tanto bueno}, jamais *\esp{tan libros}.
\end{attention}

\begin{definition}[Apocopes des adjectifs qualificatifs (rappel utile pour le thème)]
Certains adjectifs très fréquents perdent leur voyelle finale (\cle{apocope}) lorsqu'ils sont placés \textbf{devant un nom masculin singulier}. Cela ne change pas leur sens, mais c'est obligatoire à l'écrit comme à l'oral.
\begin{center}
\begin{tabular}{|l|l|l|}\hline
\rowcolor{popblueL} Adjectif & Forme apocopée (devant masc. sg.) & Exemple \\ \hline
\esp{bueno} & \esp{buen} & \esp{un buen libro} (mais \esp{un libro bueno}) \\ \hline
\esp{malo} & \esp{mal} & \esp{un mal momento} \\ \hline
\esp{grande} & \esp{gran} & \esp{un gran hombre} (mais \esp{un hombre grande} = physiquement grand) \\ \hline
\esp{primero} & \esp{primer} & \esp{el primer día} \\ \hline
\esp{tercero} & \esp{tercer} & \esp{el tercer piso} \\ \hline
\esp{alguno} & \esp{algún} & \esp{algún problema} \\ \hline
\esp{ninguno} & \esp{ningún} & \esp{ningún problema} \\ \hline
\esp{uno} & \esp{un} & \esp{un amigo} (article indéfini) \\ \hline
\esp{ciento} & \esp{cien} & \esp{cien años} (mais \esp{ciento uno}) \\ \hline
\end{tabular}
\end{center}
\textbf{Attention :} \esp{muy}, \esp{mucho}, \esp{demasiado}, \esp{bastante}, \esp{tan}, \esp{poco} ne s'apocopent \textbf{jamais}.
\end{definition}

\begin{exemplebox}[Tan, Tanto et apocopes en contexte]
\begin{itemize}
\item \esp{Es tan buen amigo que nunca falla.} (apocope \esp{buen} + \esp{tan}) « C'est un si bon ami qu'il ne déçoit jamais. »
\item \esp{No tengo tanto dinero como tú.} (accord \esp{tanto} + nom masc. sg.) « Je n'ai pas autant d'argent que toi. »
\item \esp{Hay tantas chicas como chicos en la clase.} (accord \esp{tantas} + nom fém. pl.) « Il y a autant de filles que de garçons dans la classe. »
\item \esp{El primer día fue muy duro.} (apocope \esp{primer}) « Le premier jour a été très dur. »
\item \esp{No hay ningún problema.} (apocope \esp{ningún}) « Il n'y a aucun problème. »
\end{itemize}
\end{exemplebox}

\courssub{L'échelle d'intensité et la comparaison avec le français}

\begin{popfigure}
\begin{center}
\begin{tikzpicture}[font=\small]
\draw[-{Stealth}, thick, popline] (0,0) -- (11.5,0);
\foreach \x/\palabra/\ejemplo/\col in {1/poco/poco dinero/popmuted, 4/bastante/bastante bueno/popgreen, 7/muy -- mucho/muy alto -- trabaja mucho/popblue, 10/demasiado/demasiado caro/poporange}{
  \fill[\col] (\x,0) circle (0.09);
  \node[above, align=center, \col] at (\x,0.15) {\textbf{\palabra}};
  \node[below, align=center, popdark] at (\x,-0.15) {\emph{\ejemplo}};
}
\node[align=center, popmuted] at (5.5,-1.3) {intensité croissante : « peu » $<$ « assez » $<$ « très / beaucoup » $<$ « trop »};
\end{tikzpicture}
\end{center}
\legende{Frise d'intensité : de \esp{poco} à \esp{demasiado}, avec un exemple à chaque cran.}
\end{popfigure}

\begin{popfigure}
\begin{center}
\begin{tikzpicture}[font=\small]
\node[draw, rounded corners, fill=popblueL, align=center, minimum width=5.6cm, minimum height=1.2cm] (adj) at (0,0) {\textbf{Devant ADJECTIF ou ADVERBE}\\ \textbf{invariable}};
\node[draw, rounded corners, fill=popgreenL, align=center, minimum width=5.6cm, minimum height=1.2cm] (nom) at (7.6,0) {\textbf{Avec un NOM}\\ \textbf{variable (accord G\&N)}};
\node[align=left, anchor=north west] at (-2.9,-1.1) {\esp{muy} : \esp{muy alto}, \esp{muy bien}\\ \esp{tan} : \esp{tan alto}, \esp{tan bien}\\ \esp{demasiado} : \esp{demasiado caro}\\ \esp{bastante} : \esp{bastante inteligente}\\ \esp{mucho} : seulement + comparatif\\ \hspace{0.6cm}\esp{mucho más alto}};
\node[align=left, anchor=north west] at (4.7,-1.1) {\esp{mucho/a/os/as} : \esp{mucha gente}\\ \esp{demasiado/a/os/as} : \esp{demasiadas cosas}\\ \esp{bastante(s)} : \esp{bastantes libros}\\ \hspace{0.6cm}(pas de féminin en \esp{-a})\\ \esp{tanto/a/os/as} : \esp{tanta paciencia}};
\end{tikzpicture}
\end{center}
\legende{Deux colonnes pour ne plus se tromper : invariable devant adjectif/adverbe, accord devant un nom.}
\end{popfigure}

\begin{center}
\begin{tabularx}{\linewidth}{|X|X|X|}\hline
\rowcolor{popblueL} Français & Español & Remarque \\ \hline
très bon / très vite & \esp{muy bueno / muy rápidamente} & \esp{muy} invariable, jamais avec un verbe \\ \hline
il travaille beaucoup & \esp{trabaja mucho} & \esp{mucho} invariable avec le verbe \\ \hline
beaucoup de travail & \esp{mucho trabajo} & accord avec le nom \\ \hline
beaucoup de grâces (merci beaucoup) & \esp{muchas gracias} & féminin pluriel \\ \hline
beaucoup plus grand & \esp{mucho más alto} & \esp{mucho} + comparatif \\ \hline
si grand / aussi grand & \esp{tan grande / tan grande como} & \esp{tan} invariable \\ \hline
tant de travail / autant de travail & \esp{tanto trabajo} & \esp{tanto} accordé \\ \hline
trop cher / il parle trop & \esp{demasiado caro / habla demasiado} & invariable \\ \hline
trop de bruit / trop de gens & \esp{demasiado ruido / demasiada gente} & accord avec le nom \\ \hline
assez grand / assez bien & \esp{bastante alto / bastante bien} & invariable \\ \hline
assez de livres & \esp{bastantes libros} & pluriel seulement, pas de féminin \\ \hline
\end{tabularx}
\end{center}

\begin{aretenir}[Comparaison des deux langues]
En français, « très » et « beaucoup » ne s'accordent \textbf{jamais} : « elles sont très contentes », « beaucoup de fleurs ». En espagnol, l'accord existe devant un nom : \esp{muchas flores}, \esp{demasiadas flores}, \esp{bastantes flores}, \esp{tantas flores}. Le traducteur francophone doit donc activer un réflexe que sa propre langue ne lui donne pas : regarder le nom qui suit et accorder en genre et en nombre (sauf \esp{bastante} : nombre seulement).
\end{aretenir}

\begin{exoresolu}[Thème : phrases à traduire]
Traduisez en espagnol : 1) Elle est très fatiguée. 2) Il y a trop de bruit. 3) Il est assez grand. 4) Nous avons beaucoup de travail. 5) Elle chante beaucoup mieux que moi. 6) Ce film est si intéressant ! 7) Il a autant d'amis que moi. 8) C'est un très bon ami.
\tcblower
\textbf{Solution détaillée.}
\begin{enumerate}
\item « fatiguée » est un \textbf{adjectif} : il faut \esp{muy}. \esp{Está muy cansada.}
\item « bruit » est un \textbf{nom} masculin singulier : \esp{demasiado} s'accorde. \esp{Hay demasiado ruido.}
\item « grand » est un \textbf{adjectif} : \esp{bastante} invariable. \esp{Es bastante alto.}
\item « travail » est un \textbf{nom} masculin singulier : \esp{Tenemos mucho trabajo.}
\item « mieux » est un \textbf{comparatif} : \esp{mucho} invariable. \esp{Ella canta mucho mejor que yo.}
\item « si intéressant » : adjectif $\rightarrow$ \esp{tan}. \esp{¡Es tan interesante!}
\item « autant d'amis » : nom masculin pluriel $\rightarrow$ \esp{tantos}. \esp{Tiene tantos amigos como yo.}
\item « très bon ami » : adjectif \esp{bueno} devant nom masc. sg. $\rightarrow$ apocope \esp{buen} + \esp{muy}. \esp{Es un muy buen amigo.}
\end{enumerate}
\end{exoresolu}

\begin{exercice}[À vous de jouer]
Traduisez : 1) Le film est très intéressant. 2) Les élèves parlent beaucoup. 3) Il y a trop de voitures à Douala. 4) Elle a assez d'argent. 5) Ce marché est beaucoup plus grand que l'autre. 6) Il fait très chaud aujourd'hui. 7) Tu es aussi gentil que ton frère. 8) Nous avons tant de devoirs ! 9) Il court si vite. 10) C'est un grand homme (figure importante).
\end{exercice}

\begin{corrige}[Exercice : adverbes d'intensité]
1) \esp{La película es muy interesante.} 2) \esp{Los alumnos hablan mucho.} 3) \esp{Hay demasiados coches en Douala.} (nom masculin pluriel) 4) \esp{Tiene bastante dinero.} 5) \esp{Este mercado es mucho más grande que el otro.} 6) \esp{Hace mucho calor hoy.} — attention : « il fait très chaud » se dit \esp{hace mucho calor}, car \esp{calor} est un \textbf{nom} en espagnol, pas un adjectif ! 7) \esp{Eres tan amable como tu hermano.} 8) \esp{Tenemos tantos deberes!} (nom masc. pl.) 9) \esp{Corre tan rápido.} (adverbe $\rightarrow$ \esp{tan}) 10) \esp{Es un gran hombre.} (apocope de \esp{grande} devant nom masc. sg. pour signifier « illustre » ; \esp{un hombre grande} = physiquement grand).
\end{corrige}

\begin{attention}[Le piège de « il fait très chaud » et apparentés]
Le français utilise un adjectif (« chaud », « froid », « faim », « soif »), l'espagnol utilise un \textbf{nom} (\esp{calor}, \esp{frío}, \esp{hambre}, \esp{sed}). On dit donc :
\begin{itemize}
\item \esp{Hace mucho calor / frío.} (Jamais *\esp{hace muy calor})
\item \esp{Tengo mucha hambre / sed / sueño / prisa.} (Jamais *\esp{tengo muy hambre})
\item \esp{Es mucho trabajo.} (C'est beaucoup de travail / C'est très dur comme travail)
\end{itemize}
Même logique : \esp{¿Tienes mucho frío?} « As-tu très froid ? »
\end{attention}

\begin{savaistu}[Muy/mucho/tan/tanto au Bac : des points faciles à gagner]
Les correcteurs du baccalauréat camerounais le confirment chaque année : la confusion \esp{muy}/\esp{mucho}, l'oubli de l'accord de \esp{tanto}/\esp{demasiado} devant nom, et l'emploi de *\esp{tanto} devant adjectif figurent parmi les fautes les plus fréquentes — et les plus sanctionnées — dans les épreuves de thème et de version. Or ce sont des règles sans exception : adjectif/adverbe $\rightarrow$ \esp{muy} / \esp{tan} (invariables) ; verbe $\rightarrow$ \esp{mucho} / \esp{tanto} / \esp{demasiado} / \esp{bastante} (invariables) ; nom $\rightarrow$ \esp{mucho/a/os/as}, \esp{tanto/a/os/as}, \esp{demasiado/a/os/as}, \esp{bastante(s)} (accord). La maîtriser parfaitement, c'est sécuriser des points faciles dans chaque copie de traduction. Entraînez-vous avec le test du « mot qui suit » jusqu'à ce qu'il devienne un réflexe automatique.
\end{savaistu}

\coursec{Tan, tanto et les apocopes}

Dans la section précédente, nous avons étudié les quatre grands adverbes d'intensité : \esp{muy} (très), \esp{mucho} (beaucoup), \esp{demasiado} (trop) et \esp{bastante} (assez). Nous avons vu que \esp{muy} modifie un adjectif ou un adverbe, tandis que \esp{mucho} accompagne un verbe ou un nom. Il reste maintenant à explorer le couple le plus délicat de cette famille : \cle{tan} et \cle{tanto}, qui traduisent « si », « tellement » et « tant ». Ces deux mots obéissent à une logique voisine de celle de \esp{muy}/\esp{mucho}, mais avec une difficulté supplémentaire : le phénomène de l'\cle{apocope}, c'est-à-dire la chute d'une syllabe devant certains mots. C'est aussi l'occasion de découvrir que \esp{bueno}, \esp{grande} et \esp{malo} subissent le même sort.

\courssub{Observation : tan ou tanto ?}

Lisons d'abord ce court dialogue entre deux lycéens de Yaoundé qui préparent l'examen d'espagnol.

\begin{textebox}[Diálogo : antes del examen]
— Estoy \textbf{tan} nervioso... ¡El examen es \textbf{tan} difícil!

— No exageres. Has estudiado \textbf{tanto} este mes que vas a triunfar.

— Es verdad, he leído \textbf{tantos} textos y he hecho \textbf{tantas} traducciones...

— Además, hablas español \textbf{tan} bien como tu hermana mayor.

— Gracias. ¡Qué día \textbf{tan} importante para nosotros!

— Sí. Y recuerda: nuestro profesor es un \textbf{gran} maestro, nos ha \textbf{preparados} muy bien.
\end{textebox}

\textbf{Traduction.} « Je suis si nerveux... L'examen est tellement difficile ! — N'exagère pas. Tu as tellement étudié ce mois-ci que tu vas réussir. — C'est vrai, j'ai lu tant de textes et fait tant de traductions... — En plus, tu parles espagnol aussi bien que ta sœur aînée. — Merci. Quelle journée si importante pour nous ! — Oui. Et souviens-toi : notre professeur est un grand maître, il nous a très bien préparés. »

Observons la répartition des formes :

\begin{center}
\begin{tabularx}{\linewidth}{|l|X|X|}
\hline
\rowcolor{popblueL} \textbf{Forme} & \textbf{Mot qui suit} & \textbf{Exemple du dialogue} \\
\hline
\esp{tan} & adjectif & \esp{tan nervioso}, \esp{tan difícil} \\
\hline
\esp{tan} & adverbe & \esp{tan bien} \\
\hline
\esp{tanto} & verbe & \esp{has estudiado tanto} \\
\hline
\esp{tanto(s)/tanta(s)} & nom & \esp{tantos textos}, \esp{tantas traducciones} \\
\hline
\end{tabularx}
\end{center}

La loi est déjà visible : \esp{tan} devant adjectif et adverbe, \esp{tanto} avec le verbe, \esp{tanto/tanta/tantos/tantas} devant un nom. C'est exactement le même partage qu'entre \esp{muy} (invariable, devant adjectif/adverbe) et \esp{mucho} (variable devant un nom).

\courssub{Tan : « si », « tellement » devant adjectif ou adverbe}

\begin{propriete}[Forme et emploi de tan]
\esp{Tan} est \textbf{invariable} et se place \textbf{devant un adjectif ou un adverbe}. Il traduit « si », « tellement », « aussi » (dans la comparaison).
\begin{itemize}
\item \esp{tan} + adjectif : \esp{El monte Camerún es tan alto.} « Le mont Cameroun est si haut. »
\item \esp{tan} + adverbe : \esp{Corre tan rápido.} « Il court si vite. »
\end{itemize}
\end{propriete}

\begin{exemplebox}[Tan en contexte]
\esp{El partido fue tan emocionante que nadie se fue antes del final.} « Le match fut si passionnant que personne ne partit avant la fin. »

\esp{Habla tan despacio que todos la entienden.} « Elle parle si lentement que tout le monde la comprend. »

\esp{No sabía que Malabo era tan bonito.} « Je ne savais pas que Malabo était si beau. »
\end{exemplebox}

\textbf{La comparaison d'égalité : tan... como.} Pour dire « aussi... que », l'espagnol emploie \esp{tan} + adjectif/adverbe + \esp{como} :

\esp{Pedro es tan alto como su hermano.} « Pedro est aussi grand que son frère. »

\esp{No corro tan rápido como tú.} « Je ne cours pas aussi vite que toi. »

\esp{Esta novela es tan interesante como la película.} « Ce roman est aussi intéressant que le film. »

\begin{attention}[Le piège du francophone]
En français, on dit « si grand que », « aussi grand que » ; en espagnol, une seule forme : \esp{tan... como}. On n'écrit jamais \esp{*muy... como} ni \esp{*tanto alto como}. Devant un adjectif, c'est toujours \esp{tan}, jamais \esp{tanto}.
\end{attention}

\courssub{Tanto : « tant », « tellement » avec verbe ou nom}

\begin{propriete}[Les trois emplois de tanto]
\begin{enumerate}
\item \textbf{Avec un verbe} : \esp{tanto} est \textbf{invariable} et signifie « tellement, tant ». \esp{Trabaja tanto.} « Il travaille tellement. »
\item \textbf{Devant un nom} : \esp{tanto} est un \textbf{adjectif indéfini variable} : \esp{tanto dinero} (tant d'argent), \esp{tanta gente} (tant de monde), \esp{tantos libros} (tant de livres), \esp{tantas personas} (tant de personnes).
\item \textbf{Comparaison d'égalité avec un nom} : \esp{tanto... como} : \esp{Tengo tantos amigos como tú.} « J'ai autant d'amis que toi. »
\end{enumerate}
\end{propriete}

\begin{exemplebox}[Tanto en contexte]
\esp{Hay tanta gente en el mercado de Douala que no se puede caminar.} « Il y a tant de monde au marché de Douala qu'on ne peut pas marcher. »

\esp{Come tanto que siempre tiene hambre otra vez.} « Il mange tellement qu'il a toujours faim à nouveau. »

\esp{No bebas tanto café.} « Ne bois pas tant de café. »

\esp{Leo tanto en español como en francés.} « Je lis autant en espagnol qu'en français. »
\end{exemplebox}

\textbf{Tanto... como entre deux éléments.} La structure \esp{tanto... como} sert aussi à coordonner deux termes : « aussi bien... que », « tant... que » :

\esp{El español se habla tanto en España como en Guinea Ecuatorial.} « L'espagnol se parle tant en Espagne qu'en Guinée équatoriale. »

\esp{Me gustan tanto el fútbol como el baloncesto.} « J'aime autant le football que le basket-ball. »

\courssub{Les apocopes : tanto → tan, bueno → buen, grande → gran, malo → mal}

L'\cle{apocope} est la chute de la dernière syllabe (ou de la voyelle finale) d'un mot devant un autre mot. Le français la connaît aussi : « beau » devient « bel » devant voyelle (\emph{un bel homme}), « grand » se prononce différemment dans \emph{un grand homme}. L'espagnol la pratique de façon systématique et écrite.

\begin{propriete}[Règle de l'apocope]
\begin{itemize}
\item \esp{tanto} perd sa dernière syllabe devant un adjectif ou un adverbe : on écrit toujours \esp{tan} (jamais \esp{*tanto alto}).
\item \esp{bueno} → \esp{buen} et \esp{malo} → \esp{mal} devant un \textbf{nom masculin singulier} : \esp{un buen día} (un bon jour), \esp{un mal estudiante} (un mauvais élève). Mais : \esp{una buena idea}, \esp{días buenos}.
\item \esp{grande} → \esp{gran} devant \textbf{tout nom singulier} (masculin ou féminin) : \esp{un gran escritor}, \esp{una gran ciudad}. Au pluriel, pas d'apocope : \esp{grandes escritores}.
\end{itemize}
\end{propriete}

\begin{popfigure}
\begin{center}
\begin{tikzpicture}[
  forma/.style={rectangle, rounded corners, draw=popblue, fill=popblueL, minimum width=2.1cm, minimum height=0.8cm, font=\bfseries},
  apo/.style={rectangle, rounded corners, draw=poporange, fill=poporangeL, minimum width=1.6cm, minimum height=0.8cm, font=\bfseries},
  cond/.style={rectangle, draw=popgreen, fill=popgreenL, rounded corners, align=center, font=\small, minimum height=0.9cm},
  ej/.style={rectangle, draw=poppurple, fill=poppurpleL, rounded corners, align=center, font=\small\itshape, minimum height=0.9cm},
  flecha/.style={-{Stealth[length=2.5mm]}, thick, popdark}
]
\node[forma] (t1) at (0,0) {tanto};
\node[apo] (t2) at (2.9,0) {tan};
\node[cond, align=center] (t3) at (6.4,0){devant adjectif\\ou adverbe};
\node[ej, align=center] (t4) at (10.2,0){tan alto\\tan rápido};
\draw[flecha] (t1) -- (t2);
\draw[flecha] (t2) -- (t3);
\draw[flecha] (t3) -- (t4);

\node[forma] (b1) at (0,-1.6) {bueno};
\node[apo] (b2) at (2.9,-1.6) {buen};
\node[cond, align=center] (b3) at (6.4,-1.6){devant nom\\masculin singulier};
\node[ej, align=center] (b4) at (10.2,-1.6){un buen amigo\\un buen día};
\draw[flecha] (b1) -- (b2);
\draw[flecha] (b2) -- (b3);
\draw[flecha] (b3) -- (b4);

\node[forma] (g1) at (0,-3.2) {grande};
\node[apo] (g2) at (2.9,-3.2) {gran};
\node[cond, align=center] (g3) at (6.4,-3.2){devant tout nom\\singulier};
\node[ej, align=center] (g4) at (10.2,-3.2){un gran escritor\\una gran nación};
\draw[flecha] (g1) -- (g2);
\draw[flecha] (g2) -- (g3);
\draw[flecha] (g3) -- (g4);

\node[forma] (m1) at (0,-4.8) {malo};
\node[apo] (m2) at (2.9,-4.8) {mal};
\node[cond, align=center] (m3) at (6.4,-4.8){devant nom\\masculin singulier};
\node[ej, align=center] (m4) at (10.2,-4.8){un mal momento\\mal tiempo};
\draw[flecha] (m1) -- (m2);
\draw[flecha] (m2) -- (m3);
\draw[flecha] (m3) -- (m4);
\end{tikzpicture}
\end{center}
\legende{Les quatre apocopes : forme pleine, forme apocopée, condition d'emploi et exemples.}
\end{popfigure}

\textbf{La nuance de sens de gran.} Placé \textbf{avant} le nom, \esp{gran} exprime la \textbf{qualité, la valeur morale, l'importance} ; placé \textbf{après}, \esp{grande} décrit la \textbf{taille physique} :

\esp{un gran hombre} « un grand homme » (un homme remarquable) ; \esp{un hombre grande} « un homme grand » (de grande taille).

\esp{una gran ciudad} « une grande ville » (importante) ; \esp{una ciudad grande} « une ville grande » (étendue).

\begin{savaistu}[Gran : la place change le sens]
Cette particularité rappelle le français, où l'adjectif change de sens selon sa place : \emph{un grand homme} (illustre) ne veut pas dire la même chose qu'\emph{un homme grand} (de haute taille). L'espagnol et le français fonctionnent ici de la même manière : l'adjectif placé avant le nom prend une valeur figurée, subjective ; placé après, il reste descriptif et objectif. Autres paires françaises comparables : \emph{un pauvre homme} (à plaindre) / \emph{un homme pauvre} (sans argent) ; en espagnol : \esp{un pobre hombre} / \esp{un hombre pobre}. Retenir ce parallélisme est un excellent réflexe de traduction.
\end{savaistu}

\textbf{Cas particulier : ¡Qué + nom + tan + adjectif !} Pour traduire « Quel jour si beau ! », l'espagnol place \esp{tan} entre le nom et l'adjectif :

\esp{¡Qué día tan bonito!} « Quelle belle journée ! » (litt. « quel jour si beau »)

\esp{¡Qué partido tan emocionante!} « Quel match passionnant ! »

\esp{¡Qué chica tan inteligente!} « Quelle fille intelligente ! »

\begin{aretenir}[Tableau récapitulatif]
\begin{center}
\begin{tabularx}{\linewidth}{|l|X|X|}
\hline
\rowcolor{popblueL} \textbf{Français} & \textbf{Español} & \textbf{Remarque} \\
\hline
si / tellement + adjectif & \esp{tan} + adjectif & invariable \\
\hline
aussi... que & \esp{tan... como} & comparaison d'égalité \\
\hline
tellement + verbe & verbe + \esp{tanto} & invariable \\
\hline
tant de + nom & \esp{tanto/a/os/as} + nom & variable \\
\hline
tant... que (coordination) & \esp{tanto... como} & entre deux éléments \\
\hline
un grand écrivain & \esp{un gran escritor} & valeur, avant le nom \\
\hline
un homme grand & \esp{un hombre grande} & taille, après le nom \\
\hline
Quel beau jour ! & \esp{¡Qué día tan bonito!} & qué + nom + tan + adj. \\
\hline
\end{tabularx}
\end{center}
\end{aretenir}

\begin{exoresolu}[Thème : traduisez en espagnol]
Énoncé. Traduisez : 1. Il court si vite que personne ne peut le suivre. 2. Il y a tant de monde au stade ! 3. Yaoundé est aussi belle que Malabo. 4. García Márquez est un grand écrivain. 5. Quel match si passionnant ! 6. Elle parle tant espagnol que français.
\tcblower
Solution détaillée.

1. « si vite » = adverbe → \esp{tan} : \esp{Corre tan rápido que nadie puede seguirlo.} (et non \esp{*tanto rápido}).

2. « tant de monde » = nom féminin singulier \esp{gente} → \esp{tanta} : \esp{¡Hay tanta gente en el estadio!}

3. Comparaison d'égalité devant adjectif → \esp{tan... como} : \esp{Yaundé es tan bonita como Malabo.}

4. « grand » au sens de valeur, devant nom masculin singulier → apocope : \esp{García Márquez es un gran escritor.} (\esp{un escritor grande} signifierait « un écrivain de grande taille ».)

5. Structure exclamative : \esp{¡Qué partido tan emocionante!}

6. Coordination de deux noms → \esp{tanto... como} : \esp{Habla tanto español como francés.}
\end{exoresolu}

\begin{attention}[Erreurs fréquentes des francophones]
\begin{itemize}
\item Écrire \esp{*tanto alto} : devant un adjectif, c'est toujours \esp{tan alto}. \esp{Tanto} ne précède jamais directement un adjectif.
\item Oublier l'accord : \esp{*tanto personas} est faux ; on écrit \esp{tantas personas}.
\item Calquer « très » : « si beau » ne se traduit pas par \esp{*muy bonito} dans la comparaison ; « aussi beau que » = \esp{tan bonito como}.
\item Placer \esp{gran} après le nom : \esp{*un hombre gran} n'existe pas ; après le nom, on retrouve la forme pleine \esp{grande} (taille).
\item Oublier l'apocope : \esp{*un bueno amigo} est une faute ; on écrit \esp{un buen amigo}.
\end{itemize}
\end{attention}

\begin{exercice}[À vous de jouer]
Traduisez en espagnol : 1. Ce film est tellement long ! 2. J'ai autant de devoirs que toi. 3. C'est une grande nation. 4. Il a eu un mauvais moment. 5. Quelle ville si animée ! 6. Elle chante aussi bien que sa sœur.
\end{exercice}

\begin{corrige}[Exercice : tan, tanto et les apocopes]
1. \esp{¡Esta película es tan larga!} 2. \esp{Tengo tantos deberes como tú.} 3. \esp{Es una gran nación.} 4. \esp{Ha pasado un mal momento.} 5. \esp{¡Qué ciudad tan animada!} 6. \esp{Canta tan bien como su hermana.}
\end{corrige}

Nous disposons désormais de tout l'arsenal des adverbes d'intensité : \esp{muy}, \esp{mucho}, \esp{demasiado}, \esp{bastante}, \esp{tan} et \esp{tanto}, avec leurs apocopes. La dernière section de la leçon vous proposera un entraînement complet au thème et à la version, où ces outils seront mobilisés avec les traductions du pronom « on ».

\coursec{Entraînement à la traduction : thème et version}

Après avoir étudié \esp{tan}, \esp{tanto} et les apocopes (\esp{tan bueno}, \esp{tanta gente}, \esp{gran}, \esp{buen}...), nous disposons de toutes les pièces du puzzle : les cinq visages espagnols du « on » français (\esp{se} impersonnel, troisième personne du pluriel, \esp{uno}, \esp{la gente}, \esp{nosotros}) et les adverbes d'intensité (\esp{muy}, \esp{mucho}, \esp{demasiado}, \esp{bastante}, \esp{tan}, \esp{tanto}). Il est temps de réunir ces deux savoir-faire dans l'épreuve reine du Bac : la \cle{traduction}. Cette dernière section vous donne une méthode en quatre étapes, un thème guidé phrase par phrase, une version corrigée et une grille d'auto-évaluation.

\courssub{La méthode de thème en quatre étapes}

Traduire du français vers l'espagnol (le \cle{thème}) n'est pas un acte d'inspiration : c'est une démarche contrôlée. L'erreur classique du candidat pressé est de traduire mot à mot, au fil de la phrase, et de découvrir trop tard qu'il a écrit \esp{uno trabaja} là où le contexte exigeait \esp{se trabaja}, ou \esp{muy gente} au lieu de \esp{mucha gente}.

\begin{methode}[Réussir le thème au Bac : les quatre étapes]
\begin{enumerate}
\item \textbf{Repérer} : dans le texte français, \textbf{encadrez au crayon tous les « on » et tous les adverbes d'intensité} (très, trop, beaucoup, assez, si, tant) AVANT de commencer à traduire. Ce balisage visuel vous empêche d'en oublier un seul.
\item \textbf{Analyser le sens} : pour chaque « on », demandez-vous : s'agit-il d'une vérité générale (\esp{se} ou \esp{uno}), d'une collectivité précise (\esp{la gente}, \esp{nosotros}), ou d'un « ils » indéterminé (3\textsuperscript{e} personne du pluriel) ? Pour chaque intensif : porte-t-il sur un adjectif ou un adverbe (\esp{muy}, \esp{tan}, \esp{demasiado}, \esp{bastante}) ou sur un verbe ou un nom (\esp{mucho}, \esp{tanto}, \esp{demasiado} + nom) ?
\item \textbf{Choisir l'équivalence} : appliquez le tableau des correspondances vu dans les sections précédentes. Rédigez la phrase espagnole complète, sans laisser aucun mot français sans équivalent.
\item \textbf{Vérifier} : relisez en contrôlant trois points : l'accord du verbe avec \esp{se} (singulier dans la tournure impersonnelle : \esp{se vive} ; pluriel dans la tournure passive avec complément pluriel : \esp{se venden casas}), l'accord de \esp{mucho/tanto/demasiado} employés comme adjectifs (\esp{mucha gente}, \esp{tantos problemas}, \esp{demasiadas cosas}), et les apocopes (\esp{tan} devant adjectif ou adverbe ; \esp{gran}, \esp{buen}, \esp{mal} devant nom masculin singulier).
\end{enumerate}
\end{methode}

\begin{popfigure}
\begin{tikzpicture}[
node distance=0.55cm,
etape/.style={rectangle, rounded corners=4pt, draw=popblue, fill=popblueL, text width=5.6cm, align=center, font=\small\bfseries, minimum height=1cm},
detail/.style={rectangle, draw=popmuted, fill=popcream, text width=5.6cm, align=center, font=\scriptsize},
fleche/.style={-{Stealth[length=3mm]}, thick, poporange}
]
\node[etape] (e1) {1. REPÉRER};
\node[detail, below=0.15cm of e1] (d1) {Encadrer tous les « on » et les adverbes d'intensité dans le texte français};
\node[etape, below=0.55cm of d1] (e2) {2. ANALYSER LE SENS};
\node[detail, below=0.15cm of e2] (d2) {« On » général, collectif ou indéterminé ? Intensité sur adjectif, verbe ou nom ?};
\node[etape, below=0.55cm of d2] (e3) {3. CHOISIR L'ÉQUIVALENCE};
\node[detail, below=0.15cm of e3] (d3) {\esp{se} / \esp{uno} / \esp{la gente} / 3\textsuperscript{e} pl. / \esp{nosotros} ; \esp{muy} / \esp{mucho} / \esp{tan} / \esp{tanto}...};
\node[etape, below=0.55cm of d3] (e4) {4. VÉRIFIER};
\node[detail, below=0.15cm of e4] (d4) {Accord du verbe avec \esp{se}, \esp{muy} vs \esp{mucho}, accords, apocopes, accents};
\draw[fleche] (e1) -- (d1);
\draw[fleche] (d1) -- (e2);
\draw[fleche] (e2) -- (d2);
\draw[fleche] (d2) -- (e3);
\draw[fleche] (e3) -- (d3);
\draw[fleche] (d3) -- (e4);
\end{tikzpicture}
\legende{Organigramme : la méthode de thème en quatre étapes}
\end{popfigure}

\begin{aretenir}[Le réflexe du jour du Bac]
Avant d'écrire la moindre phrase espagnole, le texte français doit être couvert de marques : les « on » encadrés, les intensifs soulignés. Cinq minutes de repérage valent mieux que vingt minutes de ratures.
\end{aretenir}

\courssub{Thème guidé, phrase par phrase}

Appliquons la méthode à cinq phrases qui combinent « on » et adverbes d'intensité. Pour chacune, nous justifions le choix.

\textbf{Phrase 1.} « On travaille beaucoup ici, mais on est très bien payé. »

\emph{Analyse} : vérité générale, contexte professionnel impersonnel → \esp{se} + verbe à la 3\textsuperscript{e} personne du singulier. « Beaucoup » porte sur le verbe \esp{trabajar} → \esp{mucho} (invariable). « Très » porte sur l'adverbe \esp{bien} → \esp{muy}.

\begin{exemplebox}[Thème corrigé — phrase 1]
\esp{Se trabaja mucho aquí, pero se está muy bien pagado.}

« On travaille beaucoup ici, mais on est très bien payé. »

Remarquez : \esp{se está... pagado} avec \esp{pagado} au masculin singulier (tournure impersonnelle, sans sujet précis) ; \esp{mucho} après le verbe, \esp{muy} devant l'adverbe \esp{bien}.
\end{exemplebox}

\textbf{Phrase 2.} « En été, on se lève très tôt au village. »

\emph{Analyse} : habitude collective d'un groupe dont le locuteur fait partie → \esp{nosotros} possible ; mais la vérité générale invite plutôt à une tournure généralisée. Attention : \esp{se levanta} seul serait ambigu (on lirait un réflexif : « il/elle se lève ») ; on préfère donc \esp{la gente} ou \esp{nosotros} : \esp{En verano, la gente se levanta muy temprano en el pueblo.} ou \esp{En verano, nos levantamos muy temprano en el pueblo.} « Très » + adverbe → \esp{muy temprano}.

\textbf{Phrase 3.} « On mange trop de viande dans les grandes villes. »

\emph{Analyse} : constat sociologique général → \esp{se come} ou \esp{uno come}. « Trop de » + nom → \esp{demasiado} adjectif, qui s'accorde en genre et en nombre : \esp{demasiada carne}. Traduction : \esp{Se come demasiada carne en las grandes ciudades.}

\textbf{Phrase 4.} « Quand on est jeune, on croit que tout est si facile. »

\emph{Analyse} : « on » renvoyant à n'importe qui, valeur de « chacun » → \esp{uno} est idéal : \esp{Cuando uno es joven, cree que todo es tan fácil.} « Si » devant adjectif → \esp{tan} (apocope de \esp{tanto}). Remarquez l'emploi de \esp{ser} (\esp{es joven}, \esp{es fácil}) : il s'agit de caractéristiques, non d'états passagers.

\textbf{Phrase 5.} « On a beaucoup parlé de la paix à la radio hier. »

\emph{Analyse} : « on » = les journalistes, les gens, indéterminé → 3\textsuperscript{e} personne du pluriel sans sujet exprimé : \esp{Ayer hablaron mucho de la paz en la radio.} « Beaucoup » sur le verbe → \esp{mucho} invariable, placé après le verbe. Le passé composé français « on a parlé » se rend ici par le \emph{pretérito indefinido} \esp{hablaron}, action passée datée (\esp{ayer}).

\begin{center}
\begin{tabularx}{\linewidth}{|X|X|X|}
\hline
\rowcolor{popblueL} Français & Español & Choix justifié \\
\hline
On travaille beaucoup ici & \esp{Se trabaja mucho aquí} & général → \esp{se} ; verbe + beaucoup → \esp{mucho} \\
\hline
On est très bien payé & \esp{Se está muy bien pagado} & très + adverbe → \esp{muy} \\
\hline
On se lève très tôt & \esp{La gente se levanta muy temprano} & collectif → \esp{la gente} ; \esp{muy} + adverbe \\
\hline
On mange trop de viande & \esp{Se come demasiada carne} & \esp{demasiado} adjectif accordé : \esp{demasiada} \\
\hline
On croit que tout est si facile & \esp{Uno cree que todo es tan fácil} & chacun → \esp{uno} ; si + adjectif → \esp{tan} \\
\hline
On a beaucoup parlé de la paix & \esp{Hablaron mucho de la paz} & indéterminé → 3\textsuperscript{e} pl. ; \esp{mucho} invariable \\
\hline
\end{tabularx}
\end{center}

\begin{exoresolu}[Thème guidé : une phrase complète]
Énoncé : traduisez en espagnol « Ici, on vit très calmement et on ne court jamais trop. »
\tcblower
Solution détaillée, étape par étape :
\begin{enumerate}
\item \textbf{Repérage} : deux « on » ; deux intensifs : « très » (devant l'adverbe « calmement ») et « trop » (modifiant le verbe « court »).
\item \textbf{Analyse} : « on » = vérité générale sur un lieu → \esp{se} + 3\textsuperscript{e} personne du singulier. « Très » + adverbe → \esp{muy}. « Trop » seul, sans nom → \esp{demasiado} adverbe, invariable.
\item \textbf{Traduction} : \esp{Aquí se vive muy tranquilamente y nunca se corre demasiado.}
\item \textbf{Vérification} : \esp{vive} et \esp{corre} au singulier après \esp{se} ; \esp{muy} devant l'adverbe \esp{tranquilamente} ; \esp{demasiado} invariable car il modifie un verbe, pas un nom ; \esp{nunca} placé devant le verbe, sans autre négation. Phrase correcte.
\end{enumerate}
\end{exoresolu}

\courssub{Version : un texte à traduire}

La \cle{version} (espagnol → français) exige le chemin inverse : reconnaître sous les formes espagnoles les valeurs du « on » et des intensifs, puis restituer un français naturel, sans calque.

\begin{textebox}[La vida en mi barrio — texte de version]
En mi barrio se vive muy tranquilo. La gente se conoce y se ayuda mucho. Uno puede pasear hasta muy tarde, aunque a veces hay demasiado ruido los sábados. Se dice que es el barrio más alegre de la ciudad, y no es tan difícil creerlo. Cuando uno llega por primera vez, la gente lo saluda con una sonrisa, y enseguida se siente como en casa. Aquí se respeta mucho a los ancianos y los niños juegan en la calle sin peligro.
\end{textebox}

\textbf{Glossaire} : \esp{el barrio} : le quartier ; \esp{tranquilo} : tranquille, paisible ; \esp{se conoce} : se connaît (réciproque) ; \esp{pasear} : se promener ; \esp{hasta muy tarde} : jusqu'à très tard ; \esp{aunque} : bien que ; \esp{el ruido} : le bruit ; \esp{alegre} : joyeux ; \esp{saludar} : saluer ; \esp{la sonrisa} : le sourire ; \esp{enseguida} : aussitôt ; \esp{los ancianos} : les personnes âgées ; \esp{sin peligro} : sans danger.

\begin{corrige}[Version corrigée et justifiée]
« Dans mon quartier, on vit très tranquillement. Les gens se connaissent et s'entraident beaucoup. On peut se promener jusqu'à très tard, bien qu'il y ait parfois trop de bruit le samedi. On dit que c'est le quartier le plus joyeux de la ville, et ce n'est pas si difficile de le croire. Quand on arrive pour la première fois, les gens vous saluent avec un sourire, et aussitôt on se sent comme à la maison. Ici, on respecte beaucoup les personnes âgées et les enfants jouent dans la rue sans danger. »

Justifications des choix :
\begin{itemize}
\item \esp{se vive} : \esp{se} impersonnel + singulier → « on vit ».
\item \esp{la gente se conoce y se ayuda} : ici \esp{se} est \textbf{réciproque} (les gens se connaissent \emph{mutuellement}) → « les gens se connaissent et s'entraident », et non « on ».
\item \esp{uno puede pasear} : \esp{uno} = chacun, n'importe qui → « on peut ».
\item \esp{se dice} : tournure impersonnelle figée → « on dit ».
\item \esp{se siente} : retour à l'impersonnel (le sujet reste \esp{uno}) → « on se sent ».
\item \esp{muy tranquilo}, \esp{muy tarde} : \esp{muy} + adjectif/adverbe → « très ». \esp{mucho} (deux fois) modifie un verbe → « beaucoup ». \esp{demasiado ruido} : \esp{demasiado} + nom masculin singulier → « trop de bruit ». \esp{tan difícil} : \esp{tan} + adjectif → « si difficile ».
\end{itemize}
\end{corrige}

\begin{attention}[En version, « se » + verbe n'est pas toujours « on »]
Avant de traduire tout \esp{se} par « on », vérifiez la nature du pronom :
\begin{itemize}
\item \esp{Se levanta a las seis.} = « Il/Elle se lève à six heures » (\textbf{réflexif}, sujet précis sous-entendu).
\item \esp{Se levantan temprano.} = « Ils/Elles se lèvent tôt » (\textbf{réflexif} pluriel, sujet précis).
\item \esp{Se vive bien aquí.} = « On vit bien ici » (\textbf{impersonnel}, aucun sujet précis possible).
\item \esp{Se venden coches.} = « On vend des voitures » / « Des voitures sont vendues » (\textbf{passif} avec \esp{se}, verbe accordé au pluriel avec le complément).
\end{itemize}
Test rapide : si le verbe a déjà un sujet précis (une personne, un nom), le \esp{se} est réflexif ou réciproque ; s'il n'y a aucun sujet possible, c'est l'impersonnel → « on » ; si le verbe s'accorde avec un complément nominal, c'est le passif → « on » ou le passif français.
\end{attention}

\courssub{Auto-évaluation : la grille de vérification}

Après chaque traduction, passez votre production au crible de cette grille. Chaque case cochée est un point gagné.

\begin{center}
\begin{tabularx}{\linewidth}{|X|X|}
\hline
\rowcolor{popgreenL} Point de contrôle & Question à se poser \\
\hline
Accord du verbe avec \esp{se} & Mon verbe est-il au singulier dans l'impersonnel (\esp{se vive}) ? Dans le passif, est-il accordé au complément (\esp{se venden casas}) ? \\
\hline
Choix de la traduction de « on » & Général (\esp{se}/\esp{uno}), collectif (\esp{la gente}/\esp{nosotros}) ou indéterminé (3\textsuperscript{e} pl.) ? \\
\hline
\esp{muy} vs \esp{mucho} & Devant adjectif ou adverbe : \esp{muy} ? Après un verbe ou devant un nom : \esp{mucho} ? \\
\hline
Accords des adjectifs d'intensité & \esp{mucha gente}, \esp{tantos problemas}, \esp{demasiadas cosas} : genre et nombre respectés ? \\
\hline
Apocopes & \esp{tan} (et non \esp{tanto}) devant adjectif/adverbe ? \esp{gran}, \esp{buen}, \esp{mal} devant nom masculin singulier ? \\
\hline
Accents et ponctuation & \esp{está}, \esp{también}, \esp{además} accentués ? ¿ et ¡ en place ? \\
\hline
\end{tabularx}
\end{center}

\begin{exercice}[À vous de jouer]
Traduisez en espagnol : 1) « On dit que ce marché est très animé, mais il y a trop de monde le samedi. » 2) « Quand on voyage beaucoup, on apprend tant de choses. » 3) Traduisez en français : \esp{En este pueblo se trabaja mucho y la gente está muy orgullosa de sus tradiciones.}
\end{exercice}

\begin{corrige}[Exercice « À vous de jouer »]
1) \esp{Se dice que este mercado es muy animado, pero hay demasiada gente los sábados.} (\esp{se dice} : impersonnel ; \esp{muy} + adjectif ; \esp{demasiada} accordé avec \esp{gente}, féminin singulier ; « le samedi » habituel → \esp{los sábados}.)

2) \esp{Cuando uno viaja mucho, aprende tantas cosas.} (\esp{uno} : valeur de « chacun » ; \esp{mucho} invariable après le verbe ; \esp{tantas} accordé avec \esp{cosas}, féminin pluriel.)

3) « Dans ce village, on travaille beaucoup et les gens sont très fiers de leurs traditions. » (\esp{se trabaja} : impersonnel ; \esp{la gente está muy orgullosa} : « les gens sont très fiers », \esp{estar} + adjectif accordé avec \esp{la gente}, féminin singulier.)
\end{corrige}

\courssub{Bilan de la leçon}

Cette leçon vous a donné deux clés de traduction : les cinq visages du « on » et la famille des intensifs. La carte mentale ci-dessous résume l'ensemble ; recopiez-la dans votre cahier et complétez-la avec vos propres exemples.

\begin{popfigure}
\begin{tikzpicture}[
mindmap, grow cyclic,
concept color=popblueL,
every node/.style={concept, font=\small},
level 1/.append style={level distance=3.4cm, sibling angle=90, font=\scriptsize},
level 2/.append style={level distance=2.1cm, sibling angle=45, font=\tiny}
]
\node[concept color=poporangeL, font=\small\bfseries] {Traduire « on » et l'intensité}
child[concept color=popblueL] { node {\esp{se} impersonnel} child { node {\esp{se vive bien}} } child { node {verbe au singulier} } }
child[concept color=popgreenL] { node {\esp{uno} / \esp{la gente}} child { node {\esp{uno cree que...}} } child { node {\esp{la gente dice}} } }
child[concept color=popgoldL] { node {3\textsuperscript{e} pl. / \esp{nosotros}} child { node {\esp{hablaron de ti}} } child { node {\esp{nos levantamos}} } }
child[concept color=poppinkL] { node {Intensité} child { node {\esp{muy} + adj./adv.} } child { node {\esp{mucho} + verbe/nom} } child { node {\esp{tan}/\esp{tanto}, \esp{demasiado}, \esp{bastante}} } };
\end{tikzpicture}
\legende{Carte mentale finale : les traductions de « on » et les adverbes d'intensité}
\end{popfigure}

\begin{aretenir}[L'essentiel à retenir]
\begin{itemize}
\item « On » général → \esp{se} + 3\textsuperscript{e} pers. du singulier ; « on » = chacun → \esp{uno} ; « on » = les gens → \esp{la gente} ; « on » = nous → \esp{nosotros} ; « on » = ils indéterminé → 3\textsuperscript{e} pers. du pluriel.
\item Devant un adjectif ou un adverbe : \esp{muy} (très), \esp{tan} (si), \esp{demasiado} (trop), \esp{bastante} (assez) — invariables.
\item Après un verbe : \esp{mucho}, \esp{tanto}, \esp{demasiado}, \esp{bastante} — invariables. Devant un nom : \esp{mucho}, \esp{tanto}, \esp{demasiado} s'accordent (\esp{mucha gente}, \esp{tantos libros}, \esp{demasiadas cosas}) ; \esp{bastante} ne prend que la marque du nombre (\esp{bastantes libros}).
\item En version, tout \esp{se} n'est pas « on » : cherchez d'abord le sujet réel de la phrase (réflexif, réciproque, impersonnel ou passif).
\item La méthode gagnante : repérer → analyser → choisir → vérifier.
\end{itemize}
\end{aretenir}

\newpage

\coursec{Activité d'intégration}

\coursec{Leçon E23 : Traduction — Les traductions du « on » ; les adverbes d'intensité}

\courssub{Observation : le « on » français et ses visages espagnols}

\begin{textebox}[Exemples authentiques : le « on » dans la presse et la vie quotidienne]
\esp{En Camerún, \textbf{se come} muy bien en los mercados populares.} (Vérité générale) \\
\esp{\textbf{Han inaugurado} el nuevo estadio de Japoma.} (On a inauguré / Ils ont inauguré — 3e pers. pl.) \\
\esp{\textbf{Uno} no debe tirar basura en la calle.} (Règle générale, registre soutenu) \\
\esp{\textbf{La gente} dice que Yaoundé es la ciudad de las siete colinas.} (Les gens disent / On dit) \\
\esp{\textbf{Nosotros} vamos a organizar el viaje la próxima semana.} (Nous, on va organiser...)
\end{textebox}

\begin{methode}[Comment identifier la valeur de « on » ?]
\begin{enumerate}
\item \textbf{Remplacez « on » par « les gens » ou « quelqu'un » :} si la phrase garde son sens $\rightarrow$ \esp{se} impersonnel, \esp{uno} ou \esp{la gente}.
\item \textbf{Remplacez « on » par « nous » :} si la phrase garde son sens $\rightarrow$ \esp{nosotros}.
\item \textbf{Transformez en passif :} « On entend les vendeuses » $\rightarrow$ « Les vendeuses sont entendues » $\rightarrow$ \esp{se} impersonnel (accord au pluriel) ou 3e pers. pl. active.
\item \textbf{Regardez le contexte :} discours officiel $\rightarrow$ \esp{uno} ; familiarité $\rightarrow$ \esp{la gente} / \esp{nosotros} ; énoncé général $\rightarrow$ \esp{se}.
\end{enumerate}
\end{methode}

\courssub{La construction impersonnelle avec « se »}

\begin{propriete}[Le « se » impersonnel : formation et accord]
\textbf{Structure :} \esp{Se} + verbe à la \textbf{3e personne du singulier} (par défaut) \textbf{OU} du \textbf{pluriel} (si le sujet logique est pluriel et placé après le verbe). \\
\textbf{Emploi :} Vérités générales, modes d'emploi, recettes, tournures passives réflexives, publicités, pancartes. \\
\textbf{Sujet logique :} Il n'y a \textbf{pas de sujet exprimé} avant le verbe. Le complément d'objet direct (s'il est placé après) détermine l'accord.
\begin{center}
\begin{tabularx}{\linewidth}{|X|X|X|}
\hline
\rowcolor{popblueL} \textbf{Français} & \textbf{Espagnol} & \textbf{Analyse} \\
\hline
On mange bien ici. & \esp{Se come bien aquí.} & Verbe au singulier (pas de COD après). \\
\hline
On vend des mangues. & \esp{Se venden mangos.} & COD \esp{mangos} (pluriel) après le verbe $\rightarrow$ accord pluriel. \\
\hline
On parle espagnol. & \esp{Se habla español.} & COD \esp{español} (singulier) $\rightarrow$ singulier. \\
\hline
Ici, on respecte les aînés. & \esp{Aquí se respeta a los mayores.} & COD animé \esp{a los mayores} (préposition \esp{a}) $\rightarrow$ singulier par défaut (ou pluriel \esp{se respetan} selon régions, mais singulier norme standard). \\
\hline
\end{tabularx}
\end{center}
\end{propriete}

\begin{attention}[Pièges fréquents pour francophones]
\begin{itemize}
\item \textbf{Ne pas confondre \esp{se} impersonnel et \esp{se} réflexif :} \esp{Se lava} (il se lave / on lave) vs \esp{Se lava las manos} (il se lave les mains). Le contexte décide.
\item \textbf{Accord au pluriel :} \esp{Se venden casas} (On vend des maisons) — \textbf{jamais} \esp{Se vende casas}.
\item \textbf{« Se » + infinitif :} \esp{Se puede comer} (On peut manger), \esp{Se debe estudiar} (On doit étudier). Le \esp{se} reste invariable.
\item \textbf{Faux ami :} \esp{On} ne se traduit \textbf{jamais} par \esp{uno} conjugué à la 3e pers. du singulier \emph{si un sujet suit} (ex: \esp{Uno come bien} = « Quelqu'un mange bien », pas « On mange bien » en général).
\end{itemize}
\end{attention}

\courssub{Les autres traductions de « on » : 3e personne du pluriel, uno, la gente, nosotros}

\begin{aretenir}[Tableau récapitulatif : les 5 visages de « on »]
\begin{center}
\begin{tabularx}{\linewidth}{|X|X|X|X|}
\hline
\rowcolor{popblueL} \textbf{Traduction} & \textbf{Valeur / Registre} & \textbf{Exemple espagnol} & \textbf{Français} \\
\hline
\esp{Se} + V (3e pers.) & Général, passif, neutre, écrit & \esp{En Douala se baila mucho.} & À Douala, on danse beaucoup. \\
\hline
3e pers. pluriel (ils/elles) & Action anonyme, oral, journalistique & \esp{Han cortado el agua en el barrio.} & On a coupé l'eau dans le quartier. \\
\hline
\esp{Uno} + V (3e sg) & Règle morale, conseil, soutenu & \esp{Uno debe respetar las normas.} & On doit / Il faut respecter les règles. \\
\hline
\esp{La gente} + V (3e sg) & « Les gens », concret, courant & \esp{La gente compra en el mercado Mokolo.} & Les gens / On achète au marché Mokolo. \\
\hline
\esp{Nosotros} + V (1e pl) & « Nous » (insistance, opposition) & \esp{Nosotros, los cameruneses, somos hospitalarios.} & Nous, les Camerounais, on est hospitaliers. \\
\hline
\end{tabularx}
\end{center}
\end{aretenir}

\begin{exemplebox}[Nuances contextuelles : Yaoundé, un samedi matin]
\esp{\textbf{Se ve} mucho movimiento en la avenida Kennedy.} (Constatation générale, impersonnel) \\
\esp{\textbf{Han puesto} nuevos semáforos en la rotonda.} (Info anonyme, « ils ont mis ») \\
\esp{\textbf{Uno} siente el calor nada más salir.} (Expérience universelle, « on » = n'importe qui) \\
\esp{\textbf{La gente} prefiere el taxi moto (\emph{bendskin}) para ir rápido.} (Habitude collective observée) \\
\esp{\textbf{Nosotros} vamos a tomar un \emph{foléré} fresco, ¿y ustedes?} (Nous, on va boire..., opposition à « vous »)
\end{exemplebox}

\courssub{Les adverbes d'intensité : muy, mucho, demasiado, bastante}

\begin{propriete}[Le grand écart \textbf{muy} vs \textbf{mucho} : la règle du mot modifié]
La question unique à se poser : \textbf{« Quel mot est modifié ? »}
\begin{center}
\begin{tabularx}{\linewidth}{|X|X|X|X|}
\hline
\rowcolor{popblueL} \textbf{Adverbe} & \textbf{Modifie un...} & \textbf{Forme} & \textbf{Exemple traduit} \\
\hline
\esp{Muy} & Adjectif / Adverbe & \textbf{Invariable} & \esp{muy caliente} (très chaud), \esp{muy bien} (très bien) \\
\hline
\esp{Mucho} & Verbe & \textbf{Invariable} & \esp{Trabaja mucho.} (Il travaille beaucoup.) \\
\hline
\esp{Mucho/a/os/as} & Nom ( = adjectif indéfini) & \textbf{Variable (G\&N)} & \esp{Mucha gente}, \esp{muchos libros}, \esp{mucha agua} \\
\hline
\esp{Demasiado} & Adjectif / Adverbe & \textbf{Invariable} & \esp{demasiado caro}, \esp{demasiado rápido} \\
\hline
\esp{Demasiado/a/os/as} & Nom & \textbf{Variable (G\&N)} & \esp{demasiados turistas}, \esp{demasiada prisa} \\
\hline
\esp{Bastante} & Adjectif / Adverbe / Verbe & \textbf{Invariable} & \esp{bastante grande}, \esp{bastante bien}, \esp{come bastante} \\
\hline
\esp{Bastante/a/os/as} & Nom & \textbf{Variable (G\&N)} & \esp{bastante dinero}, \esp{bastantes amigos} \\
\hline
\end{tabularx}
\end{center}
\end{propriete}

\begin{attention}[Erreurs fatales au Bac : \esp{muy} vs \esp{mucho}]
\begin{itemize}
\item \textbf{Interdit :} \esp{*Es muy calor}, \esp{*Tiene muy hambre}, \esp{*Come muy}. \\
\textbf{Correct :} \esp{Hace mucho calor}, \esp{Tiene mucha hambre}, \esp{Come mucho}. (Noms : \esp{calor, hambre} $\rightarrow$ \esp{mucho/a} ; Verbe : \esp{comer} $\rightarrow$ \esp{mucho}).
\item \textbf{Interdit :} \esp{*Mucho bueno}, \esp{*Mucho rápido}. \\
\textbf{Correct :} \esp{Muy bueno}, \esp{Muy rápido}. (Adjectifs/Adverbes $\rightarrow$ \esp{muy}).
\item \textbf{Accord :} \esp{Mucha gente} (fém. sg), \esp{Muchos problemas} (masc. pl). \textbf{Jamais} \esp{mucho gente}.
\end{itemize}
\end{attention}

\courssub{Tan, tanto et les apocopes}

\begin{propriete}[\textbf{Tan} vs \textbf{Tanto} : l'apocope obligatoire]
\begin{itemize}
\item \textbf{Tan} (apocope de \esp{tanto}) s'emploie \textbf{uniquement devant adjectif ou adverbe}. Sens : « si », « tant » (conséquence souvent introduite par \esp{que}).
\item \textbf{Tanto} (forme pleine) s'emploie \textbf{devant verbe} (invariable) ou \textbf{devant nom} (variable : \esp{tanto/a/os/as}).
\end{itemize}
\begin{center}
\begin{tabular}{|l|l|l|}
\hline
\rowcolor{popblueL} \textbf{Contexte} & \textbf{Forme} & \textbf{Exemple} \\
\hline
Devant adjectif / adverbe & \esp{Tan} & \esp{tan bueno}, \esp{tan rápido}, \esp{tan bien} \\
\hline
Devant verbe & \esp{Tanto} & \esp{Estudia tanto que se cansa.} \\
\hline
Devant nom (masc. sg) & \esp{Tanto} & \esp{tanto dinero} \\
\hline
Devant nom (fém. sg) & \esp{Tanta} & \esp{tanta gente} \\
\hline
Devant nom (masc. pl) & \esp{Tantos} & \esp{tanto libros} \\
\hline
Devant nom (fém. pl) & \esp{Tantas} & \esp{tanta casas} \\
\hline
\end{tabular}
\end{center}
\end{propriete}

\begin{savaistu}[Autres apocopes espagnoles (rappel Terminale)]
Certains adjectifs perdent leur voyelle finale devant un nom masculin singulier :
\begin{center}
\begin{tabular}{|l|l|l|}
\hline
\rowcolor{poporangeL} \textbf{Forme pleine} & \textbf{Apocope} & \textbf{Exemple} \\
\hline
\esp{Bueno} & \esp{Buen} & \esp{un buen amigo} \\
\hline
\esp{Malo} & \esp{Mal} & \esp{un mal momento} \\
\hline
\esp{Uno} & \esp{Un} & \esp{un libro} \\
\hline
\esp{Primero} & \esp{Primer} & \esp{el primer día} \\
\hline
\esp{Tercero} & \esp{Tercer} & \esp{el tercer piso} \\
\hline
\esp{Grande} & \esp{Gran} & \esp{un gran hombre} (devant n'importe quel genre, sens figuré) \\
\hline
\esp{Cualquiera} & \esp{Cualquier} & \esp{cualquier persona} \\
\hline
\esp{Ninguno} & \esp{Ningún} & \esp{ningún problema} \\
\hline
\esp{Alguno} & \esp{Algún} & \esp{algún día} \\
\hline
\end{tabular}
\end{center}
\emph{Attention :} \esp{Grande} $\rightarrow$ \esp{Gran} seulement devant singulier (masc/fém) et sens subjectif (« grand » au sens noble). Sinon \esp{grande} après le nom : \esp{un hombre grande} (physiquement grand).
\end{savaistu}

\begin{exemplebox}[Intensité et conséquence : \esp{tan... que} / \esp{tanto... que}]
\esp{El pescado a la brasa está \textbf{tan} bueno \textbf{que} los turistas repiten.} (Adj $\rightarrow$ \esp{tan}) \\
\esp{Caminamos \textbf{tanto} \textbf{que} nos dolían los pies.} (Verbe $\rightarrow$ \esp{tanto}) \\
\esp{Hay \textbf{tanta} gente en el mercado \textbf{que} no se puede pasar.} (Nom fém. sg $\rightarrow$ \esp{tanta}) \\
\esp{Hace \textbf{tan} buen tiempo \textbf{que} apetece pasear.} (\esp{buen} apocope de \esp{bueno} + \esp{tan})
\end{exemplebox}

\courssub{Entraînement à la traduction : thème et version — « Guía turístico por un día »}

\begin{experience}[Situation d'intégration : Concours de l'Office du Tourisme]
\textbf{Contexte :} L'Office du tourisme du Cameroun lance le concours \emph{« Guía turístico por un día »}. Les lycéens hispanisants doivent proposer une page de guide en espagnol pour accueillir des touristes hispanophones (Espagne, Guinée équatoriale, Amérique latine) dans une ville camerounaise.

\textbf{Texte de départ (français) — à adapter et traduire :}
\begin{textebox}[Texte source : Journaliste camerounais, Yaoundé, 2024]
Dans notre ville, on se lève tôt : dès six heures, on entend les vendeuses du marché qui installent leurs étals. On y trouve de tout, et les produits sont très frais. On mange très bien ici : le poisson braisé est si bon que les touristes en redemandent ! Mais attention : on marche beaucoup sous le soleil, et il fait parfois trop chaud à midi ; on est assez fatigué en fin de journée. Les gens sont très accueillants : on vous salue avec le sourire et on vous aide volontiers. Chez nous, on aime tant la musique qu'on danse dans la rue le week-end. On peut aussi visiter les monuments du centre-ville : on y apprend beaucoup de choses sur notre histoire. Nous, on est fiers de notre ville, et on vous attend nombreux !
\end{textebox}
\end{experience}

\begin{methode}[Consignes de l'activité de groupe (2 h)]
\begin{enumerate}
\item \textbf{Repérage (15 min) :} Soulignez les 10 occurrences de \esp{on}. Entourez les 6 adverbes d'intensité (\emph{très, trop, beaucoup, assez, si, tant}). Pour chaque \esp{on}, déterminez : Indéfini / « Nous » / Passif.
\item \textbf{Tableau de préparation (20 min) :} Remplissez le tableau ci-dessous (un par groupe).
\begin{center}
\begin{tabularx}{\linewidth}{|X|X|X|X|}
\hline
\rowcolor{popblueL} \textbf{Phrase source (extrait)} & \textbf{Valeur du « on »} & \textbf{Traduction choisie} & \textbf{Justification} \\
\hline
On se lève tôt & Indéfini (gens) & \esp{La gente se levanta} / \esp{Se levanta} & Habitude collective observable \\
\hline
On entend les vendeuses & Passif / Impersonnel & \esp{Se oyen las vendedoras} / \esp{Oyen las vendedoras} & Tournure passive réflexive (COD pluriel) \\
\hline
On y trouve de tout & Vérité générale & \esp{Se encuentra de todo} & \esp{se} impersonnel standard \\
\hline
On mange très bien & Vérité générale & \esp{Se come muy bien} & \esp{muy} + adv \esp{bien} \\
\hline
... est si bon & Intensité (adj) & \esp{tan bueno} & Apocope \esp{tan} + \esp{que} \\
\hline
On marche beaucoup & Verbe & \esp{Se camina mucho} & \esp{mucho} modifie verbe \\
\hline
Il fait trop chaud & Adj/Nom & \esp{Hace demasiado calor} & \esp{demasiado} + nom \\
\hline
On est assez fatigué & Adj & \esp{Uno llega bastante cansado} & \esp{bastante} + adj ; \esp{uno} soutenu \\
\hline
On vous salue / aide & 3e pers. pl. (gens) & \esp{La gente saluda} / \esp{Te saludan} & Sujet \esp{la gente} ou 3e pl. anonyme \\
\hline
On aime tant la musique & Verbe & \esp{Nos gusta tanto la música} & \esp{tanto} + verbe \esp{gustar} \\
\hline
On danse & « Nous » & \esp{Bailamos} / \esp{Nosotros bailamos} & 1e pl. \\
\hline
On peut visiter & Impersonnel capacité & \esp{Se puede visitar} & \esp{se} + infinitif \\
\hline
On y apprend & Indéfini / Passif & \esp{Uno aprende} / \esp{Se aprende} & \esp{uno} ou \esp{se} \\
\hline
Nous, on est fiers & « Nous » insisté & \esp{Nosotros estamos orgullosos} & \esp{nosotros} + \esp{estar} + adj \\
\hline
On vous attend & « Nous » (accueil) & \esp{Les esperamos} & 1e pl. objet \esp{les} (formel) \\
\hline
\end{tabularx}
\end{center}
\item \textbf{Rédaction du guide (35 min) :} 8--10 phrases en espagnol. \textbf{Contraintes obligatoires :} minimum 1 \esp{se} impersonnel, 1 \esp{la gente}, 1 \esp{nosotros}, 1 \esp{muy}, 1 \esp{mucho}, 1 \esp{demasiado}, 1 \esp{tan}. Adaptez à VOTRE ville (Douala, Bafoussam, Limbé, Kribi, Briqueterie, Mokolo...). Ajoutez 2 phrases libres (conseil + invitation).
\item \textbf{Relecture croisée (15 min) :} Échangez les textes. Check-list : \esp{muy} (adj/adv) vs \esp{mucho} (verbe/nom) ; accords \esp{demasiado/a/os/as}, \esp{bastante(s)} ; \esp{tan} (adj/adv) vs \esp{tanto} (verbe/nom) ; concordance \esp{nosotros}.
\item \textbf{Présentation orale \& Vote (25 min) :} Lecture expressive. Grille d'écoute : noter les 5 types de \esp{on} et les 7 adverbes entendus. Vote justifié en espagnol : \esp{Voto por el grupo 3 porque han usado bien el « se » impersonal y « tan... que ».}
\end{enumerate}
\end{methode}

\courssub{Proposition de production et corrigé commenté}

\begin{exoresolu}[Modèle de page de guide : Yaoundé, la ville aux sept collines]
Rédiger en espagnol une page de guide (8--10 phrases) respectant toutes les contraintes.
\tcblower
\textbf{Production modèle :}

\esp{En Yaoundé, la gente se levanta muy temprano: a las seis, ya se oyen las voces de las vendedoras del mercado Mokolo. En este mercado se encuentra de todo y los productos son muy frescos. Aquí se come muy bien: el pescado a la brasa es tan bueno que los turistas siempre piden más. Pero cuidado: se camina mucho bajo el sol y a veces hace demasiado calor al mediodía; uno llega bastante cansado al final del día. La gente es muy acogedora: te saluda con una sonrisa y te ayuda con gusto. A nosotros nos gusta tanto la música que bailamos en la calle el fin de semana. También se pueden visitar los monumentos del centro, donde uno aprende muchas cosas sobre nuestra historia. Nosotros estamos muy orgullosos de nuestra ciudad y les esperamos con los brazos abiertos. ¡No se pierdan esta experiencia única!}

\textbf{Traduction française :}
« À Yaoundé, les gens se lèvent très tôt : à six heures, on entend déjà les voix des vendeuses du marché Mokolo. Dans ce marché, on trouve de tout et les produits sont très frais. Ici on mange très bien : le poisson braisé est si bon que les touristes en redemandent toujours. Mais attention : on marche beaucoup sous le soleil et il fait parfois trop chaud à midi ; on arrive assez fatigué en fin de journée. Les gens sont très accueillants : ils vous saluent avec le sourire et vous aident volontiers. Nous aimons tant la musique que nous dansons dans la rue le week-end. On peut aussi visiter les monuments du centre, où l'on apprend beaucoup de choses sur notre histoire. Nous sommes très fiers de notre ville et nous vous attendons les bras ouverts. Ne manquez pas cette expérience unique ! »

\textbf{Commentaire linguistique détaillé :}
\begin{enumerate}
\item \esp{La gente se levanta} : « On » = « Les gens » (sujet collectif concret). \esp{La gente} (sg) + verbe pronominal \esp{levantarse}. \esp{Muy temprano} : \esp{muy} + adv.
\item \esp{Se oyen las voces} : \esp{se} impersonnel (passif réflexif). COD \esp{las voces} (pluriel) après verbe $\rightarrow$ accord pluriel \esp{oyen}. \emph{Pas} \esp{la gente oyen} (sujet explicite déjà utilisé).
\item \esp{Se encuentra de todo} : \esp{se} impersonnel vérité générale. \esp{De todo} (indéfini) ne déclenche pas l'accord pluriel.
\item \esp{Son muy frescos} : \esp{muy} + adj \esp{frescos} (accord \emph{productos}). \esp{Ser} (caractéristique intrinsèque).
\item \esp{Se come muy bien} : \esp{se} impersonnel + \esp{muy} + adv \esp{bien}. Invariable.
\item \esp{Es tan bueno que} : \esp{tan} (apocope de \esp{tanto}) devant adj \esp{bueno}. Conséquence \esp{que} + indicatif.
\item \esp{Se camina mucho} : \esp{se} impersonnel + \esp{mucho} modifie verbe \esp{caminar}. \textbf{Jamais} \esp{muy}.
\item \esp{Hace demasiado calor} : \esp{Hacer} impersonnel météo. \esp{Demasiado} + nom masc sg \esp{calor} $\rightarrow$ invariable (adjectif intensif). Si adj : \esp{demasiado caliente}.
\item \esp{Uno llega bastante cansado} : \esp{Uno} (indéfini soutenu) pour « on » générique. \esp{Bastante} + adj \esp{cansado}. \esp{Llegar} (résultat) vs \esp{estar}.
\item \esp{La gente es muy acogedora} : \esp{Ser} (qualité définissante). \esp{Muy} + adj \esp{acogedora} (accord \emph{la gente} fém. sg).
\item \esp{Te saluda... te ayuda} : 3e pers. sg (\esp{la gente} sujet) ou 3e pl anonyme (\esp{saludan/ayudan}). Ici cohérence avec \esp{la gente} sg.
\item \esp{A nosotros nos gusta tanto la música} : \esp{A nosotros} (insistance « nous »). \esp{Gustar} : sujet \esp{la música} (sg) $\rightarrow$ \esp{gusta}. \esp{Tanto} modifie verbe $\rightarrow$ forme pleine.
\item \esp{Bailamos} / \esp{Nosotros bailamos} : 1e pl. \esp{Nosotros} explicite pour contraster avec touristes.
\item \esp{Se pueden visitar} : \esp{se} impersonnel + modal \esp{poder} + infinitif. Accord pluriel (\esp{monumentos} sujet patient).
\item \esp{Uno aprende muchas cosas} : \esp{Uno} (indéfini soutenu) variation stylistique. \esp{Muchas} (variable fém. pl) modifie \esp{cosas}.
\item \esp{Nosotros estamos muy orgullosos} : \esp{Estar} (état/sentiment). \esp{Orgullosos} (masc. pl accord \emph{nosotros}). \esp{Muy} + adj.
\item \esp{Les esperamos} : \esp{Les} (COD formel « vous » pl) + \esp{esperamos} (1e pl). \esp{Con los brazos abiertos} (locution).
\item \esp{No se pierdan} : Impératif \emph{ustedes} (formel) pronominal \esp{perderse} (manquer/perdre l'occasion). \esp{Se} impersonnel impossible à l'impératif.
\end{enumerate}
\end{exoresolu}

\courssub{Exercices d'application individuelle (Thème / Version)}

\begin{exercice}[Thème : Phrases à traduire en espagnol]
Traduisez en espagnol en justifiant le choix du « on » et de l'adverbe.
\begin{enumerate}
\item On dit que le cacao du Cameroun est le meilleur du monde. (Vérité générale / Passif)
\item Nous, on a beaucoup marché pour arriver au marché Mokolo. (Nous / Verbe)
\item Il fait trop chaud pour jouer au football à midi. (Adj/Nom météo)
\item Les touristes sont si contents qu'ils reviennent l'an prochain. (Tan... que)
\item On ne doit pas jeter les déchets par terre. (Règle morale / Uno)
\item Dans ce quartier, on danse beaucoup le week-end. (Habitude / Se / Verbe)
\item Avez-vous assez d'argent pour le taxi ? (Nom / Bastante variable)
\item On vous attend avec impatience à Douala ! (Nous / Accueil)
\end{enumerate}
\end{exercice}

\begin{corrige}[Exercice Thème]
\begin{enumerate}
\item \esp{Se dice que el cacao de Camerún es el mejor del mundo.} (\esp{se} impersonnel passif ; \esp{el mejor} superlatif).
\item \esp{Nosotros hemos caminado mucho para llegar al mercado Mokolo.} (\esp{Nosotros} = nous ; \esp{mucho} modifie verbe \esp{caminado} ; \esp{hemos caminado} passé composé récent).
\item \esp{Hace demasiado calor para jugar al fútbol al mediodía.} (\esp{Hace} météo ; \esp{demasiado} + nom \esp{calor} invariable ; \esp{al} = \esp{a + el}).
\item \esp{Los turistas están tan contentos que vuelven el año que viene.} (\esp{tan} + adj \esp{contentos} ; \esp{están} état ; \esp{vuelven} présent habitude future).
\item \esp{Uno no debe tirar basura en el suelo.} (\esp{Uno} règle morale ; \esp{debe} obligation ; \esp{tirar} infinitif ; \esp{en el suelo}).
\item \esp{En este barrio se baila mucho el fin de semana.} (\esp{se} impersonnel habitude ; \esp{mucho} modifie verbe \esp{baila}).
\item \esp{¿Tienes bastante dinero para el taxi?} (\esp{Bastante} + nom masc sg \esp{dinero} $\rightarrow$ invariable).
\item \esp{Les esperamos con ilusión en Douala.} (\esp{Les} COD formel pl ; \esp{esperamos} 1e pl ; \esp{con ilusión} « avec impatience »).
\end{enumerate}
\end{corrige}

\begin{exercice}[Version : Traduisez en français]
\begin{enumerate}
\item \esp{En Guinea Ecuatorial se habla español y se come muy bien.}
\item \esp{La gente de Malabo es muy amable; te ayudan siempre.}
\item \esp{Uno no puede entender la cultura sin vivirla.}
\item \esp{Hemos visitado tantos monumentos que estamos demasiado cansados.}
\item \esp{¡No bebas tanto refresco, hace mucho calor!}
\end{enumerate}
\end{exercice}

\begin{corrige}[Exercice Version]
\begin{enumerate}
\item En Guinée équatoriale, on parle espagnol et on mange très bien. (\esp{se} impersonnel $\times$2 ; \esp{muy} + \esp{bien}).
\item Les gens de Malabo sont très aimables ; ils t'aident toujours. (\esp{la gente} sg $\rightarrow$ \esp{es} ; \emph{te ayudan} 3e pl anonyme ou \esp{la gente} considéré comme pl. notionnel).
\item On ne peut pas comprendre la culture sans la vivre. (\esp{Uno} indéfini soutenu ; \esp{puede} + infinitif).
\item Nous avons visité tant de monuments que nous sommes trop fatigués. (\esp{Tantos} variable masc pl \esp{monumentos} ; \esp{demasiado} invariable + adj \esp{cansados} ; \esp{estamos} état).
\item Ne bois pas tant de soda, il fait très chaud ! (\esp{Tanto} variable masc sg \esp{refresco} ; \esp{mucho} modifie nom \esp{calor} / verbe \esp{hace} impersonnel).
\end{enumerate}
\end{corrige}

\courssub{Bilan et prolongement}

\begin{aretenir}[Synthèse pour le Bac : « On » et Intensité]
\begin{itemize}
\item \textbf{« On » :} 5 choix contextuels. \textbf{Se} (général, passif) > \textbf{3e pl} (anonyme, oral) > \textbf{Uno} (soutenu, moral) > \textbf{La gente} (concret, « les gens ») > \textbf{Nosotros} (= nous).
\item \textbf{Intensité :} \textbf{Muy} (adj/adv) | \textbf{Mucho} (verbe / nom variable). \textbf{Demasiado / Bastante} : même logique (invariable adj/adv/verbe, variable nom). \textbf{Tan} (adj/adv) vs \textbf{Tanto} (verbe / nom variable).
\item \textbf{Apocopes :} \esp{Tan, buen, mal, un, primer, tercer, gran, cualquier, ningún, algún}.
\end{itemize}
\end{aretenir}

\begin{savaistu}[Cameroun \& Guinée équatoriale : un espace hispanophone commun]
La Guinée équatoriale (Malabo, Bata) est le seul pays africain où l'espagnol est langue officielle. Les liens culturels avec le Cameroun sont forts : frontière commune, communautés \emph{Fang}, \emph{Bubi}, commerce transfrontalier (cacao, bois, poisson), musique (\emph{makossa}, \emph{bikutsi} vs \emph{balélé}). Maîtriser le « on » et l'intensité permet de rédiger des guides touristiques bilingues pour la \emph{CEMAC} (Communauté Économique et Monétaire de l'Afrique Centrale) et d'accueillir les visiteurs hispanophones avec une langue précise : \esp{Les esperamos en Camerún, ¡donde la gente es muy acogedora y se come tan rico que querrán volver!}
\end{savaistu}

\newpage

\coursec{Méthodes et exercices résolus}

\coursec{Leçon E23 : Traduction – Le « on » français et les adverbes d'intensité}

\courssub{Observation : Un reportage sur les médias et la paix}

\begin{textebox}[Reportage original : « La voz de los medios en la construcción de la paz »]
En Camerún, como en muchos países, \textbf{se espera mucho} de los periodistas en tiempos de crisis. \textbf{La gente quiere} informaciones verdaderas, pero \textbf{a veces se publican noticias demasiado sensacionalistas} que dañan la convivencia. \textbf{Uno se pregunta} si los medios son \textbf{bastante responsables} al difundir rumores. \textbf{Se dice que} la paz \textbf{se construye} también con palabras rigurosas y \textbf{tan} necesarias como el pan. \textbf{No hay tanto} espacio para la improvisación: \textbf{se requiere} rigor, ética y \textbf{mucha} paciencia. \textbf{Es muy cierto} que una prensa libre y seria es el mejor antídoto contra el odio.
\end{textebox}

\textbf{Traduction et glossaire :} « Au Cameroun, comme dans de nombreux pays, on attend beaucoup des journalistes en temps de crise. Les gens veulent des informations vraies, mais parfois sont publiées des nouvelles trop sensationnalistes qui nuisent à la coexistence. On se demande si les médias sont assez responsables en diffusant des rumeurs. On dit que la paix se construit aussi avec des mots rigoureux et si nécessaires que le pain. Il n'y a pas tant d'espace pour l'improvisation : on exige rigueur, éthique et beaucoup de patience. C'est très vrai qu'une presse libre et sérieuse est le meilleur antidote contre la haine. » \\
\emph{Glossaire :} \esp{convivencia} (coexistence, vie en société), \esp{sensacionalista} (sensationnaliste), \esp{difundir} (diffuser, propager), \esp{improvisación} (improvisation), \esp{rigor} (rigueur), \esp{antídoto} (antidote).

\begin{experience}[Activité orale : Débat « Médias et paix »]
En binômes, préparez un dialogue de 2 minutes. \textbf{Rôle A :} Journaliste à la CRTV (Yaoundé), vous défendez le travail de vérification des faits. \textbf{Rôle B :} Blogueur influent sur les réseaux sociaux, vous privilégiez la rapidité et l'émotion. Utilisez au moins 3 traductions de « on » (\esp{se}, 3\textsuperscript{e} pl., \esp{uno}, \esp{la gente}, \esp{nosotros}) et 4 adverbes d'intensité (\esp{muy, mucho, demasiado, bastante, tan, tanto}). Présentez devant la classe.
\end{experience}

\begin{savaistu}[Le paysage médiatique : Cameroun, Espagne, Guinée équatoriale]
Au Cameroun, le paysage médiatique est bilingue (français/anglais) : CRTV (public), radios privées (Radio Balafon, STV), presse écrite (\emph{Mutations}, \emph{Le Jour}). En Espagne, on distingue presse nationale (\emph{El País}, \emph{El Mundo}), régionale et gratuite (\emph{20 Minutos}). En Guinée équatoriale, seul pays africain hispanophone, les médias sont majoritairement publics (\emph{TVGE}, \emph{Radio Nacional}) et le espagnol y est langue d'administration et d'éducation. Le module 5 insiste sur le rôle des médias dans la \emph{cultura de la paz} (résolution pacifique des conflits, lutte contre les \esp{noticias falsas} / \emph{fake news}).
\end{savaistu}

\courssub{Les multiples visages du « on » français en espagnol}

\begin{propriete}[Les cinq traductions principales de « on »]
\begin{center}
\begin{tabularx}{\linewidth}{|X|X|X|X|}
\hline
\rowcolor{popblueL}
\textbf{Sens de « on »} & \textbf{Structure espagnole} & \textbf{Verbe} & \textbf{Exemple} \\
\hline
Vérité générale, impersonnel & \esp{se} + verbe & 3\textsuperscript{e} pers. sg. & \esp{En España \textbf{se come} tarde.} \\
\hline
Sujet indéfini (« quelqu'un », « ils ») & 3\textsuperscript{e} pers. pl. (sans sujet) & 3\textsuperscript{e} pers. pl. & \esp{\textbf{Han robado} mi móvil.} \\
\hline
« Nous » (locuteur inclus) & \esp{nosotros} (souvent omis) & 1\textsuperscript{re} pers. pl. & \esp{\textbf{Hemos llegado} temprano.} \\
\hline
Chacun, l'individu (généralité morale) & \esp{uno} & 3\textsuperscript{e} pers. sg. & \esp{\textbf{Uno} no sabe lo que pasará.} \\
\hline
Les gens (familier, sujet collectif) & \esp{la gente} + verbe & 3\textsuperscript{e} pers. sg. & \esp{\textbf{La gente dice} que llueve.} \\
\hline
\end{tabularx}
\end{center}
\emph{Remarque :} \esp{uno} exige les possessifs \esp{su/suyo} (\esp{Uno debe hacer su deber}). \esp{La gente} est féminin singulier (\esp{la gente está cansada}).
\end{propriete}

\begin{methode}[Choisir la bonne traduction de « on » : démarche en 3 étapes]
\begin{enumerate}
\item \textbf{Remplacer mentalement « on »} par : « les gens en général » / « quelqu'un » / « nous » / « chacun » / « les gens (familier) ».
\item \textbf{Sélectionner la structure} correspondante dans le tableau ci-dessus.
\item \textbf{Vérifier l'accord} : \esp{se} + singulier (pas de sujet) ; \esp{la gente} + singulier ; 3\textsuperscript{e} pl. sans sujet exprimé ; \esp{uno} + singulier ; \esp{nosotros} + pluriel.
\end{enumerate}
\end{methode}

\begin{exoresolu}[Thème guidé : traduire « on » selon le contexte]
Traduisez en espagnol :
\begin{enumerate}
\item Au Cameroun, on parle français et anglais.
\item On m'a dit que le match commençait à seize heures.
\item Hier soir, on a regardé le journal télévisé ensemble.
\item On ne sait jamais ce que l'avenir nous réserve.
\item On dit que les médias influencent beaucoup les jeunes.
\end{enumerate}
\tcblower
\textbf{Raisonnement phrase par phrase :}
\begin{enumerate}
\item « On » = les gens en général $\rightarrow$ \esp{se} impersonnel + singulier. \esp{Hablar} reste au singulier malgré le COD pluriel. \textbf{Traduction :} \esp{En Camerún \textbf{se habla} francés e inglés.} (\esp{e} devant \esp{inglés} pour éviter l'hiatus).
\item « On m'a dit » = agent inconnu $\rightarrow$ 3\textsuperscript{e} pers. pl. Discours indirect : imparfait de l'indicatif. \textbf{Traduction :} \esp{\textbf{Me han dicho} que el partido \textbf{empezaba} a las cuatro de la tarde.}
\item « On » = nous (réel) $\rightarrow$ 1\textsuperscript{re} pers. pl. \textbf{Traduction :} \esp{Ayer por la noche \textbf{vimos} juntos el telediario (\emph{ou} noticiero).}
\item « On » = chacun $\rightarrow$ \esp{uno}. \textbf{Traduction :} \esp{\textbf{Uno} nunca sabe lo que el futuro nos reserva.}
\item « On dit que » = formule figée $\rightarrow$ \esp{se dice que} (impersonnel) ou \esp{la gente dice que}. \textbf{Traduction :} \esp{\textbf{Se dice} que los medios de comunicación \textbf{influyen mucho} en los jóvenes.}
\end{enumerate}
\textbf{Conclusion :} analyser le sens réel avant de choisir la structure espagnole.
\end{exoresolu}

\courssub{Le « se » impersonnel vs passif réfléchi : le test décisif}

\begin{propriete}[Impersonnel ou passif réfléchi ? La règle de l'accord]
\begin{itemize}
\item \textbf{\esp{Se} impersonnel} : \textbf{pas de sujet patient}. Verbe \textbf{toujours au singulier}. \esp{Se vive bien en Yaoundé.} (On vit bien à Yaoundé.)
\item \textbf{Passif réfléchi (\esp{se} passif)} : \textbf{sujet patient explicite} (souvent après le verbe). Verbe \textbf{s'accorde} avec ce sujet (pluriel si sujet pluriel). \esp{Se venden casas.} (Des maisons sont vendues / On vend des maisons.)
\end{itemize}
\textbf{Test infaillible :} Puis-je transformer la phrase en passif français avec le nom qui suit le verbe comme sujet ? \\
\esp{Se venden coches} $\rightarrow$ \emph{Des voitures sont vendues} (OUI $\rightarrow$ Passif réfléchi, pluriel). \\
\esp{Se trabaja mucho} $\rightarrow$ \emph{*Beaucoup est travaillé} (NON $\rightarrow$ Impersonnel, singulier).
\end{propriete}

\begin{methode}[Identifier et traduire les structures en \esp{se}]
\begin{enumerate}
\item Repérer \esp{se} + verbe.
\item Chercher un nom (sujet patient) qui pourrait être le sujet du verbe.
\item \textbf{Si sujet patient présent} $\rightarrow$ Passif réfléchi $\rightarrow$ Accord (sg/pl) $\rightarrow$ Traduction : passif français (\emph{est/son fait}) ou « on » actif.
\item \textbf{Si aucun sujet patient} $\rightarrow$ Impersonnel $\rightarrow$ Singulier $\rightarrow$ Traduction : « on » / « il est... » / tournure impersonnelle.
\item \textbf{Temps composés} : \esp{se ha decidido} (impers.) vs \esp{se han decidido las fechas} (passif).
\end{enumerate}
\end{methode}

\begin{exoresolu}[Version commentée : article de presse]
Traduisez en français :
\begin{quote}
\esp{En muchos países se debate sobre el papel de las redes sociales. Se dice que informan, pero también se sabe que a veces se difunden noticias falsas. Por eso se recomienda verificar siempre las fuentes.}
\end{quote}
\tcblower
\textbf{Analyse structure par structure :}
\begin{itemize}
\item \esp{se debate} : aucun sujet patient $\rightarrow$ Impersonnel $\rightarrow$ « on débat ».
\item \esp{se dice que} / \esp{se sabe que} : formules impersonnelles figées $\rightarrow$ « on dit que » / « on sait que ».
\item \esp{se difunden noticias falsas} : \esp{noticias falsas} = sujet patient pluriel $\rightarrow$ Passif réfléchi $\rightarrow$ « de fausses nouvelles sont diffusées ».
\item \esp{se recomienda verificar} : infinitif comme sujet (impersonnel) $\rightarrow$ « on recommande de vérifier ».
\end{itemize}
\textbf{Traduction rédigée :} « Dans de nombreux pays, on débat du rôle des réseaux sociaux. On dit qu'ils informent, mais on sait aussi que de fausses informations sont parfois diffusées. C'est pourquoi l'on recommande de toujours vérifier les sources. »
\end{exoresolu}

\courssub{Les adverbes d'intensité : \esp{muy, mucho, demasiado, bastante}}

\begin{propriete}[Comportement grammatical : invariables vs variables]
\begin{center}
\begin{tabularx}{\linewidth}{|X|X|X|X|}
\hline
\rowcolor{popblueL}
\textbf{Adverbe} & \textbf{Devant adjectif / adverbe} & \textbf{Devant verbe} & \textbf{Devant nom (COD/Sujet)} \\
\hline
\esp{Muy} (très) & \esp{muy} \textbf{invariable} (\esp{muy rápido}) & \textbf{Impossible} & \textbf{Impossible} \\
\hline
\esp{Mucho} (beaucoup) & \textbf{Impossible} & \esp{mucho} \textbf{invariable} (\esp{trabaja mucho}) & \esp{mucho/a/os/as} \textbf{variable} (\esp{mucha gente, muchos libros}) \\
\hline
\esp{Demasiado} (trop) & \esp{demasiado} \textbf{invariable} (\esp{demasiado caro}) & \esp{demasiado} \textbf{invariable} (\esp{habla demasiado}) & \esp{demasiado/a/os/as} \textbf{variable} (\esp{demasiadas noticias}) \\
\hline
\esp{Bastante} (assez) & \esp{bastante} \textbf{invariable} (\esp{bastante difícil}) & \esp{bastante} \textbf{invariable} (\esp{estudia bastante}) & \esp{bastante/s} \textbf{variable} (\esp{bastantes personas}) \\
\hline
\end{tabularx}
\end{center}
\textbf{Règle d'or :} Quel mot est intensifié ? Adjectif/Adverbe $\rightarrow$ forme \textbf{invariable}. Verbe $\rightarrow$ adverbe \textbf{invariable après le verbe}. Nom $\rightarrow$ forme \textbf{variable accordée} en genre/nombre.
\end{propriete}

\begin{attention}[Piège n°1 : \esp{muy} vs \esp{mucho} — Erreur éliminatoire]
\begin{itemize}
\item \textbf{FAUX :} \esp{*Esta periodista es mucho conocida.} \quad \textbf{FAUX :} \esp{*Trabaja muy.}
\item \textbf{VRAI :} \esp{Esta periodista es \textbf{muy} conocida.} (adj.) \quad \esp{Trabaja \textbf{mucho}.} (verbe)
\item \textbf{Astuce :} \esp{Muy} = \textbf{M}odifie \textbf{A}djectif/\textbf{A}dverbe (M.A.A.). \esp{Mucho} = \textbf{M}odifie \textbf{V}erbe ou \textbf{N}om (M.V.N.).
\end{itemize}
\end{attention}

\begin{exoresolu}[Thème : \esp{muy} ou \esp{mucho} ? (et \esp{demasiado}, \esp{bastante})]
Traduisez en espagnol :
\begin{enumerate}
\item Cette journaliste est très connue et elle travaille beaucoup.
\item Il y a trop de violence à la télévision ; c'est très inquiétant.
\item Les élèves sont assez fatigués parce qu'ils ont beaucoup de devoirs.
\item Ce reportage m'a beaucoup intéressé : il est très bien documenté.
\end{enumerate}
\tcblower
\textbf{Raisonnement :}
\begin{enumerate}
\item \esp{conocida} (adj.) $\rightarrow$ \esp{muy} ; \esp{trabaja} (verbe) $\rightarrow$ \esp{mucho}. \textbf{Traduction :} \esp{Esta periodista es \textbf{muy} conocida y trabaja \textbf{mucho}.}
\item \esp{violencia} (nom f. sg.) $\rightarrow$ \esp{demasiada} ; \esp{inquietante} (adj.) $\rightarrow$ \esp{muy}. \textbf{Traduction :} \esp{Hay \textbf{demasiada} violencia en la televisión; es \textbf{muy} inquietante.}
\item \esp{cansados} (adj.) $\rightarrow$ \esp{bastante} (invar.) ; \esp{deberes} (nom m. pl.) $\rightarrow$ \esp{muchos}. \textbf{Traduction :} \esp{Los alumnos están \textbf{bastante} cansados porque tienen \textbf{muchos} deberes (\emph{ou} tareas).}
\item \esp{interesado} (verbe \esp{interesar} $\leftarrow$ me) $\rightarrow$ \esp{mucho} (adv.) ; \esp{bien documentado} (adv. + part.) $\rightarrow$ \esp{muy bien}. \textbf{Traduction :} \esp{Este reportaje me ha interesado \textbf{mucho}: está \textbf{muy bien} documentado.}
\end{enumerate}
\end{exoresolu}

\courssub{\esp{Tan, tanto} et les apocopes : formes brèves obligatoires}

\begin{propriete}[Comparaison, consécution et apocopes]
\begin{enumerate}
\item \textbf{\esp{Tan}} = « si, tellement » (invariable). Devant adjectif/adverbe : \esp{tan importante}, \esp{tan rápido}. \textbf{Apocope de \esp{tanto}.}
\item \textbf{\esp{Tanto}} = « tant, tellement ».
    \begin{itemize}
    \item Adverbe (après verbe) : invariable. \esp{Trabaja tanto.}
    \item Adjectif (devant nom) : variable. \esp{tanta información, tantos problemas}.
    \end{itemize}
\item \textbf{Structures types :}
    \begin{itemize}
    \item Comparaison d'égalité : \esp{tan... como} (adj/adv), \esp{tanto... como} (nom/verbe). \esp{Es tan listo como tú.} / \esp{Tiene tantos libros como yo.}
    \item Consécution (conséquence) : \esp{tan... que} / \esp{tanto... que}. \esp{Habla tan rápido que no le entiendo.} / \esp{Hay tanto ruido que no oigo.}
    \end{itemize}
\item \textbf{Apocopes obligatoires (devant nom masc. sg. sauf \esp{gran}) :}
\begin{center}
\begin{tabular}{|l|l|l|}
\hline
\rowcolor{poporangeL}
\textbf{Forme pleine} & \textbf{Apocope} & \textbf{Exemple} \\
\hline
\esp{bueno} & \esp{buen} & \esp{un \textbf{buen} artículo} \\
\esp{malo} & \esp{mal} & \esp{un \textbf{mal} momento} \\
\esp{grande} & \esp{gran} & \esp{una \textbf{gran} noticia} (devant \textbf{tout} nom sg.) \\
\esp{primero} & \esp{primer} & \esp{el \textbf{primer} canal} \\
\esp{tercero} & \esp{tercer} & \esp{el \textbf{tercer} piso} \\
\esp{uno} & \esp{un} & \esp{\textbf{un} día} \\
\esp{alguno} & \esp{algún} & \esp{\textbf{algún} problema} \\
\esp{ninguno} & \esp{ningún} & \esp{\textbf{ningún} libro} \\
\esp{ciento} & \esp{cien} & \esp{\textbf{cien} personas} \\
\esp{tanto} & \esp{tan} & \esp{\textbf{tan} difícil} (devant adj/adv) \\
\hline
\end{tabular}
\end{center}
\emph{Note :} \esp{gran} s'emploie devant tout nom singulier (masc. ou fém.) avec valeur de « grand, remarquable » (\esp{una gran mujer}). \esp{Grande} après le nom = taille physique (\esp{una mujer grande}).
\end{enumerate}
\end{propriete}

\begin{exoresolu}[Thème et complétion : \esp{tan}, \esp{tanto}, apocopes]
A. Traduisez :
\begin{enumerate}
\item Cette nouvelle est si importante que tous les journaux en parlent.
\item Il y a tant de chaînes qu'on ne sait pas laquelle choisir.
\item Internet est aussi rapide que la télévision pour diffuser une information.
\end{enumerate}
B. Complétez avec la forme correcte : \esp{un --- (bueno) programa}, \esp{una --- (grande) audiencia}, \esp{el --- (tercero) canal}, \esp{--- (tanto) interesante}.
\tcblower
\textbf{A. Raisonnement :}
\begin{enumerate}
\item « Si importante » (adj.) $\rightarrow$ \esp{tan} ; consécution $\rightarrow$ \esp{que}. \textbf{Traduction :} \esp{Esta noticia es \textbf{tan} importante que todos los periódicos hablan de ella.}
\item « Tant de chaînes » (nom f. pl.) $\rightarrow$ \esp{tantas}. \esp{Uno} pour « on » générique. \textbf{Traduction :} \esp{Hay \textbf{tantas} cadenas que \textbf{uno} no sabe \textbf{cuál} elegir.} (\esp{cuál} accentué : interrogatif indirect).
\item « Aussi rapide que » $\rightarrow$ \esp{tan... como}. \textbf{Traduction :} \esp{Internet es \textbf{tan} rápido \textbf{como} la televisión para difundir una información.}
\end{enumerate}
\textbf{B. Solutions :} \esp{un \textbf{buen} programa} ; \esp{una \textbf{gran} audiencia} ; \esp{el \textbf{tercer} canal} ; \esp{\textbf{tan} interesante}.
\end{exoresolu}

\courssub{Synthèse et entraînement final : version et thème mixtes}

\begin{aretenir}[Check-list avant le Bac : « on » et intensité]
\begin{itemize}
\item \textbf{« On »} : 1. \esp{se} + sg (général) 2. 3\textsuperscript{e} pl (indéfini) 3. \esp{nosotros} (nous) 4. \esp{uno} (chacun) 5. \esp{la gente} + sg (familier).
\item \textbf{\esp{Se}} : Impersonnel (sg, pas de sujet) vs Passif réfléchi (accord avec sujet patient).
\item \textbf{Intensité} : \esp{muy} (adj/adv), \esp{mucho} (verbe/nom), \esp{demasiado/bastante} (même logique que \esp{mucho}).
\item \textbf{Comparaison/Consécution} : \esp{tan} (adj/adv), \esp{tanto} (nom/verbe).
\item \textbf{Apocopes} : \esp{buen, mal, gran, primer, tercer, un, algún, ningún, cien, tan} — \textbf{obligatoires} devant nom masc. sg. (\esp{gran} : tout sg.).
\end{itemize}
\end{aretenir}

\begin{methode}[Stratégie de version complète (texte mixte)]
\begin{enumerate}
\item \textbf{Lecture globale} : saisir le thème, le ton, le registre.
\item \textbf{Surlignage ciblé} : tous les \esp{se} (impers. vs passif), tous les intensifs (\esp{muy, mucho, demasiado, bastante, tan, tanto}), les apocopes.
\item \textbf{Analyse locale} : pour chaque \esp{se}, appliquer le test du sujet patient. Pour chaque intensif, identifier le mot modifié (adj/adv/verbe/nom).
\item \textbf{Traduction phrase par phrase} : respecter la syntaxe française naturelle (passif, « on », tournures impersonnelles).
\item \textbf{Relecture finale} : fluidité, accords, sens précis (\esp{demasiado} = excès négatif ; \esp{bastante} = suffisance).
\end{enumerate}
\end{methode}

\begin{exoresolu}[Version intégrée : « Médias, paix et responsabilité »]
Traduisez en français :
\begin{quote}
\esp{En tiempos de conflicto, se espera mucho de los periodistas. La gente quiere informaciones verdaderas, pero a veces se publican noticias demasiado sensacionalistas. Uno se pregunta si los medios son bastante responsables. Se dice que la paz se construye también con palabras, y eso es muy cierto.}
\end{quote}
\tcblower
\textbf{Analyse détaillée :}
\begin{itemize}
\item \esp{se espera mucho} : Impersonnel (\esp{mucho} = adv.) $\rightarrow$ « on attend beaucoup ».
\item \esp{La gente quiere} : Sujet collectif sg $\rightarrow$ « Les gens veulent ».
\item \esp{se publican noticias} : Passif réfléchi (\esp{noticias} f. pl. = sujet) $\rightarrow$ « sont publiées des nouvelles » / « on publie des nouvelles ».
\item \esp{demasiado sensacionalistas} : \esp{demasiado} invariable devant adj. $\rightarrow$ « trop sensationnalistes ».
\item \esp{Uno se pregunta} : « On se demande » (généralité individuelle).
\item \esp{bastante responsables} : \esp{bastante} invariable devant adj. $\rightarrow$ « assez responsables ».
\item \esp{Se dice que} : Impersonnel $\rightarrow$ « On dit que ».
\item \esp{la paz se construye} : \esp{la paz} = sujet patient $\rightarrow$ Passif réfléchi / pronominal naturel $\rightarrow$ « la paix se construit ».
\item \esp{muy cierto} : \esp{muy} + adj. $\rightarrow$ « très vrai ».
\end{itemize}
\textbf{Traduction rédigée :} « En temps de conflit, on attend beaucoup des journalistes. Les gens veulent des informations véridiques, mais parfois sont publiées des nouvelles trop sensationnalistes. On se demande si les médias sont assez responsables. On dit que la paix se construit aussi avec des mots, et cela est très vrai. »
\end{exoresolu}

\begin{exercice}[Entraînement personnel : Thème et Version]
\textbf{A. Thème (Français $\rightarrow$ Espagnol). Traduisez :}
\begin{enumerate}
\item Dans ce pays, on mange très tard et on travaille beaucoup.
\item On a volé mon portable hier ; quelqu'un l'a trouvé ce matin.
\item Nous, on est allés voir un bon film au cinéma.
\item Chacun pense que sa opinion est la meilleure.
\item Il y a trop de bruit ; on n'entend pas assez bien la radio.
\item Ce journaliste est si courageux qu'il enquête sur tout.
\item Il y a tant d'informations qu'on ne sait pas quoi croire.
\item C'est un très bon reportage ; il est assez long mais très intéressant.
\end{enumerate}

\textbf{B. Version (Espagnol $\rightarrow$ Français). Traduisez :}
\begin{quote}
\esp{En la redacción, se trabaja con mucha presión. Los redactores tienen que escribir muy rápido y a veces cometen demasiados errores. La gente no entiende que no es tan fácil ser preciso. Un buen periodista sabe que la verdad es lo primero. Por eso, uno debe verificar tanto las fuentes como los datos. Se publican pocas rectificaciones, pero son muy necesarias.}
\end{quote}
\end{exercice}

\begin{corrige}[Exercice d'entraînement]
\textbf{A. Thème :}
\begin{enumerate}
\item \esp{En este país \textbf{se come} muy tarde y \textbf{se trabaja} mucho.} (Impersonnel général).
\item \esp{\textbf{Me robaron} el móvil ayer; \textbf{alguien lo encontró} esta mañana.} (3\textsuperscript{e} pl. indéfini + \esp{alguien}). \emph{Variante :} \esp{Me han robado...}
\item \esp{\textbf{Fuimos} a ver \textbf{un buen} cine.} (1\textsuperscript{re} pl. + apocope \esp{buen}).
\item \esp{\textbf{Uno} piensa que \textbf{su} opinión es la mejor.} (\esp{Uno} + possessif \esp{su}).
\item \esp{Hay \textbf{demasiado} ruido; no \textbf{se oye} \textbf{bastante bien} la radio.} (\esp{demasiado} inv. + nom ; \esp{se oye} impers. ; \esp{bastante bien} inv. + adv.).
\item \esp{Este periodista es \textbf{tan} valiente que investiga \textbf{sobre} todo.} (\esp{tan} + adj ; \esp{investigar sobre}).
\item \esp{Hay \textbf{tanta} información que \textbf{uno} no sabe qué creer.} (\esp{tanta} f. sg. ; \esp{uno} générique).
\item \esp{Es \textbf{un buen} reportaje; es \textbf{bastante} largo pero \textbf{muy} interesante.} (Apocopes \esp{buen}, \esp{bastante} inv., \esp{muy}).
\end{enumerate}

\textbf{B. Version :}
« Dans la rédaction, on travaille sous forte pression. Les rédacteurs doivent écrire très vite et commettent parfois trop d'erreurs. Les gens ne comprennent pas qu'il n'est pas si facile d'être précis. Un bon journaliste sait que la vérité est la première chose. C'est pourquoi on doit vérifier tant les sources que les données. Peu de rectifications sont publiées, mais elles sont très nécessaires. »

\textbf{Points clés du corrigé :}
\begin{itemize}
\item \esp{se trabaja} (impers. sg) ; \esp{mucha presión} (nom f. sg $\rightarrow$ \esp{mucha}).
\item \esp{demasiados errores} (nom m. pl $\rightarrow$ \esp{demasiados}).
\item \esp{tan fácil} (\esp{tan} + adj).
\item \esp{Un buen periodista} (apocope \esp{buen}).
\item \esp{lo primero} (neutre).
\item \esp{tanto las fuentes como los datos} (corrélation \esp{tanto... como} + noms $\rightarrow$ \esp{tanto} variable).
\item \esp{Se publican pocas rectificaciones} (Passif réfléchi : \esp{rectificaciones} f. pl. sujet $\rightarrow$ \esp{publican} pl.).
\item \esp{son muy necesarias} (\esp{ser} + adj. f. pl.).
\end{itemize}
\end{corrige}

\begin{attention}[Les 5 erreurs fatales au Bac : Récapitulatif ultime]
\begin{enumerate}
\item \textbf{\esp{Muy} vs \esp{Mucho}} : \esp{*muy trabajo} / \esp{*mucho cansado}. \textbf{Règle :} \esp{Muy} + Adj/Adv ; \esp{Mucho} + Verbe / Nom (accordé).
\item \textbf{\esp{Uno} universel} : Traduire « on » par \esp{uno} dans « On est allés au cinéma » (\esp{*Uno fue...}). \textbf{Correct :} \esp{Fuimos} (nosotros).
\item \textbf{Accord du \esp{se} impersonnel} : \esp{*Se viven bien} (pas de sujet). \textbf{Correct :} \esp{Se vive bien}. Mais \esp{Se venden casas} (passif, sujet \esp{casas}).
\item \textbf{Apocopes oubliées} : \esp{*un bueno día}, \esp{*el tercero canal}, \esp{*tanto difícil}. \textbf{Correct :} \esp{un buen día}, \esp{el tercer canal}, \esp{tan difícil}.
\item \textbf{Calques français} : \esp{*muy de cosas} (pour « beaucoup de choses » $\rightarrow$ \esp{muchas cosas}) ; \esp{*muy mejor} (pour « beaucoup mieux » $\rightarrow$ \esp{mucho mejor}) ; accents oubliés : \esp{television}, \esp{informacion}, \esp{nacion} $\rightarrow$ \esp{televisión}, \esp{información}, \esp{nación}.
\end{enumerate}
\end{attention}

\newpage

\coursec{Exercices}

\courssub{Je vérifie mes connaissances}

\begin{exercice}[QCM : la bonne traduction de « on »]
Pour chaque phrase française, choisis la traduction espagnole correcte parmi les quatre propositions. Une seule réponse est juste.
\begin{enumerate}
\item \textbf{On} ne peut pas fumer ici.
\begin{enumerate}
\item \esp{Nosotros no podemos fumar aquí.}
\item \esp{No se puede fumar aquí.}
\item \esp{Ellos no pueden fumar aquí.}
\item \esp{La gente no pueden fumar aquí.}
\end{enumerate}
\item \textbf{On} m'a volé mon téléphone au marché.
\begin{enumerate}
\item \esp{Uno me ha robado el teléfono en el mercado.}
\item \esp{Se me ha robado el teléfono en el mercado.}
\item \esp{Me han robado el teléfono en el mercado.}
\item \esp{La gente me ha robado el teléfono en el mercado.}
\end{enumerate}
\item Hier soir, \textbf{on} est allés au stade avec mes amis.
\begin{enumerate}
\item \esp{Anoche se fue al estadio con mis amigos.}
\item \esp{Anoche uno fue al estadio con mis amigos.}
\item \esp{Anoche la gente fue al estadio con mis amigos.}
\item \esp{Anoche fuimos al estadio con mis amigos.}
\end{enumerate}
\item \textbf{On} dit que le président visitera Douala demain.
\begin{enumerate}
\item \esp{Nosotros decimos que el presidente visitará Duala mañana.}
\item \esp{Uno dice que el presidente visitará Duala mañana.}
\item \esp{Se dice que el presidente visitará Duala mañana.}
\item \esp{Se dicen que el presidente visitará Duala mañana.}
\end{enumerate}
\end{enumerate}
\end{exercice}

\begin{exercice}[Vrai ou faux ? Justifie ta réponse]
Dis si les affirmations suivantes sont vraies ou fausses, puis justifie en une phrase.
\begin{enumerate}
\item \esp{Muy} s'emploie devant un adjectif ou un adverbe, \esp{mucho} avec un verbe ou un nom.
\item Dans la construction impersonnelle, le verbe qui suit \esp{se} est toujours au singulier, même si le nom qui suit est au pluriel.
\item \esp{La gente} est un nom féminin singulier : le verbe et l'adjectif s'accordent donc au singulier.
\item \esp{Tan} est l'apocope de \esp{tanto} devant un adjectif ou un adverbe.
\item On peut écrire \esp{*mucho bueno} pour dire « très bon ».
\end{enumerate}
\end{exercice}

\begin{exercice}[Vocabulaire : les mots de l'intensité]
Complète le tableau suivant en donnant, pour chaque adverbe français, son équivalent espagnol et un exemple espagnol de ton invention.
\begin{center}
\begin{tabularx}{\linewidth}{|X|X|X|}
\hline
\rowcolor{popblueL} Français & Español & Ejemplo \\
\hline
très (devant un adjectif) & & \\
\hline
beaucoup (avec un verbe) & & \\
\hline
trop (devant un adjectif) & & \\
\hline
assez, plutôt & & \\
\hline
si, tellement (devant un adjectif) & & \\
\hline
tant, autant (avec un nom) & & \\
\hline
\end{tabularx}
\end{center}
\end{exercice}

\begin{exercice}[Conjugaison : l'impersonnel à tous les temps]
Conjugue le verbe entre parenthèses à la forme impersonnelle (\esp{se} + 3\up{e} personne du singulier) au temps demandé.
\begin{enumerate}
\item Présent : \esp{Aquí (vender) \ldots\ldots\ldots artesanía africana.}
\item Passé composé : \esp{(Se + hablar) \ldots\ldots\ldots mucho de la paz en la reunión.}
\item Imparfait : \esp{Antes (vivir) se \ldots\ldots\ldots mejor en este barrio.}
\item Futur : \esp{Mañana (discutir) se \ldots\ldots\ldots el problema en la radio.}
\item Présent : \esp{En los exámenes no (poder) se \ldots\ldots\ldots copiar.}
\end{enumerate}
\end{exercice}

\courssub{Je m'entraîne}

\begin{exercice}[Texte à trous : muy, mucho, demasiado, bastante, tan, tanto]
Complète le texte suivant avec \esp{muy}, \esp{mucho}, \esp{mucha}, \esp{muchos}, \esp{muchas}, \esp{demasiado}, \esp{demasiada}, \esp{demasiados}, \esp{bastante}, \esp{tan} ou \esp{tanto}, selon le sens et la grammaire.
\end{exercice}

\begin{textebox}[La radio de mi barrio]
En mi barrio de Yaoundé hay una emisora de radio \ldots\ldots\ldots popular. Habla \ldots\ldots\ldots de la cultura de la paz, y sus programas son \ldots\ldots\ldots interesantes que toda la familia los escucha. Ayer hubo \ldots\ldots\ldots ruido en el estudio porque vinieron \ldots\ldots\ldots jóvenes del liceo. El presentador estaba \ldots\ldots\ldots contento: dijo que los jóvenes tienen \ldots\ldots\ldots energía y \ldots\ldots\ldots ideas nuevas. Pero algunos oyentes se quejaron: para ellos, el programa duró \ldots\ldots\ldots tiempo y la música era \ldots\ldots\ldots alta. Al final, el presentador prometió trabajar \ldots\ldots\ldots para mejorar.
\end{textebox}

\begin{exercice}[Appariements : chaque « on » a son visage]
Associe chaque phrase française (1 à 5) à la traduction espagnole qui convient (A à E), puis indique la valeur du « on » (impersonnel, nous, les gens, quelqu'un d'indéfini).
\begin{enumerate}
\item En Espagne, on dîne très tard.
\item On nous a invités au mariage de Carla.
\item Les amis et moi, on part demain pour Kribi.
\item Quand on est jeune, on croit tout savoir.
\item Ici, on parle français et anglais.
\end{enumerate}
\begin{itemize}
\item[A.] \esp{Cuando uno es joven, cree saberlo todo.}
\item[B.] \esp{Mis amigos y yo partimos mañana para Kribi.}
\item[C.] \esp{En España la gente cena muy tarde.}
\item[D.] \esp{Nos han invitado a la boda de Carla.}
\item[E.] \esp{Aquí se habla francés e inglés.}
\end{itemize}
\end{exercice}

\begin{exercice}[Reformulation : trois façons de dire]
La phrase « On a construit un pont sur le fleuve Wouri » peut se traduire de trois manières différentes en espagnol. Donne ces trois traductions (avec \esp{se} impersonnel, avec la 3\up{e} personne du pluriel, avec \esp{nosotros}), puis explique en français la nuance de sens de chacune : qui construit le pont dans chaque cas ?
\end{exercice}

\begin{exercice}[Corrige les erreurs d'élèves francophones]
Les phrases suivantes contiennent chacune une erreur typique des francophones. Réécris-les correctement et explique l'erreur en une phrase.
\begin{enumerate}
\item \esp{*Este libro es mucho bueno.}
\item \esp{*Se vende frutas en el mercado.} (sens : « on vend des fruits »)
\item \esp{*Mi hermano es tanto alto como yo.}
\item \esp{*La gente son muy amables aquí.}
\item \esp{*Me gusta muy esta canción.}
\item \esp{*Hay demasiado coches en Douala.}
\end{enumerate}
\end{exercice}

\courssub{Je me prépare au Bac}

\begin{exercice}[Texte et questions de compréhension et de langue]
Lis le texte suivant, puis réponds aux questions.
\end{exercice}

\begin{textebox}[Los medios y la paz en la escuela]
En muchos liceos de Camerún se habla cada vez más de la cultura de la paz. La gente piensa que los medios de comunicación tienen un papel muy importante: cuando uno escucha la radio o lee el periódico, aprende a respetar las opiniones de los demás. Sin embargo, a veces se difunden noticias falsas en las redes sociales, y eso crea demasiados conflictos entre los jóvenes. Por eso, en algunos establecimientos se han creado clubes de prensa escolar. Allí, los alumnos aprenden a verificar la información antes de compartirla. « Antes, uno creía todo lo que leía en Internet », explica un profesor de Yaoundé. « Ahora, los chicos son bastante más críticos ». Según él, si se educa bien a la juventud, se construye una sociedad más pacífica. La tarea es larga, pero nadie dice que sea imposible.
\end{textebox}

\textbf{I. Comprensión lectora} (réponds en espagnol, avec des phrases complètes)
\begin{enumerate}
\item ¿Qué papel tienen los medios de comunicación según la gente?
\item ¿Qué problema crean las noticias falsas?
\item ¿Para qué se han creado los clubes de prensa escolar?
\item ¿Qué cambio observa el profesor en los alumnos?
\end{enumerate}

\textbf{II. Questions de langue}
\begin{enumerate}
\item Releve dans le texte trois constructions différentes qui traduisent le « on » français et précise la valeur de chacune.
\item Justifie l'accord du verbe dans : \esp{la gente piensa}.
\item Pourquoi écrit-on \esp{demasiados conflictos} et non \esp{*demasiado conflictos} ?
\item Mets à la forme impersonnelle : \esp{Los alumnos verifican la información.} (commence par \esp{Se\ldots})
\end{enumerate}

\begin{exercice}[Thème type Bac]
Traduis en espagnol les cinq phrases suivantes. Attention : chaque « on » a une valeur différente, et les adverbes d'intensité doivent être employés correctement.
\begin{enumerate}
\item On ne travaille pas le dimanche dans ce bureau.
\item Mes cousins et moi, on aime beaucoup le football camerounais.
\item On m'a dit que ce film était très intéressant.
\item Quand on voyage trop, on est souvent fatigué.
\item Les gens boivent assez d'eau quand il fait très chaud.
\end{enumerate}
\end{exercice}

\begin{exercice}[Version type Bac]
Traduis en français le paragraphe suivant, extrait d'un article de presse scolaire. Soigne le choix des équivalents français de \esp{se}, \esp{uno} et \esp{la gente}.
\end{exercice}

\begin{textebox}[La fiesta del barrio]
En nuestro barrio se organiza cada año una gran fiesta para celebrar la paz. La gente llega muy temprano: se venden platos típicos, se baila y se canta hasta muy tarde. Cuando uno participa en esta fiesta, comprende que la convivencia es bastante más fácil de lo que parece. Los ancianos cuentan historias del pasado y los niños escuchan con demasiada atención para no olvidar nada. Al final de la noche, todos se dicen: « el año que viene volveremos ».
\end{textebox}

\begin{exercice}[Expression écrite : sujet de rédaction]
Sujet : \esp{« Los medios de comunicación y la cultura de la paz en mi país »}.

Rédige en espagnol un texte de \textbf{150 à 180 mots} dans lequel tu présentes le rôle des médias (radio, télévision, réseaux sociaux) dans la promotion de la paix au Cameroun. Tu devras employer :
\begin{itemize}
\item au moins trois traductions différentes du « on » (\esp{se} impersonnel, \esp{uno}, \esp{la gente} ou \esp{nosotros}) ;
\item au moins quatre adverbes d'intensité différents (\esp{muy}, \esp{mucho}, \esp{demasiado}, \esp{bastante}, \esp{tan}, \esp{tanto}) ;
\item un connecteur de cause (\esp{porque}, \esp{por eso}) et un connecteur d'opposition (\esp{sin embargo}, \esp{pero}).
\end{itemize}
\end{exercice}

\newpage

\coursec{Corrigés des exercices}

\begin{corrige}[Exercice 1 — QCM : la bonne traduction de « on »]
\begin{enumerate}
\item Réponse \textbf{b} : \esp{No se puede fumar aquí.} Il s'agit d'une interdiction générale, impersonnelle : la construction \esp{se} + 3\up{e} personne du singulier est la traduction naturelle. La réponse d est fausse car \esp{la gente} exige le singulier (\esp{no puede}).
\item Réponse \textbf{c} : \esp{Me han robado el teléfono en el mercado.} Le « on » désigne ici des personnes indéfinies (des voleurs) : l'espagnol utilise la 3\up{e} personne du pluriel sans sujet exprimé. La réponse b est incorrecte : \esp{se me ha robado} signifierait « mon téléphone m'a été volé (par accident) », construction de l'accident involontaire, pas du vol.
\item Réponse \textbf{d} : \esp{Anoche fuimos al estadio con mis amigos.} Le « on » français vaut ici « nous » (le locuteur est inclus : « avec mes amis ») ; l'espagnol dit simplement \esp{nosotros} ou le verbe à la 1\up{re} personne du pluriel.
\item Réponse \textbf{c} : \esp{Se dice que el presidente visitará Duala mañana.} Rumeur, rumeur publique : \esp{se dice que} + indicatif, verbe au singulier. La réponse d est fausse : devant une subordonnée introduite par \esp{que}, le verbe reste au singulier.
\end{enumerate}
\end{corrige}

\begin{corrige}[Exercice 2 — Vrai ou faux ?]
\begin{enumerate}
\item \textbf{Vrai.} \esp{Muy} modifie un adjectif ou un adverbe (\esp{muy alto}, \esp{muy bien}) ; \esp{mucho} modifie un verbe (\esp{trabaja mucho}) ou, accordé, un nom (\esp{mucha gente}).
\item \textbf{Faux.} Dans l'impersonnel avec un nom sujet, le verbe s'accorde avec ce nom : \esp{se vende fruta} mais \esp{se venden coches}. Le verbe n'est au singulier que lorsqu'il n'y a pas de nom pluriel qui joue le rôle de sujet (\esp{se vive bien}).
\item \textbf{Vrai.} \esp{La gente} est un nom collectif féminin singulier : \esp{la gente es amable}, jamais \esp{*la gente son amables}.
\item \textbf{Vrai.} Devant un adjectif ou un adverbe, \esp{tanto} s'apocope en \esp{tan} : \esp{tan alto}, \esp{tan rápidamente}.
\item \textbf{Faux.} Devant un adjectif, on emploie \esp{muy} : \esp{muy bueno}. \esp{Mucho} ne modifie jamais un adjectif.
\end{enumerate}
\end{corrige}

\begin{corrige}[Exercice 3 — Vocabulaire : les mots de l'intensité]
Tableau complété (les exemples sont des propositions ; toute phrase correcte est acceptée) :
\begin{center}
\begin{tabularx}{\linewidth}{|X|X|X|}
\hline
\rowcolor{popblueL} Français & Español & Ejemplo \\
\hline
très (devant un adjectif) & \esp{muy} & \esp{El examen fue muy difícil.} \\
\hline
beaucoup (avec un verbe) & \esp{mucho} & \esp{Mi hermana estudia mucho.} \\
\hline
trop (devant un adjectif) & \esp{demasiado} & \esp{Este café está demasiado caliente.} \\
\hline
assez, plutôt & \esp{bastante} & \esp{La película es bastante buena.} \\
\hline
si, tellement (devant un adjectif) & \esp{tan} & \esp{El partido fue tan emocionante.} \\
\hline
tant, autant (avec un nom) & \esp{tanto / tanta / tantos / tantas} & \esp{Hay tantos alumnos en la clase.} \\
\hline
\end{tabularx}
\end{center}
Rappel : \esp{demasiado} et \esp{tanto} s'accordent devant un nom (\esp{demasiada agua}, \esp{tantas ideas}) ; \esp{muy}, \esp{tan}, \esp{bastante} et \esp{mucho} (avec un verbe) sont invariables.
\end{corrige}

\begin{corrige}[Exercice 4 — Conjugaison : l'impersonnel à tous les temps]
\begin{enumerate}
\item \esp{Aquí \textbf{se vende} artesanía africana.} (présent, 3\up{e} pers. sing.)
\item \esp{\textbf{Se ha hablado} mucho de la paz en la reunión.} (passé composé : auxiliaire \esp{haber} conjugué + participe invariable)
\item \esp{Antes \textbf{se vivía} mejor en este barrio.} (imparfait)
\item \esp{Mañana \textbf{se discutirá} el problema en la radio.} (futur simple)
\item \esp{En los exámenes no \textbf{se puede} copiar.} (verbe modal \esp{poder} conjugué + infinitif)
\end{enumerate}
Point de méthode : dans la construction impersonnelle, seul le verbe conjugué porte la marque de la 3\up{e} personne du singulier ; \esp{se} se place devant le verbe conjugué (ou attaché à l'infinitif : \esp{no puede copiarse}).
\end{corrige}

\begin{corrige}[Exercice 5 — Texte à trous : muy, mucho, demasiado, bastante, tan, tanto]
Texte complété :
\begin{quote}
\esp{En mi barrio de Yaoundé hay una emisora de radio \textbf{muy} popular. Habla \textbf{mucho} de la cultura de la paz, y sus programas son \textbf{tan} interesantes que toda la familia los escucha. Ayer hubo \textbf{mucho} ruido en el estudio porque vinieron \textbf{muchos} jóvenes del liceo. El presentador estaba \textbf{muy} contento: dijo que los jóvenes tienen \textbf{mucha} energía y \textbf{muchas} ideas nuevas. Pero algunos oyentes se quejaron: para ellos, el programa duró \textbf{demasiado} tiempo y la música era \textbf{demasiado} alta. Al final, el presentador prometió trabajar \textbf{más} para mejorar.}
\end{quote}
Justifications : \esp{muy popular} / \esp{muy contento} / \esp{demasiado alta} : devant un adjectif, adverbe invariable. \esp{Habla mucho} : avec un verbe. \esp{tan interesantes que\ldots} : corrélation d'intensité + conséquence. \esp{mucho ruido}, \esp{mucha energía}, \esp{muchos jóvenes}, \esp{muchas ideas}, \esp{demasiado tiempo} : devant un nom, accord en genre et en nombre. Pour la dernière case, \esp{trabajar más} (« davantage ») est la solution la plus naturelle ; \esp{trabajar mucho} est accepté si l'élève n'a pas encore utilisé \esp{mucho} avec un verbe.
\end{corrige}

\begin{corrige}[Exercice 6 — Appariements : chaque « on » a son visage]
\begin{center}
\begin{tabularx}{\linewidth}{|c|X|X|}
\hline
\rowcolor{popblueL} N° & Traduction & Valeur du « on » \\
\hline
1 & C — \esp{En España la gente cena muy tarde.} & les gens, les Espagnols en général \\
\hline
2 & D — \esp{Nos han invitado a la boda de Carla.} & quelqu'un d'indéfini (3\up{e} pers. pluriel) \\
\hline
3 & B — \esp{Mis amigos y yo partimos mañana para Kribi.} & nous \\
\hline
4 & A — \esp{Cuando uno es joven, cree saberlo todo.} & quiconque, valeur générale \\
\hline
5 & E — \esp{Aquí se habla francés e inglés.} & impersonnel \\
\hline
\end{tabularx}
\end{center}
Remarque : dans la phrase 2, \esp{nos} (à nous) est le complément d'objet indirect ; le « on » français devient le sujet indéfini de la 3\up{e} personne du pluriel \esp{han invitado}.
\end{corrige}

\begin{corrige}[Exercice 7 — Reformulation : trois façons de dire]
\begin{enumerate}
\item Impersonnel : \esp{Se ha construido un puente sobre el río Wouri.} — On ne dit pas qui a construit le pont ; l'agent est effacé, on insiste sur le résultat (le pont existe).
\item 3\up{e} personne du pluriel : \esp{Han construido un puente sobre el río Wouri.} — On pense à des personnes indéfinies mais réelles (les ouvriers, les autorités, « ils ») sans les nommer.
\item \esp{Nosotros} : \esp{Nosotros hemos construido un puente sobre el río Wouri.} — Le locuteur s'inclut dans le groupe des constructeurs (les habitants, notre communauté) : valeur de « nous ».
\end{enumerate}
Traduction française commune : « On a construit un pont sur le fleuve Wouri », mais la nuance d'agent change selon la construction choisie.
\end{corrige}

\begin{corrige}[Exercice 8 — Corrige les erreurs d'élèves francophones]
\begin{enumerate}
\item \esp{*Este libro es mucho bueno.} $\rightarrow$ \esp{Este libro es \textbf{muy} bueno.} Devant un adjectif, « très » se dit \esp{muy} ; \esp{mucho} ne modifie jamais un adjectif.
\item \esp{*Se vende frutas en el mercado.} $\rightarrow$ \esp{Se \textbf{venden} frutas en el mercado.} Le nom pluriel \esp{frutas} joue le rôle de sujet : le verbe s'accorde au pluriel.
\item \esp{*Mi hermano es tanto alto como yo.} $\rightarrow$ \esp{Mi hermano es \textbf{tan} alto como yo.} Devant un adjectif, \esp{tanto} s'apocope en \esp{tan} (comparatif d'égalité : \esp{tan\ldots como}).
\item \esp{*La gente son muy amables aquí.} $\rightarrow$ \esp{La gente \textbf{es} muy \textbf{amable} aquí.} \esp{La gente} est un singulier collectif : verbe et adjectif au singulier.
\item \esp{*Me gusta muy esta canción.} $\rightarrow$ \esp{Me gusta \textbf{mucho} esta canción.} (ou \esp{me gusta muchísimo}). Avec un verbe, « beaucoup » se dit \esp{mucho} ; \esp{muy} est impossible.
\item \esp{*Hay demasiado coches en Douala.} $\rightarrow$ \esp{Hay \textbf{demasiados} coches en Douala.} Devant un nom masculin pluriel, \esp{demasiado} s'accorde : \esp{demasiados}.
\end{enumerate}
\end{corrige}

\begin{corrige}[Exercice 9 — Texte : compréhension et questions de langue]
\textbf{I. Comprensión lectora} (réponses modèles)
\begin{enumerate}
\item \esp{Según la gente, los medios de comunicación tienen un papel muy importante: enseñan a respetar las opiniones de los demás.}
\item \esp{Las noticias falsas crean demasiados conflictos entre los jóvenes.}
\item \esp{Se han creado los clubes de prensa escolar para que los alumnos aprendan a verificar la información antes de compartirla.}
\item \esp{El profesor observa que ahora los chicos son bastante más críticos: antes creían todo lo que leían en Internet.}
\end{enumerate}
\textbf{II. Questions de langue}
\begin{enumerate}
\item Trois constructions du « on » : \esp{se habla cada vez más} (impersonnel avec \esp{se}) ; \esp{la gente piensa} (les gens) ; \esp{cuando uno escucha la radio} (quiconque, valeur générale). On peut ajouter \esp{se difunden noticias falsas} et \esp{se han creado clubes} (passif réfléchi / impersonnel).
\item Dans \esp{la gente piensa}, le verbe est au singulier parce que \esp{gente} est un nom collectif féminin singulier, même s'il désigne une multitude de personnes.
\item On écrit \esp{demasiados conflictos} parce que \esp{demasiado} est ici un déterminant qui accompagne un nom masculin pluriel : il s'accorde en genre et en nombre (\esp{demasiado ruido}, \esp{demasiada gente}, \esp{demasiados conflictos}, \esp{demasiadas noticias}).
\item \esp{Se verifica la información.} (impersonnel avec \esp{se} + 3\up{e} personne du singulier ; \esp{la información} est singulier).
\end{enumerate}
Barème indicatif (sur 20) : compréhension 8 pts (2 par question), questions de langue 12 pts (4 pts pour la question 1, 3 pts pour la 2, 3 pts pour la 3, 2 pts pour la 4).
\end{corrige}

\begin{corrige}[Exercice 10 — Thème type Bac]
Traductions proposées :
\begin{enumerate}
\item \esp{No se trabaja los domingos en esta oficina.} (impersonnel ; « le dimanche » = \esp{los domingos})
\item \esp{Mis primos y yo queremos mucho al fútbol camerounés.} — formulation la plus naturelle : \esp{A mis primos y a mí nos gusta mucho el fútbol camerounés.} (« on » = nous ; « beaucoup » avec un verbe = \esp{mucho})
\item \esp{Me han dicho que esta película era muy interesante.} (on = quelqu'un d'indéfini : 3\up{e} pers. du pluriel ; « très » devant adjectif = \esp{muy})
\item \esp{Cuando uno viaja demasiado, está a menudo cansado.} (on = quiconque : \esp{uno} ; « trop » avec un verbe = \esp{demasiado} invariable)
\item \esp{La gente bebe bastante agua cuando hace mucho calor.} (\esp{la gente} + singulier ; \esp{bastante} invariable devant le nom en espagnol standard ; « très chaud » = \esp{mucho calor} car \esp{calor} est un nom : \esp{hace mucho calor}.)
\end{enumerate}
Points de vigilance : ne jamais écrire \esp{*muy calor} ; \esp{gustar} se construit avec un complément d'objet indirect (\esp{nos gusta}) ; \esp{demasiado} est invariable avec un verbe mais s'accorde devant un nom.
Barème indicatif (sur 10) : 2 pts par phrase (1 pt pour la traduction correcte du « on », 1 pt pour l'adverbe d'intensité et l'exactitude globale).
\end{corrige}

\begin{corrige}[Exercice 11 — Version type Bac]
Traduction proposée :
\begin{quote}
Dans notre quartier, on organise chaque année une grande fête pour célébrer la paix. Les gens arrivent très tôt : on vend des plats typiques, on danse et on chante jusqu'à très tard. Quand on participe à cette fête, on comprend que la vie en communauté est bien plus facile qu'il n'y paraît. Les anciens racontent des histoires du passé et les enfants écoutent avec une attention extrême pour ne rien oublier. À la fin de la nuit, tous se disent : « l'année prochaine, nous reviendrons ».
\end{quote}
Choix de traduction commentés :
\begin{itemize}
\item \esp{se organiza}, \esp{se venden}, \esp{se baila}, \esp{se canta} : « on » impersonnel français ; noter l'accord \esp{se venden platos} (pluriel).
\item \esp{la gente llega} : « les gens arrivent » (et non « on arrive », pour varier).
\item \esp{cuando uno participa} : « quand on participe », valeur générale.
\item \esp{bastante más fácil} : « bien plus facile / nettement plus facile ».
\item \esp{con demasiada atención} : littéralement « avec trop d'attention », à rendre par « avec une attention extrême » pour garder le sens positif du texte.
\item \esp{todos se dicen} : « tous se disent » (\esp{se} réciproque, à ne pas confondre avec l'impersonnel).
\end{itemize}
Barème indicatif (sur 20) : fidélité au sens 10 pts, choix des équivalents de \esp{se} / \esp{uno} / \esp{la gente} 5 pts, qualité du français 5 pts.
\end{corrige}

\begin{corrige}[Exercice 12 — Expression écrite : sujet de rédaction]
Proposition de rédaction modèle (environ 165 mots) :
\begin{quote}
\esp{En mi país, Camerún, los medios de comunicación tienen un papel muy importante en la promoción de la paz. La gente escucha la radio todos los días, porque es el medio más popular, sobre todo en los barrios y en los pueblos. Cuando uno oye programas sobre la convivencia y el respeto, aprende a dialogar en vez de pelear. Sin embargo, en las redes sociales se difunden a veces noticias falsas, y eso crea demasiados conflictos entre los jóvenes. Por eso, en muchos liceos se han creado clubes de prensa escolar: allí, nosotros aprendemos a verificar la información antes de compartirla. La televisión también puede ayudar mucho: hay programas bastante interesantes que muestran la riqueza de todas las culturas del país. Pero los periodistas deben trabajar más para informar con honestidad. Si se educa bien a la juventud, se construye una sociedad más pacífica. La tarea es larga, pero no es imposible: cada uno de nosotros puede ser un mensajero de la paz.}
\end{quote}
Traduction du modèle : « Dans mon pays, le Cameroun, les médias jouent un rôle très important dans la promotion de la paix. Les gens écoutent la radio tous les jours, car c'est le média le plus populaire, surtout dans les quartiers et les villages. Quand on entend des émissions sur la vie en communauté et le respect, on apprend à dialoguer au lieu de se battre. Cependant, sur les réseaux sociaux, on diffuse parfois de fausses nouvelles, et cela crée trop de conflits entre les jeunes. C'est pourquoi, dans beaucoup de lycées, on a créé des clubs de presse scolaire : là, nous apprenons à vérifier l'information avant de la partager. La télévision peut aussi beaucoup aider : il y a des émissions assez intéressantes qui montrent la richesse de toutes les cultures du pays. Mais les journalistes doivent travailler davantage pour informer avec honnêteté. Si on éduque bien la jeunesse, on construit une société plus pacifique. La tâche est longue, mais elle n'est pas impossible : chacun de nous peut être un messager de la paix. »

Vérification des contraintes dans le modèle :
\begin{itemize}
\item Trois traductions du « on » : \esp{la gente escucha}, \esp{cuando uno oye}, \esp{se difunden} / \esp{se han creado} / \esp{se educa}, \esp{nosotros aprendemos}.
\item Quatre adverbes d'intensité : \esp{muy importante}, \esp{demasiados conflictos}, \esp{bastante interesantes}, \esp{mucho} (avec \esp{ayudar}), \esp{más}.
\item Connecteurs : \esp{porque}, \esp{por eso} (cause) ; \esp{sin embargo}, \esp{pero} (opposition).
\end{itemize}
Points de vigilance pour les élèves : accorder \esp{demasiados} avec \esp{conflictos} ; verbe au singulier après \esp{la gente} et \esp{uno} ; \esp{se difunden} au pluriel devant \esp{noticias falsas} ; ne pas écrire \esp{*muy conflictos} ni \esp{*mucho importante}.
Barème indicatif (sur 20) : respect du sujet et de la longueur 4 pts, emploi correct des trois « on » 6 pts, emploi correct des adverbes d'intensité 4 pts, connecteurs et cohérence 3 pts, correction linguistique générale 3 pts.
\end{corrige}

\newpage

\coursec{Fiche bilan}

\begin{popfigure}
\begin{tikzpicture}[
    mindmap,
    grow cyclic,
    every node/.style={concept, execute at begin node=\hskip0pt},
    concept color=popblue,
    level 1/.append style={level distance=4.5cm, sibling angle=51.4},
    level 2/.append style={level distance=3cm, sibling angle=45},
    texte court/.style={concept, font=\small\bfseries, text width=2.8cm, align=center, inner sep=4pt}
]
\node[texte court, concept color=popblue, align=center]{Leçon E23\\Traduction : « on » et Intensité}
    child[concept color=poporange] { node[texte court, align=center]{« on» = \cle{se} impersonnel\\ \emph{Se vive bien}} }
    child[concept color=popgreen] { node[texte court, align=center]{« on» = 3e pers. pl.\\ \emph{Dicen que...}} }
    child[concept color=poppurple] { node[texte court, align=center]{« on» = \cle{uno} / \cle{la gente}\\ \emph{Uno nunca sabe}} }
    child[concept color=poppink] { node[texte court, align=center]{« on» = \cle{nosotros}\\ \emph{Vamos, ¿no?}} }
    child[concept color=popgold] { node[texte court, align=center]{\cle{muy} + adj/adv\\ \cle{mucho} + verbe/nom} }
    child[concept color=popdark] { node[texte court, align=center]{\cle{demasiado}, \cle{bastante}\\ \cle{tan}, \cle{tanto}, apocopes} };
\end{tikzpicture}
\legende{Carte mentale : les 5 visages du « on » français en espagnol et les adverbes d'intensité clés.}
\end{popfigure}

\begin{definition}[Lexique clé — Traduction et Intensité]
\begin{center}
\begin{tabularx}{\linewidth}{|X|X|X|}
\hline
\rowcolor{popblueL}
\textbf{Español} & \textbf{Français} & \textbf{Contexte / Note} \\
\hline
\esp{se} + verbe (3e pers. sg) & \cle{on} (impersonnel) & Généralités, coutumes, règles. \esp{Se come bien aquí.} \\
\hline
Verbe 3e pers. pl (ils) & \cle{on} (indéfini, vague) & « On dit », « on a vu ». \esp{Dicen que llueve.} \\
\hline
\esp{uno} / \esp{una} & \cle{on} (généralisant, sage) & Sujet générique, humain. \esp{Uno debe estudiar.} \\
\hline
\esp{la gente} + verbe sg & \cle{on} / \cle{les gens} (familier) & Vie courante, oral. \esp{La gente piensa que...} \\
\hline
\esp{nosotros} & \cle{on} = \cle{nous} (inclusif) & Familier, complicité. \esp{¿Vamos al cine?} \\
\hline
\esp{muy} + adj/adv/part. passé & \cle{très} & Invariable. \esp{Muy rápido, muy cansado.} \\
\hline
\esp{mucho} + verbe / \esp{mucho/a/os/as} + nom & \cle{beaucoup} / \cle{très} (verbe) & Accord si déterminant. \esp{Como mucho. Tengo mucha hambre.} \\
\hline
\esp{demasiado} & \cle{trop} (excès) & Accord possible. \esp{Demasiado ruido.} \\
\hline
\esp{bastante} & \cle{assez} / \cle{pas mal de} & Invariable en adv, variable en det. \esp{Bastante caro. Bastantes libros.} \\
\hline
\esp{tan} + adj/adv & \cle{si} / \cle{aussi} & Corrélatif \esp{tan... como}. \esp{Tan guapo como tú.} \\
\hline
\esp{tanto/a/os/as} + nom & \cle{tant de} / \cle{autant de} & Accord obligatoire. \esp{Tanta gente, tantos libros.} \\
\hline
\esp{buen, mal, gran, primer, tercer, algún, ninguno} & Apocopes & Devant nom masc. sg. \esp{Un buen libro, mal tiempo.} \\
\hline
\end{tabularx}
\end{center}
\end{definition}

\begin{propriete}[Règles de grammaire à connaître par cœur]
\begin{center}
\begin{tabularx}{\linewidth}{|>{\bfseries}X|X|X|}
\hline
\rowcolor{poporangeL}
\centering \textbf{Notion} & \textbf{Règle / Forme} & \textbf{Exemple traduit} \\
\hline
\cle{Se impersonnel} & \esp{Se} + verbe 3e sg (ou pl si COD placé avant) & \esp{Se habla español.} « On parle espagnol. » \\
\hline
\cle{Uno} & Indéfini, déclenche accords au masc. sg. & \esp{Uno se siente solo.} « On se sent seul. » \\
\hline
\cle{Muy vs Mucho} & \esp{Muy} + adj/adv (invariable). \esp{Mucho} + verbe (invariable) OU \esp{mucho/a/os/as} + nom (variable). & \esp{Es muy listo.} (très) \\ \esp{Estudia mucho.} (beaucoup) \\ \esp{Tiene mucho trabajo.} (beaucoup de) \\
\hline
\cle{Demasiado / Bastante} & Adv = invariable. Det = variable (genre/nombre). & \esp{Es demasiado tarde.} (adv) \\ \esp{Demasiados errores.} (det) \\
\hline
\cle{Tan / Tanto} & \esp{Tan} + adj/adv (invariable). \esp{Tanto/a/os/as} + nom (variable). Corrélatif : \esp{tan... como}, \esp{tanto... como}. & \esp{Tan rico como tú.} \\ \esp{Tanto dinero como yo.} \\
\hline
\cle{Apocopes} & \esp{Bueno, malo, grande, primero, tercero, alguno, ninguno} $\to$ \esp{buen, mal, gran, primer, tercer, algún, ningún} devant nom masc. sg. & \esp{Un buen amigo.} \esp{Mal resultado.} \esp{Gran hombre.} \esp{Primer paso.} \esp{Tercer piso.} \esp{Algún día.} \esp{Ningún problema.} \\
\hline
\end{tabularx}
\end{center}
\end{propriete}

\begin{methode}[Méthodes express pour le thème et la version]
\begin{enumerate}
\item \textbf{Traduire « on » :} Identifier le sens $\to$ Généralité/impersonnel = \cle{se} (verbe au singulier) ; Témoignage/« on dit » = 3e personne du pluriel (\esp{dicen, hacen, ven}) ; Conseil/morale = \cle{uno} (accords au masculin singulier) ; Familier/« les gens » = \cle{la gente} (verbe au singulier) ; « Nous » inclusif = \cle{nosotros}.
\item \textbf{Choisir \esp{muy} ou \esp{mucho} :} Test du remplacement. Si « très » + adjectif/adverbe $\to$ \esp{muy} (invariable). Si « beaucoup » + verbe $\to$ \esp{mucho} (invariable). Si « beaucoup de » + nom $\to$ \esp{mucho/a/os/as} (accordé).
\item \textbf{Intensité + nom :} Exclamatif \esp{Qué} + nom (\esp{¡Qué calor!}) ou \esp{tan} + adjectif (\esp{tan caloroso}). Quantité : \esp{tanto/a/os/as} + nom = « tant de / autant de » (accordé).
\item \textbf{Apocopes :} Vérifier si l'adjectif (\esp{bueno, malo, grande, primero, tercero, alguno, ninguno}) précède un nom masculin singulier $\to$ couper la finale et ajouter l'accent sur \esp{algún, ningún}.
\end{enumerate}
\end{methode}

\begin{aretenir}[Checklist avant le Bac]
\begin{itemize}
\item $\square$ Je sais traduire « on » par \esp{se} (impersonnel) et conjuguer le verbe au singulier.
\item $\square$ Je sais utiliser la 3e personne du pluriel pour « on » indéfini (\esp{dicen, hacen, ven}).
\item $\square$ Je distingue \esp{uno} (généralisant, formel) de \esp{la gente} (familier, verbe au singulier).
\item $\square$ J'identifie le « on » = « nous » (inclusif) $\to$ \esp{nosotros}.
\item $\square$ Je ne confonds plus \esp{muy} (invariable, + adj/adv) et \esp{mucho} (variable si + nom, invariable si + verbe).
\item $\square$ J'accorde \esp{demasiado, bastante, tanto} quand ils sont déterminants (devant nom).
\item $\square$ J'utilise \esp{tan} + adjectif/adverbe (invariable) et \esp{tanto} + nom (variable).
\item $\square$ J'applique les apocopes : \esp{buen, mal, gran, primer, tercer, algún, ningún} devant nom masc. sg.
\item $\square$ Je traduis « trop » par \esp{demasiado} (excès négatif) et non \esp{muy mucho}.
\item $\square$ Je construis les comparatifs d'égalité : \esp{tan... como} (adj/adv) et \esp{tanto... como} (nom).
\item $\square$ Je relis mes thèmes/versions pour vérifier l'accord du participe passé avec \esp{se} : masc. sg. par défaut pour le \emph{se} impersonnel pur ; accord avec le sujet (COD devenu sujet) pour le \emph{se} passif (\esp{Se han vendido las entradas}).
\end{itemize}
\end{aretenir}

\findecours

\end{document}
